\documentclass[11pt]{template/cauchy}
\usepackage[czech]{babel}

% IMPORTY
\usepackage{enumerate}
\usepackage{mlmodern}
\usepackage{template/mathsphystools}
\usepackage{amsthm,thmtools,xcolor,amsmath,amssymb, mathtools}
\usepackage[scr]{rsfso}
% \usepackage{thmstyles}
\usepackage{tabularx}
\usepackage{graphicx}
\usepackage{multicol}
\usepackage{appendix}
\usepackage{tabularx}
\usepackage{subcaption}
\usepackage{hyperref}
\usepackage{comment}
\usepackage{lscape}
\usepackage{xcolor}
\usepackage{wrapfig}
\usepackage{centernot}
\usepackage{pgfplots}
\usepackage{caption}
\usepackage{fontawesome}
\graphicspath{{images/}}
\pgfplotsset{width=10cm,compat=1.9}


% \usepackage{tikzit}
% \input{default.tikzstyles}
% je potřeba přiložit soubor default.tikzstyles


% STYLY
% definice, příklad, cvičení, věta, lemma, řešení, poznámka, důsledek, axiom, motivace

\declaretheoremstyle[
  headfont=\color{red}\normalfont\bfseries,
  bodyfont=\color{black}\normalfont,
]{def}

\declaretheoremstyle[
  headfont=\color{blue}\normalfont\bfseries,
  bodyfont=\color{black}\normalfont,
  numbered=no
]{pr}

\declaretheoremstyle[
  headfont=\color{teal}\normalfont\bfseries,
  bodyfont=\color{black}\normalfont\itshape,
  numbered=no
]{cv}

\declaretheoremstyle[
  headfont=\color{green}\normalfont\bfseries,
  bodyfont=\color{black}\normalfont\itshape,
]{veta}

\declaretheoremstyle[
  headfont=\color{lime}\normalfont\bfseries,
  bodyfont=\color{black}\normalfont\itshape,
]{lemma}

\declaretheoremstyle[
  headfont=\color{black}\normalfont\itshape,
  bodyfont=\color{black}\normalfont,
  numbered=no
]{res}

\declaretheoremstyle[
  headfont=\color{brown}\normalfont\bfseries,
  bodyfont=\color{black}\normalfont,
]{pozn}

\declaretheoremstyle[
  headfont=\color{olive}\normalfont\bfseries,
  bodyfont=\color{black}\normalfont,
]{dusl}

\declaretheoremstyle[
  headfont=\color{cyan}\normalfont\bfseries,
  bodyfont=\color{black}\normalfont,
]{axiom}

\declaretheoremstyle[
  headfont=\color{black}\normalfont\bfseries,
  bodyfont=\color{black}\normalfont,
  numbered=no
]{motivace}

\declaretheoremstyle[
  headfont=\color{olive}\normalfont\bfseries,
  bodyfont=\color{black}\normalfont,
  numbered=no
]{ozn}

\declaretheoremstyle[
  headfont=\color{black}\normalfont\bfseries,
  bodyfont=\color{black}\normalfont,
  numbered=no,
  headpunct=
]{pozorovani}

\declaretheorem[
  style=def,
  name=Definice,
]{definice}

\declaretheorem[
  style=ozn,
  name=Označení,
]{oznaceni}

\declaretheorem[
  style=pr,
  name=Příklad,
]{priklad}

\declaretheorem[
  style=cv,
  name=Cvičení,
]{cviceni}

\declaretheorem[
  style=res,
  name=Řešení,
]{reseni}

\declaretheorem[
  style=veta,
  name=Věta,
]{veta}

\declaretheorem[
  style=lemma,
  name=Lemma,
]{lemma}

\declaretheorem[
  style=pozn,
  name=Poznámka,
]{poznamka}

\declaretheorem[
  style=dusl,
  name=Důsledek,
]{dusledek}

\declaretheorem[
  style=axiom,
  name=Axiom,
]{axiom}

\declaretheorem[
  style=motivace,
  name=Motivace,
]{motivace}

\declaretheorem[
  style=pozorovani,
  name=\faEye,
]{pozorovani}

\DeclareMathOperator{\tg}{tg}
\DeclareMathOperator{\cotg}{cotg}
\DeclareMathOperator{\arctg}{arctg}
\DeclareMathOperator{\arccotg}{arccotg}

\newcommand\nm[1]{\lVert #1 \rVert}

\title{Matematická analýza \\ {\large Zápisky z přednášky doc. Miroslava Zeleného} \\ {\large verze z \today}}
\author{Dominik Doležel}

\begin{document}
\frontmatter
\maketitle
\thispagestyle{empty}
\clearpage

\pagenumbering{roman}
\tableofcontents

\clearpage

% ÚVODNÍ INFORMACE
% dobré pro poznačení informačních věcí od učitele,
% např. jeho email, jména skript apod.
%%!TEX root = ../main.tex
%
\phantomsection
\addcontentsline{toc}{section}{Úvodní informace}
\begin{center}
{\bfseries\Large Úvodní informace}\par\vspace{1em}
\end{center}

% SEM PATŘÍ TEXT


% ZNAČENÍ
% pokud nepotřebujete, zakomentujte řádek níže
% všechny změny provádějte v souboru additional/conventions.tex
%%!TEX root = ../main.tex
%
\phantomsection
\addcontentsline{toc}{section}{Značení}
\vspace{2em}
\begin{center}
{\bfseries\Large Značení}\par\vspace{1em}
\end{center}

% SEM PATŘÍ TEXT


% KAPITOLY
% \input{chapters/01_nazev_kapitoly}
% ...
\mainmatter
\section{Logika, množiny a základní číselné obory}
\subsection{Logika}

\textbf{Výrokem} nazveme jakékoliv tvrzení, o němž má smysl říci, že platí (je pravdivé), nebo že neplatí (je nepravdivé).

\begin{definice}
\textbf{Negací} $\neg A$ výroku $A$ rozumíme výrok:
\begin{center}
	\emph{Není pravda, že platí $A$.}
\end{center}
\[
\begin{array}{c|c}
A & \neg A \\
\hline
0 & 1 \\
1 & 0
\end{array}
\]
\end{definice}

\begin{definice}
\textbf{Konjunkcí} $A \wedge B$ výroků $A$ a $B$ nazveme výrok:

\begin{center}
	\emph{Platí $A$ i $B$.}
\end{center}
\end{definice}

\begin{definice}
\textbf{Disjunkcí} $A \vee B$ výroků $A$ a $B$ nazveme výrok:

\begin{center}
	\emph{Platí $A$ nebo $B$.}
\end{center}
\end{definice}

\begin{definice}
\textbf{Implikací} $A \implies B$ nazýváme výrok:

\begin{center}
	\emph{Jestliže platí výrok $A$, potom platí výrok $B$.}
\end{center}

Výroku $A$ v implikaci se říká \textbf{premisa}, výrok $B$ se nazývá \textbf{závěr}.
Výrok $A$ je \textbf{postačující podmínkou} pro platnost $B$ a $B$ je \textbf{nutnou podmínkou} pro platnost $A$.
\end{definice}

\begin{definice}
\textbf{Ekvivalencí} $A \iff B$ nazýváme výrok:

\begin{center}
	\emph{Výrok $A$ platí tehdy a jen tehdy, když platí výrok $B$.}
\end{center}

(Platnost výroku) $A$ je \textbf{nutnou a postačující} podmínkou (platnosti výroku) $B$.

\[
\begin{array}{c|c|c|c|c|c}
A & B & A\wedge B & A\vee B & A\Rightarrow B & A\Leftrightarrow B \\
\hline
0 & 0 & 0 & 0 & 1 & 1 \\
0 & 1 & 0 & 1 & 1 & 0 \\
1 & 0 & 0 & 1 & 0 & 0 \\
1 & 1 & 1 & 1 & 1 & 1
\end{array}
\]

\textbf{Výroková forma} $V$ je výraz, který má konečný počet proměnných, přičemž když za tyto proměnné dosadíme prvky z daného oboru, obdržíme výrok.
\end{definice}

\begin{definice}
Nechť $V$ je výroková forma s jednou proměnnou.
\begin{enumerate}[(i)]
\item Výrok „Pro každé $x$ platí $V(x)$.“ symbolicky zapisujeme ve tvaru
\[
\forall x : V(x).
\]
Symbol $\forall$ nazýváme \textbf{obecným kvantifikátorem}.
\item Výrok „Existuje $x$ takové, že platí $V(x)$.“ zapisujeme ve tvaru
\[
\exists x : V(x).
\]
Symbol $\exists$ nazýváme \textbf{existenčním kvantifikátorem}.
\end{enumerate}
\end{definice}

\subsection{Metody důkazů}

\begin{itemize}
\item přímý důkaz
\item nepřímý důkaz
\item důkaz sporem
\item důkaz rozborem případů
\item důkaz matematickou indukcí
\end{itemize}

\subsection{Množiny}

\paragraph{G. Cantor:} \uv{\textbf{Množinou} rozumíme každé shrnutí určitých a navzájem různých objektů, které nazýváme \textbf{prvky}, do jediného celku.}
\par Množinu definujeme výčtem prvků nebo pomocí vlastnosti, kterou musejí splňovat její prvky, tj. píšeme
$\{x \in M ; V(x)\}$, kde $M$ je množina a $V$ je výroková forma.

\begin{definice}
Řekneme, že množina $A$ je \textbf{částí množiny} $B$ (nebo $A$ je \textbf{podmnožinou} $B$), jestliže každý prvek množiny $A$ je rovněž prvkem množiny $B$. Tomuto vztahu
říkáme \textbf{inkluze} a značíme jej $A \subset B$.
Množiny $A$ a $B$ jsou si rovny ($A = B$), jestliže mají stejné prvky.
\textbf{Prázdnou množinou} nazveme množinu, která neobsahuje žádný prvek. Označíme ji symbolem $\emptyset$.
\end{definice}

\begin{definice}
\textbf{Sjednocením} množin $A$ a $B$ nazveme množinu tvořenou všemi prvky, které patří alespoň do jedné z množin $A$ či $B$. Značíme $A \cup B$.
\par Je-li $\mathcal{A}$ systém množin, pak jeho \textbf{sjednocení}
$\bigcup \mathcal{A}$
definujeme jako množinu všech prvků $a$, pro které existuje $A \in \mathcal{A}$ takové, že $a \in A$.
\end{definice}

\begin{definice}
\textbf{Průnikem} množin $A$ a $B$ nazveme množinu všech prvků, které náležejí současně do $A$ i do $B$. Značíme $A \cap B$. Mají-li množiny $A$ a $B$
prázdný průnik, řekneme o nich, že jsou \textbf{disjunktní}.
\par Je-li $\mathcal A$ neprázdný systém množin, pak jeho \textbf{průnik} $ \bigcap \mathcal A$ definujeme jako
množinu všech prvků $a$, které pro každé $A\in \mathcal A$ splňují $a\in A$.
\end{definice}

\begin{definice}
\textbf{Rozdílem} množin $A$ a $B$ nazveme množinu prvků, které patří do množiny $A$ a nepatří do množiny $B$. Značíme $A \smallsetminus B$.
\end{definice}

\begin{definice}
\textbf{Kartézským součinem} množin $A_1,\dots,A_n$ nazveme množinu všech uspořádaných $n$-tic
\[
A_1\times A_2 \times \cdots \times A_n
= \{[a_1,\dots,a_n];\, a_i \in A_i\}.
\]
\end{definice}

\begin{veta}[de Morganova pravidla]\label{veta1.1}
Nechť $X$ je množina a $\mathcal{A}$ je neprázdný systém množin. Pak platí
\[
X \smallsetminus \bigcup \mathcal{A}
= \bigcap \{X \smallsetminus A;\, A \in \mathcal{A}\}
\]
a dále
\[
X \smallsetminus \bigcap \mathcal{A}
= \bigcup \{X \smallsetminus A;\, A \in \mathcal{A}\}.
\]
\end{veta}

\subsection{Relace uspořádání a zobrazení}

\begin{definice}
\textbf{Binární relací} rozumíme libovolnou množinu uspořádaných dvojic. Pokud $R$ je
binární relace a $[a,b]\in R$, pak říkáme, že prvek $a$ je v relaci $R$ s prvkem $b$.
Často v tomto případě používáme zápis $a\,R\,b$.
\par Pokud binární relace $R$ splňuje $R\subset A\times B$, pak říkáme, že $R$ je \textbf{binární relací
mezi prvky množin} $A$ \textbf{a} $B$. Pokud $A=B$, pak říkáme, že $R$ je \textbf{binární relací na} $A$.
\end{definice}

\begin{definice}
Nechť $X$ je množina a $R$ je relace na $X$. Řekneme, že $R$ je
\begin{itemize}
\item \textbf{reflexivní}, jestliže pro každé $x\in X$ platí $[x,x]\in R$,
\item \textbf{symetrická}, jestliže pro každé $x,y\in X$ splňující $[x,y]\in R$ platí $[y,x]\in R$,
\item \textbf{tranzitivní}, jestliže pro každé $x,y,z\in X$ splňující $[x,y]\in R$ a $[y,z]\in R$
platí $[x,z]\in R$,
\item \textbf{antisymetrická}, jestliže pro každé $x,y\in X$ splňující $[x,y]\in R$ platí
$[y,x]\notin R$,
\item \textbf{slabě antisymetrická}, jestliže pro každé $x,y\in X$ splňující
$[x,y]\in R$ a $[y,x]\in R$ platí $x=y$.
\end{itemize}
\end{definice}

\begin{definice}
Nechť $R$ je relace na množině $A$. Řekneme, že $R$ je na $A$
\begin{itemize}
\item \textbf{uspořádání} (někdy také částečné uspořádání či neostré uspořádání),
jestliže je reflexivní, slabě antisymetrická a tranzitivní,
\item \textbf{ostré uspořádání}, jestliže je antisymetrická a tranzitivní,
\item \textbf{lineární uspořádání}, jestliže jde o uspořádání a pro každé $x,y\in A$ platí
$[x,y]\in R$ nebo $[y,x]\in R$.
\end{itemize}
\end{definice}

\begin{definice}
Nechť $\le$ je relace uspořádání na množině $X$ a $A\subset X$. Řekneme, že prvek
$x\in X$ je
\begin{itemize}
\item \textbf{maximálním} prvkem (\textbf{maximem}) množiny $A$, jestliže $x\in A$ a neexistuje
$a\in A$ takové, že $x\le a$ a $x\ne a$,
\item \textbf{největším} prvkem množiny $A$, jestliže $x\in A$ a pro každé $a\in A$ platí
$a\le x$.
\end{itemize}
\par Pojmy \textbf{minimální} prvek (\textbf{minimum}) množiny a \textbf{nejmenší} prvek množiny jsou definovány
zřejmým způsobem.
\end{definice}

\begin{definice}
Nechť $\le$ je relace uspořádání na množině $X$ a $A\subset X$. Řekneme, že prvek
$x\in X$ je
\begin{itemize}
\item \textbf{horní závorou} množiny $A$, jestliže pro každé $a\in A$ platí $a\le x$,
\item \textbf{dolní závorou} množiny $A$, jestliže pro každé $a\in A$ platí $x\le a$.
\end{itemize}
\par Množina $A$ je
\begin{itemize}
\item \textbf{shora omezená}, jestliže existuje prvek $x\in X$, který je horní závorou množiny $A$,
\item \textbf{zdola omezená}, jestliže existuje prvek $x\in X$, který je dolní závorou množiny $A$,
\item \textbf{omezená}, jestliže je omezená shora i zdola.
\end{itemize}
\end{definice}

\begin{definice}
Nechť $\le$ je relace uspořádání na množině $X$ a $M\subset X$. Řekneme, že prvek
$G\in X$ je \textbf{supremem} množiny $M$, jestliže platí:
\begin{enumerate}
\item $G$ je horní závorou množiny $M$,
\item je-li prvek $G'\in X$ horní závorou množiny $M$, potom $G\le G'$.
\end{enumerate}

Řekneme, že prvek $g\in X$ je \textbf{infimem} množiny $M$, jestliže platí:
\begin{enumerate}
\item $g$ je dolní závorou množiny $M$,
\item je-li prvek $g'\in X$ dolní závorou množiny $M$, potom $g'\le g$.
\end{enumerate}

\end{definice}


\begin{veta}\label{veta1.2}
Nechť $\le$ je relace uspořádání na množině $X$, $M\subset X$ je neprázdná množina
a existuje infimum a supremum množiny $M$. Potom platí
$\inf M \le \sup M.$
\end{veta}

\begin{definice}
	Binární relaci $F$ nazýváme \textbf{zobrazením}, pokud splňuje
	$$
	    \forall x \, \forall y_1\, \forall y_2: \left( \left( \left [ c,y_1 \right ] \in F \land \left [ x,y_2 \right ] \in F \right ) \implies y_1 = y_2 \right ) .
	$$
	Řekneme, že zobrazení $F$ je \textbf{zobrazení z množiny $A$ do množiny $B$}, jestliže platí $F\subset A\times B$.
\end{definice}

\begin{definice}
	Nechť $F$ je zobrazení.
	\begin{itemize}
		\item \textbf{Definičním oborem} zobrazení $F$ nazýváme množinu
		$$
		    \mathcal H(F)=\left \{ y; \, \exists x : \left [ x,y \right ] \in F \right \}.
		$$
		Pro $x \in \mathcal D(F)$ označujeme jednoznačně určený prvek $y$ splňující $\left [ x,y \right ] \in F$ symbolem $F(x)$.
		\item \textbf{Oborem hodnot} zobrazení $F$ nazýváme množinu
		$$
		    \mathcal H(f) = \left \{ y;\, \exists x:\left [ x,y \right ] \in F \right \}.
		$$
	\end{itemize}
\end{definice}

\begin{oznaceni}
    Nechť $A$ a $B$ jsou množiny a $F$ je zobrazení.
    \begin{enumerate}[(a)]
	\item Pak symbol $F:A\to B$ znamená, že $F$ je zobrazením z množiny $A$ do množiny $B$
	a $\mathcal D(F) = A$. Takové zobrazení $F$ nazýváme také \textbf{zobrazením množiny $A$ do množiny $B$}.
	\item Pokud je $B=\mathbb R $, pak  místo termínu zobrazení používáme termín \textbf{funkce}.
    \end{enumerate}
\end{oznaceni}

\begin{definice}
	Nechť $A$, $B$ jsou množiny a $f$ je zobrazení.
	\begin{itemize}
		\item \textbf{Obrazem množiny} $A$ pri zobrazení $f$ rozumíme množinu
		$$
  \left \{ y\in \mathcal H(f);\, \exists x \in \mathcal D(f) \cap A: f(x)=y \right \},
		$$
		kterou značíme $f(a)$.
		\item \textbf{Vzorem množiny} $B$ při zobrazení $f$ rozumíme množinu
		$$
		\left \{ x\in \mathcal D(f); \, f(x) \in B \right \},
		$$
		kterou značíme $f^{-1}(B)$.
\end{itemize}
\end{definice}

\begin{definice}
	Řekneme, že zobrazení $f$
	\begin{itemize}
		\item je \textbf{prosté}, jestliže platí
		$$
		    \forall x,y \in \mathcal D(f): f(x)=f(y)\implies x=y,
		$$
		\item je \textbf{na množinu} $B$, jestliže platí $\mathcal H(f)=B$,
		\item je \textbf{bijekcí množiny $A$ na množinu $B$}, jestliže $\mathcal D(f) = A$ a jde o prosté zobrazení na $B$.
\end{itemize}
\end{definice}

\begin{definice}
	Nechť $f$ je zobrazení a $C$ je množina. Pak zobrazení definované předpisem $x\mapsto f(x)$, $x \in C\cap \mathcal D(f)$, nazýváme
	\textbf{restrikcí} nebo \textbf{zúžením} zobrazení $f$ na množinu $C$ a značíme jej $f|_C$.
\end{definice}

\begin{definice}
	NechŤ $f$ a $g$ jsou zobrazení. Pak zobrazení $g\circ f$ je definováno předpisem
	$(g\circ f)(x)=g(f(x))$ pro všechna $x \in \mathcal D(f)$ taková, že $f(x)\in \mathcal D(g)$. Zobrazení
	$g\circ f$ nazýváme \textbf{složeným zobrazením} (\textbf{složením zobrazení}) $f$ a $g$, přičemž
$g$ nazýváme \textbf{vnějším} zobrazením a $f$ nazýváme \textbf{vnitřním} zobrazením.
\end{definice}

\begin{definice}
	Nechť $f:A\to B$ je prosté zobrazení. Pak zobrazení $f^{-1}:f(A)\to A$ definované pro $y\in f(A)$ předpisem $f^{-1}(y)=x$, kde $x\in A$ je jednoznačně určeno vztahem $y=f(x)$, nazýváme \textbf{inverzním zobrazením} k zobrazení $f$.
\end{definice}

\begin{definice}
	Nechť $A$ je neprázdná množina.
	\begin{enumerate}[(a)]
		\item \textbf{Konečnou posloupností} prvků $A$ rozumíme každé zobrazení množiny $\left \{ 1,\dots, n \right \} $, kde
		$n \in \mathbb N $, do množiny $A$. Pokud $k\mapsto a_k$, $k\in \left \{ 1,\dots, k \right \} $, je takové zobrazení,
		pak tuto posloupnost značíme $\left \{ a_k \right \}_{k=1}^n $. Prvek $a_k$ nazýváme $k$-tým \textbf{členem} této posloupnosti.
		\item \textbf{Nekonečnou posloupností} prvků $A$ rozumíme každé zobrazení $n\mapsto a_n$, $n\in \mathbb N $,
		množiny přirozených čísel $\mathbb N $ do množiny $A$. Takovou posloupnost
		obvykle značíme $\left \{ a_n \right \}_{n=1}^\infty $, případně jen $\left \{ a_n \right \} $. Prvek $a_n$
		nazýváme $n-$tým \textbf{členem} této posloupnosti.
	\end{enumerate}
\end{definice}

\subsection{Konečné a spočetné množiny}
\begin{definice}
	\,
	\begin{enumerate}
    \item[(a)] Řekneme, že množina $A$ \textbf{má stejnou mohutnost} jako množina $B$, jestliže existuje
    bijekce $A$ na $B$. Značíme $A\approx B$.
    \item[(b)] Řekneme, že množina $A$ \textbf{má mohutnost menší nebo rovnou} mohutnosti množiny $B$,
    jestliže existuje prosté zobrazení $A$ do $B$. Značíme $A\preceq B$.
    \item[(c)] Řekneme, že množina $A$ \textbf{má menší mohutnost} než množina $B$, jestliže existuje
    prosté zobrazení $A$ do $B$ a přitom $A$ nemá stejnou mohutnost jako $B$. Značíme $A\prec B$.
    \end{enumerate}
\end{definice}

\begin{definice}
Řekneme, že množina $X$ je \textbf{konečná}, pokud je buď prázdná, nebo existuje
$n\in\mathbb{N}$ takové, že $X$ má stejnou mohutnost jako $\{1,\dots,n\}$.
Řekneme, že množina $X$ je \textbf{nekonečná}, pokud není konečná.
Řekneme, že množina $X$ je \textbf{spočetná}, jestliže je konečná, nebo má stejnou mohutnost jako
$\mathbb{N}$. Nekonečná množina, která není spočetná, se nazývá \textbf{nespočetná}.
\end{definice}

\begin{poznamka}
	Pro počet prvků konečné množiny $X$ používáme často značení $|X|$.
	Dvě konečné množiny $X$, $Y$ mají stejnou mohutnost právě tehdy, když $|X|=|Y|$.
\end{poznamka}

\begin{veta}[Cantor-Bernstein] \label{veta1.3}
Nechť $A,B$ jsou množiny takové, že $A\preceq B$ a zároveň $B\preceq A$.
Pak $A$ a $B$ mají stejnou mohutnost.
\end{veta}

\begin{veta}[Cantor] \label{veta1.4}
Nechť $X$ je množina. Pak $X\prec \mathcal{P}(X)$, kde $\mathcal{P}(X)$ je množina všech
podmnožin množiny $X$.
\end{veta}

\begin{veta}[Vlastnosti spočetných množin]\label{veta1.5}
\,
\begin{enumerate}
\item[(a)] Podmnožina spočetné množiny je spočetná.
\item[(b)] Nechť zobrazení $f:A\to\mathbb{N}$ je prosté. Potom je množina $A$ spočetná.
\item[(c)] Sjednocení spočetně mnoha spočetných množin je spočetné.
\item[(d)] Obraz spočetné množiny je spočetná množina.
\item[(e)] Každá nekonečná množina obsahuje nekonečnou spočetnou podmnožinu.
\end{enumerate}
\end{veta}

\subsection{Číselné obory}

Množinu reálných čísel $\mathbb{R}$ lze popsat jako množinu, na níž jsou definovány operace
\textbf{sčítání} a \textbf{násobení} a relace \textbf{uspořádání} $(\le)$, přičemž jsou splněny následující tři skupiny vlastností:

\begin{enumerate}
\item[I.] vlastnosti sčítání a násobení a jejich vzájemný vztah,
\item[II.] vztah uspořádání a operací sčítání a násobení,
\item[III.] vlastnost suprema: \textit{Každá neprázdná shora omezená podmnožina $\mathbb{R}$ má supremum}.
\end{enumerate}

\begin{veta}\label{veta1.6}
Nechť $M\subset\mathbb{R}$ je neprázdná zdola omezená množina.
Pak existuje infimum množiny $M$.
\end{veta}

\begin{veta}\label{veta1.7}
	Ke každému $x \in \mathbb R $ existuje $n \in \mathbb N $ splňující $x<n$.
\end{veta}

\begin{veta}\label{veta1.8}
	Pro každé $r \in \mathbb R $ existuje právě jedno číslo $k\in \mathbb Z $ takové, že $k\leq r<k+1$.
\end{veta}

\begin{veta}\label{veta1.9}
	Nechť $a,b\in \mathbb R$, $a<b$. Pak existuje $q \in \mathbb Q $ takové, že $a<q<b$.
\end{veta}

\begin{definice}
\textbf{Množinu komplexních čísel} $\mathbb C $ definujeme jako množinu
všech uspořádaných dvojic $(a,b)$, kde $a,b \in \mathbb R $, přičemž pro komplexní čísla
$x=(a,b)$, $y=(c,d)$ definujeme operace \textbf{sčítaní} a \textbf{násobení} následovně:
\begin{itemize}
	\item $x+y=(a+c,b+d)$,
	\item $x\cdot y = (ac-bd, ad+bc)$.
\end{itemize}
\par Nechť $x=(a,b)\in \mathbb C $. Prvek $a$ nazýváme \textbf{reálnou částí} $x$,
prvek $b$ nazýváme \textbf{imaginární čásí} $x$. \textbf{Absulotní hodnotou} komplexního čísla $x$
rozumíme $\sqrt{a^2+b^2}$. Dále definujeme $0=(0,0)$, $1=(1,0)$ a $i=(0,1)$.
\textbf{Komplexně sdruženým číslem} k $x$ rozumíme číslo $\bar{x} =(a,-b)$, symbol $-x$
značí číslo $(-a,-b)$ a symbol $1/x$ značí pro $x\ne 0$ (jednoznačně určené)
číslo splňující $x\cdot \frac{1}{x}=1$.
\end{definice}

\section{Limita posloupnosti}

\subsection{Úvod}

\begin{definice}
Řekneme, že posloupnost $\{a_n\}$ je
\begin{itemize}
  \item \textbf{shora omezená}, jestliže množina všech členů této posloupnosti je shora omezená,
  \item \textbf{zdola omezená}, jestliže množina všech členů této posloupnosti je zdola omezená,
  \item \textbf{omezená}, jestliže množina všech členů této posloupnosti je omezená.
\end{itemize}
\end{definice}

\begin{definice}
Řekneme, že posloupnost $\{a_n\}$ je
\begin{itemize}
  \item \textbf{neklesající}, je-li $a_n \le a_{n+1}$ pro každé $n \in \mathbb{N}$,
  \item \textbf{rostoucí}, je-li $a_n < a_{n+1}$ pro každé $n \in \mathbb{N}$,
  \item \textbf{nerostoucí}, je-li $a_n \ge a_{n+1}$ pro každé $n \in \mathbb{N}$,
  \item \textbf{klesající}, je-li $a_n > a_{n+1}$ pro každé $n \in \mathbb{N}$.
\end{itemize}
Posloupnost $\{a_n\}$ je \textbf{monotónní}, pokud splňuje některou z výše uvedených podmínek.
Posloupnost $\{a_n\}$ je \textbf{ryze monotónní}, pokud je rostoucí či klesající.
\end{definice}

\subsection{Vlastní limita posloupnosti}

\begin{definice}
Nechť $\{a_n\}$ je posloupnost reálných čísel a $A \in \mathbb{R}$.
Řekneme, že posloupnost $\{a_n\}$ \textbf{má limitu rovnou} $A$, jestliže platí
\[
\forall \varepsilon \in \mathbb{R}, \ \varepsilon > 0 \ \exists n_0 \in \mathbb{N}
\ \forall n \in \mathbb{N}, \ n \ge n_0 : |a_n - A| < \varepsilon.
\]
\end{definice}

\begin{veta}[Jednoznačnost limity]\label{veta2.1}
Nechť $\{a_n\}$ je posloupnost, která má limitu rovnou $A \in \mathbb{R}$ a zároveň má
limitu rovnou $B \in \mathbb{R}$. Potom platí $A = B$.
\end{veta}

\begin{definice}
Řekneme, že posloupnost $\{a_n\}$ \textbf{konverguje} (je \textbf{konvergentní}), jestliže existuje
$A \in \mathbb{R}$ takové, že $\lim a_n = A$, neboli platí
\[
\exists A \in \mathbb{R} \ \forall \varepsilon \in \mathbb{R}, \ \varepsilon > 0
\ \exists n_0 \in \mathbb{N} \ \forall n \in \mathbb{N}, \ n \ge n_0 :
|a_n - A| < \varepsilon.
\]
Je-li posloupnost konvergentní, říkáme též, že má \textbf{vlastní limitu}.
Jestliže posloupnost nemá vlastní limitu, pak říkáme, že \textbf{diverguje}
(je \textbf{divergentní}).
\end{definice}

\begin{lemma}\label{lemma 2.2}
Nechť $\{a_n\}$ je posloupnost a $A \in \mathbb{R}$. Pak $\lim a_n = A$ právě tehdy,
když existuje $K \in \mathbb{R}$, $K > 0$, takové, že platí
\[
\forall \varepsilon \in \mathbb{R}, \ \varepsilon > 0 \ \exists n_0 \in \mathbb{N}
 \forall n \in \mathbb{N}, \ n \ge n_0 : |a_n - A| < K\varepsilon.
\]
\end{lemma}

\begin{veta} \label{veta2.3}
Nechť $\lim a_n = A \in \mathbb{R}$. Potom $\lim |a_n| = |A|$.
\end{veta}

\begin{veta}\label{veta2.4}
Nechť $\{a_n\}$ je konvergentní posloupnost. Potom je $\{a_n\}$ omezená.
\end{veta}

\begin{definice}
Nechť $\{a_n\}$ je posloupnost a $\{n_k\}_{k=1}^{\infty}$ je rostoucí posloupnost
přirozených čísel. Pak $\{a_{n_k}\}_{k=1}^{\infty}$ nazýváme \textbf{vybranou posloupností}
z posloupnosti $\{a_n\}$, případně \textbf{podposloupností} posloupnosti $\{a_n\}$.
\end{definice}

\begin{veta}\label{veta2.5}
Nechť $\{a_n\}$ je posloupnost, $\lim a_n = A \in \mathbb{R}$ a
$\{n_k\}_{k=1}^{\infty}$ je rostoucí posloupnost přirozených čísel.
Potom $\lim_{k \to \infty} a_{n_k} = A.$
\end{veta}

\begin{veta}[Aritmetika limit] \label{veta2.6}
Nechť $\{a_n\}$ a $\{b_n\}$ jsou posloupnosti, $\lim a_n = A \in \mathbb{R}$ a
$\lim b_n = B \in \mathbb{R}$. Potom platí:
\begin{itemize}
  \item[(a)] $\lim (a_n + b_n) = A + B$,
  \item[(b)] $\lim (a_n \cdot b_n) = A \cdot B$,
  \item[(c)] je-li $B \ne 0$ a pro každé $n \in \mathbb{N}$ platí $b_n \ne 0$, pak
  \[
  \lim \frac{a_n}{b_n} = \frac{A}{B}.
  \]
\end{itemize}
\end{veta}

\begin{veta}\label{veta2.7}
Nechť $\{a_n\}$ a $\{b_n\}$ jsou posloupnosti, $\lim a_n = 0$ a $\{b_n\}$ je omezená.
Pak $\lim a_n b_n = 0.$
\end{veta}

\begin{veta}[Limita a uspořádání] \label{veta2.8}
Nechť $A, B \in \mathbb{R}$, $\{a_n\}$ a $\{b_n\}$ jsou posloupnosti splňující
$\lim a_n = A$ a $\lim b_n = B$.
\begin{itemize}
  \item[(a)] Nechť $A < B$. Potom existuje $n_0 \in \mathbb{N}$ takové, že pro každé
  $n \in \mathbb{N}$, $n \ge n_0$, platí $a_n < b_n$.
  \item[(b)] Nechť existuje $n^\ast \in \mathbb{N}$ takové, že pro každé
  $n \in \mathbb{N}$, $n \ge n^\ast$, platí $a_n \ge b_n$. Potom $A \ge B$.
\end{itemize}
\end{veta}

\begin{veta}[O dvou strážnících] \label{veta2.9}
Nechť $\{a_n\}$, $\{b_n\}$ a $\{c_n\}$ jsou posloupnosti splňující:
\begin{itemize}
  \item[(a)] existuje $n_0 \in \mathbb{N}$ takové, že pro každé $n \in \mathbb{N}$,
  $n \ge n_0$, platí $a_n \le c_n \le b_n$,
  \item[(b)] $\{a_n\}$ a $\{b_n\}$ jsou konvergentní a platí
  $\lim a_n = \lim b_n$.
\end{itemize}
Potom je $\{c_n\}$ konvergentní a platí
$\lim c_n = \lim a_n$.
\end{veta}

\subsection{Nevlastní limita posloupnosti}

\begin{definice}
Řekneme, že posloupnost $\{a_n\}$ \textbf{má limitu rovnou} $\infty$ (čteme plus nekonečno),
jestliže
\[
\forall K \in \mathbb{R} \ \exists n_0 \in \mathbb{N}
\ \forall n \in \mathbb{N}, \ n \ge n_0 : a_n > K.
\]
Řekneme, že posloupnost $\{a_n\}$ \textbf{má limitu rovnou} $-\infty$ (čteme mínus nekonečno),
jestliže
\[
\forall K \in \mathbb{R} \ \exists n_0 \in \mathbb{N}
\ \forall n \in \mathbb{N}, \ n \ge n_0 : a_n < K.
\]
\par Má-li posloupnost limitu rovnou plus nebo mínus nekonečnu, říkáme, že \textbf{má nevlastní
limitu}. Jestliže má posloupnost limitu rovnou $\infty$, pak říkáme, že \textbf{diverguje} k
$\infty$. Jestliže má posloupnost limitu rovnou $-\infty$, pak říkáme, že \textbf{diverguje}
k $-\infty$.
\end{definice}

\begin{veta}[Jednoznačnost limit podruhé]
Nechť $\left \{ a_n \right \} $ je posloupnost, která má limitu rovnou
$A\in \mathbb R ^*$ a zároveň má limitu rovnou $B\in \mathbb R ^*$. Potom platí $A= B$.
\end{veta}

\begin{poznamka}
	Označení. Symbol $\lim_{n\to \infty} a_n$ bude značit vlastní nebo nevlastní limitu posloupnosti
	$\left \{ a_n \right \} $, pokud existuje. Někdy budeme používat jen symbol $\lim a_n$. Někdy také píšeme
	$a_n \to A, n\to \infty$.
\end{poznamka}

\begin{poznamka}
	$$
	\lim a_n \begin{cases}
		\textrm{existuje} \begin{cases}
			\textrm{je vlastní, tj. } \lim a_n\in \mathbb R \\
			\textrm{je nevlastní, tj. } \lim a_n = \begin{cases}
				\infty \\
				-\infty
			\end{cases}
		\end{cases}\\
		\textrm{neexistuje}
		\end{cases}
	$$
	Pokud má vlastní limitu, říkáme, že konverguje. V ostatních případech diverguje.
\end{poznamka}

\begin{veta}[Aritmetika limit podruhé]
    Necht’ $\left \{ a_n \right \} $ a $\left \{ b_n \right \} $ jsou posloupnosti, $ \lim a_n = A \in \mathbb R ^*$ a
    $\lim b_n = B \in \mathbb R ^*$. Potom platí:
    \begin{enumerate}[(a)]
    	\item $\lim (a_n + b_n) = A + B$, pokud je výraz na pravé straně definován,
    	\item $\lim (a_n\cdot b_n) = A\cdot B$, pokud je výraz na pravé straně definován,
    	\item je-li $b_n \ne 0$ pro všechna $n \in N$, pak $\lim \frac{a_n}{b_n} = \frac{A}{B}$, pokud
    	je výraz na pravé straně definován.
    \end{enumerate}
\end{veta}

\begin{proof}
	Dokážeme případ $(a)$, ostatní analogicky. Pokud $A,B\in \mathbb R $, pak znění plyne z věty \ref{veta2.6}.
	Díky komutativitě stačí dokázat tvrzení v následujících případech:
	\begin{enumerate}[(i)]
		\item $A=B=\infty$,
		\item $A=\infty, B\in \mathbb R$,
		\item $A\in \mathbb R, B=-\infty$,
		\item $A = B = -\infty$.
	\end{enumerate}
	Provedeme pouze důkaz případu (ii). Zvolme $K\in \mathbb R $. Nalezneme $n_1,n_2\in \mathbb N $ splňující
	\begin{align*}
	    &\forall n \in \mathbb N, n \geq n_1: |b_n - B| < 1,\\
		&\forall n\in \mathbb N , n \geq n_2: a_n > K - B + 1.
	\end{align*}
	Položme $n_0 = \max \left \{ n_1,n_2 \right \} $. Potom pro každé $n\in \mathbb N ,n\geq n_0,$ platí
	$$
	    a_n + b_n > (K-B+1) + (B-1) = K.
	$$
\end{proof}

\begin{poznamka}
	U zkoušky budou vyžadovány i zbývající části tohoto důkazu.
\end{poznamka}

\begin{priklad}\,
\begin{enumerate}[(1)]
	\item $\lim n^2 = \lim(n\cdot n)=\infty \cdot \infty = \infty$,
	\item Nechť $a_n=n+\alpha$, $b_n = -n$. Pak $\lim a_n = \infty$, $\lim b_n = -\infty$, $\lim(a_n+ b_n)=\alpha$,
	neboť je to konstantní posloupnost, $\alpha \in \mathbb R $.
\end{enumerate}
\end{priklad}

\begin{poznamka}
	Platí následující. Nechť $\left \{ a_n \right \} $ a $\left \{ c_n \right \} $ jsou posloupnosti
	splňující
	\begin{enumerate}[(a)]
		\item existuje $n_0\in \mathbb N $ takové, že pro každé $n\in \mathbb N , n\geq n_0$ platí
		$a_n\leq c_n$,
		\item $\lim a_n=\infty$.
	\end{enumerate}
	Potom $\lim c_n = \infty$.
\end{poznamka}

\begin{veta}\label{veta2.12}
Nechť $\left \{ a_n \right \} $ a $\left \{ b_n \right \} $ jsou posloupnosti splňující:
\begin{itemize}
	\item pro každé $n \in \mathbb N $ platí $b_n \ne 0$,
	\item $\lim a_n = A\in \mathbb R ^*$ a $A>0$,
	\item $\lim b_n = 0$,
	\item existuje $n_0 \in \mathbb N $ takové, že pro každé $n \in \mathbb N , n \geq n_0$, platí $b_n > 0$.
\end{itemize}
Potom $\lim \frac{a_n}{b_n}=\infty$.
\end{veta}

\begin{proof}
	Rozlišíme dvě možnosti:
	\begin{enumerate}[(1)]
		\item $A\in \mathbb R , A>0$. Zvolme $K\in \mathbb R $. Položme $\varepsilon = \frac{1}{2}A$.
		Potom $\varepsilon \in \mathbb R , \varepsilon > 0$. Nalezneme $n_1 \in \mathbb N $ takové, že
		$$
		    \forall n\in \mathbb N , n \geq n_1: |a_n - A|< \varepsilon = \frac{1}{2}A,
		$$
		tj. $a_n \in \left( \frac{1}{2}A, \frac{3}{2}A \right ) $. Položme $L = \max \left \{ K,1 \right \} $.
		Nalezneme $n_2\in \mathbb N $ takové, že
		$$
		    \forall n \in \mathbb N , n \geq n_2: |b_n - 0| < \frac{\frac{1}{2}A}{L},
		$$
		tj. $b_n \in \left( -\frac{A}{2L}, \frac{A}{2L} \right ) $. Položme
		$n^* = \max \left \{ n_0, n_1, n_2 \right \} $. Potom pro každé $n\in \mathbb N , n \geq n^*$ platí
		$$
		\frac{a_n}{b_n}>\frac{\frac{1}{2}A}{b_n} > \frac{\frac{1}{2}A}{\frac{A}{2L}} = L \geq K.
		$$
		\item $A = \infty$. Zvolme $K\in \mathbb R $. Nalezneme $n_1\in \mathbb N $ takové, že
		$$
		    \forall n \in \mathbb N , n \geq n_1: a_n > 1.
		$$
		Položme $L = \max \left \{ K,1 \right \} $. Nalezneme  $n_2\in \mathbb N $ takové, že
		$$
		    \forall n \in \mathbb N , n \geq n_2: |_n - 0| < \frac{1}{L}.
		$$
		Položme $n^* = \max \left \{ n_0, n_1, n_2 \right \} $. Potom pro každé $n \in \mathbb N , n \geq n^*$
		platí
		$$
            \frac{a_n}{b_n} > \frac{1}{b_n} > \frac{1}{\frac{1}{L}} = l \geq K,
		$$
		což jsme chtěli dokázat. \qedhere
	\end{enumerate}
\end{proof}

\begin{poznamka}[Supremum a infimum v $\mathbb R ^*$]
	Musíme rozlišovat $(\mathbb R , \leq)$ a $(\mathbb R ^*, \leq)$. V $\mathbb R ^*$
	platí $\sup \emptyset = -\infty$ a podobně $\inf \emptyset  = \infty$. Je-li $M \subset \mathbb R $
	a $M$ není shora omezená, je $\sup M = \infty$ a podobně je-li $M\subset \mathbb R $ a $M$ není zdola omezená,
	je $\inf M = -\infty$.
\end{poznamka}

\subsection{Hlubší věty o limitách}
\begin{veta}\label{v213} % 2.13
	Nechť $\left \{ a_n \right \} $ je monotónní posloupnost. Potom existuje $\lim a_n$.
	Je-li $\left \{ a_n \right \} $ neklesající, pak $\lim a_n = \sup \left \{ a_n, n \in \mathbb N  \right \} $.
	Je-li $\left \{ a_n \right \} $ nerostoucí, pak $\lim a_n = \inf \left \{ a_n, n \in \mathbb N  \right \} $.
\end{veta}

\begin{proof}
	Položme $A=\sup \left \{ a_n, n\in \mathbb N  \right \} $ a předpokládejme, že $\left \{ a_n \right \} $
	je neklesající. Rozlišíme následující možnosti:
	\begin{enumerate}[(i)]
		\item $\left \{ a_n \right \} $ není shora omezená. Potom $A = \infty$. Zvolme $K\in \mathbb R $.
		Nalezneme $n_0\in \mathbb N $ takové, že $a_{n_0} > K$. Potom pro každé $n\in \mathbb N , n\geq n_0$
		platí
		$$
		    a_n \geq a_{n_0} > K.
		$$
		\item $\left \{ a_n \right \} $ je shora omezená. Potom $A\in \mathbb R $. Zvolme
		$\varepsilon \in \mathbb R ,\varepsilon > 0$. Číslo $A-\varepsilon$, které je menší než $A$, není
		horní závorou množiny $\left \{ a_n, n\in \mathbb N  \right \} $. Nalezneme $n_0 \in \mathbb N $
		takové, že $a_{n_0} > A-\varepsilon$. Potom pro každé $n \in \mathbb N , n \geq n_0$ platí
		$$
		    A - \varepsilon < a_{n_0}\leq a_n \leq A < A+ \varepsilon.
		$$
		Odtud plyne $\lim a_n = A$.
	\end{enumerate}
	Pokud je $\left \{ a_n \right \} $ nerostoucí, lze postupovat analogicky, ale supremum je nahrazeno
	infimem. Je možné také přejít k posloupnosti $\left \{ -a_n \right \} $.
\end{proof}

\begin{poznamka}
	Z důkazu předchozí věty plyne, že pro monotónní posloupnost $\left \{ a_n \right \} $ platí:
	$$
	    \left \{  a_n \right \} \textrm{ je } \begin{cases}
			\textrm{neklesající} \begin{cases}
				\textrm{je shora omezená, } & \textrm{pak } \lim a_n \in \mathbb R  \\
				\textrm{není shora omezení, } & \textrm{pak }\lim a_n = \infty
			\end{cases}\\
			\textrm{nerostoucí} \begin{cases}
				\textrm{je zdola omezená, }& \textrm{pak }\lim a_n \in \mathbb R\\
				\textrm{není zdola omezená, }& \textrm{pak }\lim a_n = -\infty
			\end{cases}
		\end{cases}
	$$
\end{poznamka}

\begin{priklad}
	Nechť $q\in (0,1)$. Potom platí $\lim_{n\to \infty} q^n = 0.$ Pro každé $n\in \mathbb N $ platí
	$q^n \geq q^{n+1}\geq 0$. Podle věty \ref{v213} existuje $A\in \mathbb R $ takové, že
	$\lim_{n\to \infty} q^n=A$. Podle věty \ref{veta2.8}
	platí $A\geq 0$.
	$$
	    A= \lim_{n\to\infty} q^{n+1} = \lim_{n\to\infty} q\cdot q^n = q\cdot A,
	$$
	odtud plyne $A=0$. Použili jsme větu \ref{veta2.5}.
	Pro úplnost uveďme ještě toto:
	$$
	    \lim_{n\to\infty} q^n \begin{cases}
			\infty & \textrm{pro } q >1, \\
			1 & \textrm{pro } q =1, \\
			0 & \textrm{pro } q \in (-1,1), \\
			\textrm{neexistuje} & \textrm{pro } q \leq 1.
		\end{cases}
	$$
	\begin{enumerate}[(a)]
		\item $\lim_{n \to \infty} q^n = \lim_{n \to \infty} \frac{1}{\left( \frac{1}{q} \right )^n } = \infty$ dle věty \ref{veta2.12}
		\item $\lim_{n \to \infty} q^n = \lim_{n \to \infty} 1=1$
		\item Pro každé $n\in \mathbb N $ platí $- |q|^n\leq q^n \leq |q|^n$, což konverguje k nule, takže i
		$- |q|^n$ konverguje k nule, a tedy $\lim_{n \to \infty} q^n = 0.$
		\item $\lim_{n \to \infty} q^{2n}=\lim_{n \to \infty} \left( q^2 \right )^n = \begin{cases}
			1 & q = -1,\\
			\infty & q < -1.
		\end{cases} $ Podobně $\lim_{n \to \infty} q^{2n+1} = \lim_{n \to \infty} q\cdot \left( q^2 \right )^n = \begin{cases}
			-1 & q = -1,\\
			-\infty & q < -1
		\end{cases}$
	\end{enumerate}
	Věta \ref{veta2.5} dává tvrzení.
\end{priklad}

\begin{poznamka}[Varování]
	V předchozím výpočtu je krok ověření existence limity podstatný.
	Mějme posloupnost $a_1 = -1, a_{n+1} = -a_n$. Pak
	\begin{align*}
	\lim a_{n+1} &= \lim (-a_n) = -A\\
	& = \lim a_n = A\\
		& \implies \lim a_n = 0,
	\end{align*}
	ale tato limita neexistuje! Tento výpočet je špatně.
\end{poznamka}

\begin{veta}[Bolzano--Weierstrass]\label{bolzano-weierstrass} %2.14
	Nechť $\left \{ a_n \right \} $ je posloupnost. Potom existuje vybraná posloupnost
	$\left \{ a_{n_k} \right \}_{k=1}^\infty $, která je konvergentní.
\end{veta}

\begin{proof}
	Nalezneme posloupnosti $\left \{ \alpha_k \right \}, \left \{ \beta_k \right \}  $ splňující:
	\begin{enumerate}[(a)]
		\item pro každé $k\in \mathbb N $ platí $\alpha_k \leq \beta_k$,
		\item pro každé $k \in \mathbb N $ platí $\left [ \alpha_{k+1}, \beta_{k+1} \right ] =
		\left [ \alpha_k, \frac{1}{2}\left( \alpha_k + \beta_k \right )  \right ] $ nebo
		$\left [ \alpha_{k+1}, \beta_{k+1} \right ] =
		\left [\frac{1}{2}\left( \alpha_k + \beta_k \right ), \beta_k  \right ] $,
		\item pro každé $k \in \mathbb N $ platí, že množina
		$\left \{ j\in \mathbb N ,a_j \in \left [ \alpha_k, \beta_k \right ]  \right \} $ je
		nekonečná.
	\end{enumerate}
	\paragraph{Konstrukce $\left \{ \alpha_k \right \}, \left \{ \beta_k \right \}  $.} Posloupnost
	$\left \{ a_n \right \} $ je omezená, takže nalezneme $\alpha_1, \beta_1 \in \mathbb R $ splňující
	$$
	    \forall n \in \mathbb N : \alpha_1 \leq a_n \leq \beta_1.
	$$
	Postupujeme indukcí. Předpokládejme, že již máme $\alpha_k, \beta_k$. Množina
	$$
	    S = \left \{ j \in \mathbb N , a_j \in \left [ \alpha_k, \beta_k \right ]  \right \}
	$$
	je nekonečná. Dále definujme
	\begin{align*}
		L &= \left \{ j\in \mathbb N , a_j \in \left [ \alpha_k, \frac{1}{2}\left( \alpha_k, \beta_k \right )  \right ]  \right \},  \\
		P &= \left \{ j \in \mathbb N , a_j \in \left [ \frac{1}{2}\left( \alpha_k, \beta_k \right ), \beta_k  \right ]  \right \}.
	\end{align*}
	Platí $S$ je nekonečná a $S = L\cup P.$ Potom $L$ nebo $P$ je nekončená.
	Pokud je $L$ nekonečná, klademe $\alpha_{k+1}=\alpha_k$, $\beta_{k+1}=\frac{1}{2}\left( \alpha_k + \beta_k \right ) $.
	Pokud je $L$ konečná, potom je $P$ nekonečná a klademe $\alpha_{k+1}=\frac{1}{2}\left( \alpha_k+\beta_k \right ) $,
	$\beta_{k+1}=\beta_k$.
	\paragraph{Limita pomocných posloupností.} Z konstrukce plyne, že $\left \{ a_k \right \} $ je neklesající a podobně
	$\left \{ \beta_k \right \} $ je nerostoucí. Dále platí, že
	$$
	    \forall k \in \mathbb N: \alpha_1 \leq \alpha_k \leq \beta_k \leq \beta_1.
	$$
	Máme dvě monotónní posloupnosti, které jsou omezené, takže mají podle věty \ref{v213} limity
	\begin{align*}
		\lim_{n \to \infty} \alpha_k &= \alpha \in \left [ \alpha_1, \beta_1 \right ],  \\
		\lim_{n \to \infty} \beta_k &= \beta \in \left [ \alpha_1, \beta_1 \right ].
	\end{align*}
	Platí
	$$
        \beta - \alpha=\lim_{n \to \infty} \left( \beta_k - \alpha_k \right ) = \lim_{n \to \infty} \left( \beta_1 - \alpha_1 \right ) \cdot 2^{-k+1} = 0
	$$
	Odtud plyne $\alpha = \beta$.
	\paragraph{Konstrukce vybrané posloupnosti.} Položme $n_1=1$. Konstrukce indukcí. Předpokládejme, že
	již máme přirozená čísla $n_1 < n_2 < \dots < n_k$ a pro každé $j\in \mathbb N ,j\leq k$
	platí $a_{n_j} \in \left [ \alpha_j,\beta_j \right ] $. Budeme hledat $n_{k+1}$. Množina
	$\left \{ i\in \mathbb N , a_i \in \left [ \alpha_{k+1}, \beta_{k+1} \right ]  \right \} $
	je nekonečná podle (c). Potom i množina
	$$\left \{ i \in \mathbb N , i > n_0 \land a_i \in \left [ \alpha_{k+1}, \beta_{k+1} \right ]  \right \} $$
	je nekonečná. Máme tedy nalézt $n_{k+1} \in \mathbb N $ takové, že $n_{k+1} > n_k$ a
	$a_{n_{k+1}}\in \left [ \alpha_{k+1}, \beta_{k+1}\right ] $. Potom pro každé $k \in \mathbb N $ platí
	$$
	    \alpha_k \leq a_{n_k} \leq \beta_k,
	$$
	ale $\alpha_k$ i $\beta_k$ konvergují k $\alpha$, takže podle věty o dvou strážnících platí
	$$
        \lim_{n \to \infty} a_{n_k} = \alpha\in \mathbb R,
	$$
	což jsme chtěli dokázat.
\end{proof}

\begin{definice}\label{limsup}
	Nechť $\{a_n\}$je posloupnost reálných čísel. Pak definujeme
	$$
	    \limsup_{n\to \infty} a_n = \begin{cases}
			\lim_{n\to\infty} \sup \left \{ a_k, k \geq n \right \}, & \textrm{jestliže je } {a_n} \textrm{ shora omezená},\\
			\infty, & \textrm{jestliže není } \left \{ a_n \right \} \textrm{ shora omezená}.
		\end{cases}
	$$
	Tuto hodnotu nazýváme \textbf{limes superior} posloupnosti $\left \{ a_n \right \} $.
	Obdobně definujeme \textbf{limes inferior} posloupnosti $\left \{ a_n \right \} $ předpisem
	$$
	    \liminf_{n\to \infty} a_n = \begin{cases}
			\lim_{n\to\infty} \inf \left \{ a_k, k \geq n \right \}, & \textrm{jestliže je } {a_n} \textrm{ zdola omezená},\\
			-\infty, & \textrm{jestliže není } \left \{ a_n \right \} \textrm{ zdola omezená}.
		\end{cases}
	$$
\end{definice}

\begin{poznamka}\,
    \begin{enumerate}[(i)]
    	\item Definice \ref{limsup} je korektní díky větě o limitě monotónní posloupnosti \ref{v213}.
        \item Místo $\limsup_{n\to\infty} a_n$ někdy píšeme $\overline{\lim_{n\to\infty}}$ a místo $\liminf a_n$
        $\underline{\lim_{n\to\infty}}$.
    \end{enumerate}
\end{poznamka}

\begin{priklad}\,
\begin{enumerate}[(i)]
	\item $\limsup_{n\to\infty} (-1)^n = 1$, $\liminf_{n\to\infty} = -1$.
	\item $\limsup_{n\to\infty} (-n) = -\infty$, $\liminf_{n\to\infty} (-n) = -\infty$.
\end{enumerate}
\end{priklad}

\begin{veta}[Limita, limes superior a limes inferior]\label{veta2.15}
	Nechť $\left \{ a_n \right \}$ je posloupnost a $A\in \mathbb R ^*$. Potom
$\lim a_n = A$ právě tehdy, když $\limsup a_n = \liminf a_n = A$.
\end{veta}

\begin{proof}\,
\begin{enumerate}[(i)]
	\item[$\implies$] \begin{itemize}
		\item $A \in \mathbb R $: Podle věty \ref{veta2.4}
		$\left \{ a_n \right \} $ omezená. Označme $b_n = \sup \left \{ a_n, k \geq n \right \}$,
		$c_n = \inf \left \{ a_k, k \geq n \right \} $. Zvolme $\varepsilon \in \mathbb R , \varepsilon> 0.$
		Nalezneme $n_0 \in \mathbb N $ takové, že
		$$
		\forall n \in \mathbb N, n \geq n_0: |a_n- A| < \varepsilon.
		$$
		Potom platí
		$$
		    \forall n\in \mathbb N, n \geq n_0: A - \varepsilon < a_n < A + \varepsilon.
		$$
		Odtud plyne, že pro každé $n \in \mathbb N, n \geq n_0$ platí
		$$
		    A - \varepsilon \leq c_n \leq b_n \leq A+ \varepsilon.
		$$
		Podle věty \ref{veta2.8}
		plyne, že
		$$
		    A- \varepsilon \leq \underbrace{\lim c_n}_{\liminf a_n} \leq \underbrace{\lim b_n}_{\limsup a_n} \leq A+\varepsilon,
		$$
		odkud plyne $\limsup_{n\to\infty} a_n = \limsup_{n\to\infty} a_n = A.$
		\item $A = \infty$: $\left \{ a_n \right \} $ není omezená shora, ale je omezená zdola.
		Dostáváme $\limsup a_n = \infty$ přímo z definice. Zvolme $K\in \mathbb R $. Nalezneme $n_0 \in \mathbb N$
		takové, že
		$$
		    \forall n \in \mathbb N , n \geq n_0: a_n > K + 1.
		$$
		Potom pro každé $n \in \mathbb N, n \geq n_0$ platí
		$$
		    c_n \geq K + 1 > K.
		$$
		Je tedy $\liminf a_n = \lim c_n = \infty$.
		\item $A = -\infty$: Proveďte samostatně.
\end{itemize}
\item[$\impliedby$] \begin{itemize}
	\item $A \in \mathbb R $: $\limsup a_n = \liminf a_n = A$. Zde máme, že $\left \{ a_n \right \} $
	je omezená posloupnost. Potom pro každé $n \in \mathbb N $ platí $b_n, c_n\in \mathbb R $ a
	$$
	    c_n \leq a_n \leq b_n.
	$$
	Protože $c_n$ konverguje k $\liminf a_n = A$ a $b_n$ k $\limsup a_n = A$, je podle věty o dvou strážnících
	\ref{veta2.9}
	$\lim a_n = A$.
	\item $A = \infty$: $\left \{ a_n \right \} $ není shora omezená, ale je zdola omezená.
	Potom pro každé $n \in \mathbb N $ platí $c_n \in \mathbb R $ a platí $c_n \leq a_n$.
	Protože $c_n$ konverguje k $\liminf a_n = \infty$, je podle věty
	o dvou strážnících \ref{veta2.9}
	i $\lim a_n = \infty$.
	\item $A = -\infty$: Samostatně. \qedhere
\end{itemize}
\end{enumerate}
\end{proof}

\begin{veta}\label{veta2.16}
	Nechť $\left \{ x_n \right \} $ a $\left \{ y_n \right \} $ jsou postloupnosti. Předpokládejme, že platí
	$$
	    \exists n_0 \in \mathbb N \, \forall n \in \mathbb N , n \geq n_0: x_n \leq y_n.
	$$
	Potom $\limsup x_n \leq \limsup y_n$ a $\liminf x_n \leq \liminf y_n$.
\end{veta}

\begin{proof}
	Dokážeme jen druhou nerovnost, první se provede obdobně. Pokud $\left \{ x_n \right \} $ není zdola
	omezená, potom $\liminf x_n = -\infty$ a dokazovaná nerovnost zřejmě platí. Předpokládejme tedy,
	že $\left \{ x_n \right \} $ je zdola omezená. Potom je zdola omezená i posloupnost $\left \{ y_n \right \}$.
	Potom pro každé $n \in \mathbb N , n \leq n_0$ platí
	$$
	    \inf \left \{ x_k, k \geq n \right \} \leq \inf \left \{ y_k, k \geq n \right \}.
	$$
	Protože levá strana konverguje k $\liminf x_n$ a pravá k $\liminf y_n$, musí platit i
	$\liminf x_n \leq \liminf y_n$.
\end{proof}

\begin{definice}
Nechť $\{a_n\}$ je posloupnost.
Řekneme, že $A \in \mathbb R ^*$ je \textbf{hromadná hodnota} posloupnosti $\{a_n\}$,
jestliže existuje vybraná posloupnost $\{a_{n_k}\}_{k=1}^{\infty}$
z posloupnosti $\{a_n\}$ taková, že platí
$\lim_{k\to \infty} a_{n_k} = A.$ Množinu všech hromadných hodnot posloupnosti ${a_n}$
značíme $H \left( \left \{ a_n \right \}  \right ) $.
\end{definice}

\begin{priklad}\,
\begin{enumerate}[(i)]
	\item $H \left( \left \{ (-1)^n \right \}  \right )  = \left \{ -1,1 \right \} $
	\item $H \left( \left \{ \sin n \right \}  \right )  = \left [ -1,1 \right ] $
\end{enumerate}
\end{priklad}

\begin{veta}[Limes superior, limes inferior a hromadné hodnoty]\label{veta2.17}
	Nechť $\left \{ a_n \right \} $ je posloupnost. Potom $\limsup a_n, \liminf a_n \in H \left( \left \{ a_n \right \}  \right ) $
	a pro každé $y \in H \left( \left \{ a_n \right \}  \right ) $ platí $\liminf a_n \leq y \leq \limsup a_n$.
\end{veta}

\begin{proof}
	Označme $A = \limsup a_n$. Rozlišíme následující případy:
	\begin{enumerate}[(i)]
		\item $A \in \mathbb R $: Označme $b_n = \sup \left \{ a_k, k \geq n \right \} $.
		Potom platí $b_n \in \mathbb R $ pro každé $n \in \mathbb N $ a
		$\left \{ b_n \right \} $ je nerostoucí posloupnost splňující $\lim b_n = A$. Nalezneme
		rostoucí posloupnost přirozených čísel $\left \{ n_k \right \}_{k=1}^\infty $
		a rostoucí posloupnost přirozených čísel $\left \{ m_k \right \}_{k=1}^\infty $
		takové, že $m_k \leq n_k$ a
		\begin{equation}\label{rovnicelimsup}
		    \left | b_{m_n} - a_{n_k}\right |< \frac{2}{k}.
		\end{equation}
		\begin{itemize}
			\item $k = 1$: Položme $m_1 = 1$. Podle definice je $b_{m_1} = \sup \left \{ a_k, k \geq m_1 \right \}$.
			Nalezneme $n_1 \geq m_1$ splňující $b_{m_1} \geq a_{n_1} > b_{m_1} - \frac{2}{k}$.
			\item Indukční krok z $k$ na $k+1$: Máme
			\begin{align*}
				n_1 < n_2 < n_3 < \dots & < n_k, \\
				m_1 < m_2 < m_3 < \dots & < m_k.
			\end{align*}
			Nalezneme $m_{k+1}\in \mathbb N $ takové, že $m_{k+1}>n_k$:
			$$
			    b_{m_{k+1}} = \sup \left \{ a_j, j \geq m_{k+1} \right \}.
			$$
			Nalezneme
			$$
			    n_{k+1} >\geq a_{n_{k+1}} > b_{m_{k+1}} - \frac{2}{k+1}.
			$$
			Pak podle \ref{rovnicelimsup} je
			$$
			    |b_{m_{k+1}} - a_{m_{k+1}}| < \frac{2}{k+1}.
			$$
			platí $\lim_{k\to\infty} b_{m_k} = A = \lim_{k\to\infty} \left( b_{m_k}-\frac{2}{k+1} \right ) $,
			takže $\lim a_{n_k} = A.$
		\end{itemize}
		\item $A = \infty$: $\left \{ a_n \right \} $ není shora omezená.
		\begin{itemize}
			\item $k = 1$: Nalezneme $n_1 \in \mathbb N $ takové, že $a_{n_1} > 1$.
			\item Indukční krok z $k$ na $k+1$: Máme $n_1 < n_2 < n_3 < \dots < n_k$.
			Množina $\left \{ a_j, j > n_k \right \}$ není shora omezená, takže nalezneme
			$n_{k+1} \in \mathbb N , n_{k+1} > n_k$ takové, že $a_{n_{k+1}} > k+1$.
			Pro každé $k \in \mathbb N $ platí $a_{n_k} > k$, ale $k$ diverguje do $\infty$, takže
			podle věty o dvou strážnících \ref{veta2.9}
			$\lim a_{n_k} = \infty$.
		\end{itemize}
		\item $A = -\infty$: $\left \{ a_n \right \} $ není zdola omezená. Máme tedy $\liminf a_n = -\infty$.
		Podle věty o dvou strážnících \ref{veta2.9}
		je $\lim a_n = -\infty$, a tedy $-\infty \in H \left(  \left \{ a_n \right \}  \right ) $.
	\end{enumerate}
	\paragraph{Důkaz nerovnosti pro $\liminf a_n$.}
	Označme $y\in H \left( \left \{ a_n \right \}  \right ) $. Pokud $\liminf a_n = -\infty$,
	pak $\liminf a_n = -\infty \leq y$. Předpokládejme tedy, že $\liminf a_n < -\infty$. Potom
	platí
	$$
	    \liminf a_n = \lim_{n\to \infty} \inf \left \{ a_k, k\geq n \right \}.
	$$
	Nalezneme rostoucí posloupnost přirozených čísel $\left \{ n_k \right \}_{k=1}^\infty $ takovou,
	že $\lim_{n\to \infty} a_{n_k}=y$. Pro každé $k \in \mathbb N $ platí
	$$
	    \inf \left \{ a_j, j\geq n_k \right \} \leq a_{n_k}.
	$$
	Protože $a_{n_k}$ konverguje k $y$ a $\inf \left \{ a_j, j\geq n_k \right \}$ konverguje
	k $\liminf a_n$, platí $\liminf a_n \leq y$, což jsme chtěli dokázat.
	Pro limes superior postupujeme obdobně.
\end{proof}

\begin{dusledek}
	Nechť $\left \{ a_n \right \} $ je posloupnost. Potom je $H \left( \left \{ a_n \right \}  \right ) $
	neprázdná.
\end{dusledek}

\begin{proof}
	$\limsup a_n \in H \left( \left \{ a_n \right \}  \right ) $
\end{proof}

\begin{definice}[Bolzanova--Cauchyova podmínka]
	Řekneme, že posloupnost $\left \{ a_n \right \} $ splňuje \textbf{Bolzanovu--Cauchyovu podmínku},
	jestliže
	$$
	    \forall \varepsilon \in \mathbb R, \varepsilon > 0 \, \exists n_0 \in \mathbb N\, \forall m,n \in \mathbb N , n\geq n_0, m\geq n_0: |a_n-a_m| < \varepsilon.
	$$
\end{definice}

\begin{veta} \label{veta2.18}
	Posloupnost $\left \{ a_n \right \} $ má vlastní limitu právě tehdy, když splňuje Bolzanovu--Cauchyovu podmínku.
\end{veta}

\begin{proof}\,
\begin{enumerate}
	\item[$\implies$] Předpokládáme, že $\lim a_n = A\in \mathbb R $. Zvolme $\varepsilon \in \mathbb R ,\varepsilon > 0$.
	K němu nalezneme $n_0 \in \mathbb N $ takové, že
	\begin{equation}\label{rovnicebolzano}
	\forall n \in \mathbb N , n \geq n_0: |a_n - A| < \frac{\varepsilon}{2}.
\end{equation}
	Potom pro každé $m,n \in \mathbb N , m,n \geq n_0$ platí
	$$
	    |a_n - a_m| = |a_n - A+A-a_m| \leq |a_n - A| + |a_m - A|.
	$$
	Pak podle (\ref{rovnicebolzano}) je
	$$
	|a_n - A| + |a_m - A| < \frac{\varepsilon}{2} + \frac{\varepsilon}{2} = \varepsilon.
	$$
	\item[$\impliedby$] Zvolme $\varepsilon \in \mathbb R , \varepsilon > 0$. Podle BC podmínky
	nalezneme $n \in \mathbb N $ takové, že
	$$
	\forall m,n \in \mathbb N , n\geq n_0, m\geq n_0: |a_n-a_m| < \varepsilon.
	$$
	Odtud plyne
	\begin{align*}
	&\forall n \in \mathbb N, n \geq n_0: |a_n - a_{n_0}| < \varepsilon, \\
	&	\forall n \in \mathbb N, n \geq n_0: a_{n_0} - \varepsilon < a_n < a_{n_0} +\varepsilon.
	\end{align*}
	Dále máme
	$$
	    a_{n_0} - \varepsilon \leq \liminf a_n \leq \limsup a_n \leq a_{n_0} + \varepsilon.
	$$
	Dostáváme tak, že $\limsup a_n, \liminf a_n \in \mathbb R $ a $0 \leq \limsup a_n - \liminf a_n \leq 2\varepsilon$.
	Platí tedy $\limsup a_n = \liminf a_n \in \mathbb R $, a tedy podle věty \ref{veta2.15}
	existuje vlastní limita $a_n$.\qedhere
\end{enumerate}
\end{proof}

\section{Limita a spojitost funkce}
\subsection{Základní pojmy}
\begin{definice}
\textbf{Funkce} $f$ \textbf{jedné reálné proměnné} (dále jen \textbf{funkce}) je zobrazení
$f: M\to \mathbb R $, kde $M$ je podmnožinou množiny reálných čísel.
\end{definice}

\begin{definice}
	Funkce $f: J\to \mathbb R $ je \textbf{rostoucí} na intervalu $J$, jestliže pro
	každou dvojici $x_1, x_2 \in J$, $x_1 < x_2$ platí nerovnost $f(x_1)<f(x_2)$. Analogicky
	definujeme funkci \textbf{klesající} (\textbf{neklesající}, \textbf{nerostoucí}) na intervalu $J$.
\end{definice}

\begin{definice}
\textbf{Monotónní funkcí} (resp. \textbf{ryze monotónní funkcí}) na intervalu $J$ rozumíme funkci,
která je neklesající nebo nerostoucí (resp. rostoucí nebo klesající) na $J$.
\end{definice}

\begin{definice}
Nechť $f$ je funkce a $M \subset \mathcal D(f)$. Řekneme, že funkce $f$ je
\begin{itemize}
  \item \textbf{lichá}, jestliže pro každé $x \in \mathcal D(f)$ platí $-x \in \mathcal D(f)$ a
  $f(-x) = -f(x)$,
  \item \textbf{sudá}, jestliže pro každé $x \in \mathcal D(f)$ platí $-x \in \mathcal D(f)$ a
  $f(-x) = f(x)$,
  \item \textbf{periodická s periodou} $a \in \mathbb{R}$, $a > 0$, jestliže pro každé
  $x \in \mathcal D(f)$ platí $x + a \in \mathcal D(f)$, $x - a \in \mathcal D(f)$ a
  $f(x + a) = f(x - a) = f(x)$,
  \item \textbf{shora omezená} na $M$, jestliže existuje číslo $K \in \mathbb{R}$ takové,
  že pro všechna $x \in M$ je $f(x) \le K$,
  \item \textbf{zdola omezená} na $M$, jestliže existuje číslo $K \in \mathbb{R}$ takové,
  že pro všechna $x \in M$ je $f(x) \ge K$,
  \item \textbf{omezená} na $M$, jestliže existuje číslo $K \in \mathbb{R}$ takové,
  že pro všechna $x \in M$ je $|f(x)| \le K$,
  \item \textbf{konstantní} na $M$, jestliže pro všechna $x, y \in M$ platí
  $f(x) = f(y)$.
\end{itemize}
\end{definice}

\subsection{Limita funkce}

\begin{definice}
	Nechť $c \in \mathbb R $. Potom \textbf{okolím bodu} $c$ rozumíme každou množinu tvaru
	$B(c, \varepsilon) = (c-\varepsilon, c+\varepsilon)$, kde $\varepsilon \in \mathbb R , \varepsilon > 0$.
	\textbf{Okolím} bodu $\infty$ rozumíme každou množinu tvaru $B(\infty, \varepsilon) = \left( \frac{1}{\varepsilon}, \infty \right ) $,
	kde $\varepsilon \in \mathbb R , \varepsilon > 0$. \textbf{Okolím} bodu $-\infty$
	rozumíme každou množinu $B(-\infty, \varepsilon)= \left( -\infty, -\frac{1}{\varepsilon} \right ) $, kde $\varepsilon \in \mathbb R $, $\varepsilon >0$.
\end{definice}

\begin{definice}
	Nechť $c \in \mathbb R $. Potom \textbf{prstencovým okolím} bodu $c$ rozumíme
	každou množinu tvaru $(c-\varepsilon, c+\varepsilon) \smallsetminus \left \{ c \right \} $,
	kde $\varepsilon \in \mathbb R ,\varepsilon > 0$, a kterou značíme
	$P(c,\varepsilon)$. \textbf{Prstencovým okolím} bodu $\infty$ rozumíme každou množinu tvaru $\left( \frac{1}{\varepsilon},\infty \right ) $,
	kde $\varepsilon \in \mathbb R , \varepsilon > 0$, a kterou značíme $P(\infty, \varepsilon)$.
	\textbf{Prstencovým okolím} bodu $-\infty$ rozumíme každou množinu tvaru
	$\left( -\infty, -\frac{1}{\varepsilon} \right ) $, kde $\varepsilon \in \mathbb R ,\varepsilon > 0$,
	a kterou značíme $P(-\infty, \varepsilon)$.
\end{definice}

\begin{pozorovani}
	Nechť $c_1,c_2 \in \mathbb R ^*, c_1 \ne c_2$. Potom existuje $\varepsilon \in \mathbb R , \varepsilon > 0$
	takové, že $B(c_1, \varepsilon) \cap B(c_2,\varepsilon) = \emptyset$.
\end{pozorovani}

\begin{definice}
	Řekneme, že číslo $A \in \mathbb R ^*$ je \textbf{limitou funkce $f$ v bodě $c \in \mathbb R ^*$},
	jestliže
	$$
	    \forall \varepsilon \in \mathbb R, \varepsilon > 0 \, \exists \delta \in \mathbb R , \delta > 0 \, \forall x \in P(c, \delta): f(x)\in B(A,\varepsilon).
	$$
\end{definice}

\begin{veta}\label{3.1}% 3.1
	Nechť funkce $f$ má v bodě $c\in \mathbb R ^*$ limitu $A\in \mathbb R ^*$ a limitu
	$B\in \mathbb R ^*$. Potom platí $A=B$.
\end{veta}

\begin{proof}
	Pro spor předpokládejme, že $A\ne B$. Nalezneme $\varepsilon \in \mathbb R , \varepsilon >0$
	takové, že $B(A,\varepsilon)\cap B(B,\varepsilon) = \emptyset$. Nalezneme $\delta_1, \delta_2 \in \mathbb R ,
	\delta_1, \delta_2 > 0$ takové, že platí
	$$
	    \forall x \in P(c,\delta_1): f(x) \in B(A,\varepsilon),\\
		\forall x \in P(c,\delta_2): f(x) \in B(B,\varepsilon).
	$$
	Položme $\delta_3 = \min \left \{ \delta_1, \delta_2 \right \} $. Zvolme $z\in P(c,\delta_3)$.
	Potom platí $c\in P(c,\delta_1)$, a tedy $f(z)\in B(A,\varepsilon).$ Podobně
	máme $z\in P(c,\delta_2)$, a tedy $f(z)\in B(B,\varepsilon).$ Odtud plyne
	$$
	    f(z)\in B(A, \varepsilon) \cap B(B,\varepsilon) = \emptyset,
	$$
	což je spor.
\end{proof}

\begin{poznamka}\,
\begin{enumerate}[(1)]
	\item Pokud existuje limita funkce $f$ v bodě $c\in \mathbb R ^*$, pak ji značíme symbolem
	$\lim_{x\to c} f(x)$.
	\item Pokud existuje $\lim_{x\to c} f(x)$, potom je $f$ definována na jistém prstencovém okolí
	bodu $c$.
	\item Limitu funkce můžeme počítat ve vlastním bodě $c$, tj. $c\in \mathbb R $, nebo v nevlastním
	bodě $c$, tj. $c\in \left \{ \infty, -\infty \right \} $. Výsledek může být opět vlastní nebo
	nevlastní.
	\item Nechť $c, A\in \mathbb R .$ Potom platí, že $\lim_{x\to c}f(x)=A$ právě tehdy, když
	$$
	    \forall \varepsilon \in \mathbb R , \varepsilon > 0 \, \exists \delta \in \mathbb R , \delta > 0 \, \forall x \in \mathbb R , 0 < |x-c|<\delta:|f(x)-A|<\varepsilon.
	$$
	\item Nechť $c,A\in \mathbb R , K\in \mathbb R , K>0$. Potom $\lim_{x\to c}f(x)=A$ právě tehdy, když platí
	$$
	\forall \varepsilon \in \mathbb R , \varepsilon > 0 \, \exists \delta \in \mathbb R , \delta > 0 \, \forall x \in \mathbb R , 0 < |x-c|<\delta:|f(x)-A|<K\varepsilon.
	$$
\end{enumerate}
\end{poznamka}

\begin{definice}
	Nechť $c\in \mathbb R $.
	\begin{itemize}
  \item \textbf{Pravým okolím} bodu $c$ rozumíme každý interval
  $[c, c + \varepsilon)$, kde $\varepsilon > 0$, a který značíme $B^+(c,\varepsilon)$,
  \item \textbf{levým okolím} bodu $c$ rozumíme každý interval
  $(c - \varepsilon, c]$, kde $\varepsilon > 0$, a který značíme $B^-(c,\varepsilon)$,
  \item \textbf{pravým prstencovým okolím} bodu $c$ rozumíme každý interval
  $(c, c + \varepsilon)$, kde $\varepsilon > 0$, a který značíme $P^+(c,\varepsilon)$,
  \item \textbf{levým prstencovým okolím} bodu $c$ rozumíme každý interval
  $(c - \varepsilon, c)$, kde $\varepsilon > 0$, a který značíme $P^-(c,\varepsilon)$.
\end{itemize}
Pokračujme s definicí pro nevlastní hodnoty.
\begin{itemize}
  \item \textbf{Pravým okolím bodu $-\infty$} rozumíme každý interval
  $\left(-\infty, -\frac{1}{\varepsilon}\right)$, kde $\varepsilon > 0$,
  a který značíme $B^+(-\infty,\varepsilon)$,
  \item \textbf{levým okolím bodu $\infty$} rozumíme každý interval
  $\left(\frac{1}{\varepsilon}, \infty\right)$, kde $\varepsilon > 0$,
  a který značíme $B^-(\infty,\varepsilon)$,
  \item \textbf{pravým prstencovým okolím bodu $-\infty$} rozumíme každý interval
  $\left(-\infty, -\frac{1}{\varepsilon}\right)$, kde $\varepsilon > 0$,
  a který značíme $P^+(-\infty,\varepsilon)$,
  \item \textbf{levým prstencovým okolím bodu $\infty$} rozumíme každý interval
  $\left(\frac{1}{\varepsilon}, \infty\right)$, kde $\varepsilon > 0$,
  a který značíme $P^-(\infty,\varepsilon)$.
\end{itemize}
\end{definice}

\begin{definice}
Nechť $A \in \mathbb{R}^\star$, $c \in \mathbb{R} \cup \{-\infty\}$.
Řekneme, že funkce $f$ má v bodě $c$ \textbf{limitu zprava} rovnou $A$, jestliže
\[
\forall \varepsilon \in \mathbb{R},\ \varepsilon > 0\ \exists \delta \in \mathbb{R},\
\delta > 0\ \forall x \in P^+(c,\delta) : f(x) \in B(A,\varepsilon).
\]
Analogicky definujeme pojem \textbf{limity zleva} v bodě
$c \in \mathbb{R} \cup \{\infty\}$.
\end{definice}

\begin{poznamka}
	Věta \ref{3.1} platí analogicky i pro limitu zleva a limitu zprava. Značení: $\lim_{x\to c_-}f(x)$, resp.
	$\lim_{x\to c_+}f(x)$.
\end{poznamka}

\begin{priklad}\,
\begin{itemize}
	\item Nechť $c\in \mathbb R ^*, A \in \mathbb R ,\delta_0\in \mathbb R ,\delta_0>0$ splňují
	$\forall x \in P(c,\delta_0):f(x)=A.$
	Potom $\lim_{x\to c}f(x)=A.$
	\item Nechť $c\in \mathbb R ^*$. Potom platí $\lim_{x\to c}x=c$. Zvolme
	$\varepsilon \in \mathbb R ,\varepsilon >0$. Položme $\delta=\varepsilon$. Potom platí
	$P(c,\delta)=P(c,\varepsilon)\subseteq B(c,\varepsilon),$ a tedy
	$\forall x \in P(c,\delta):x\in B(c,\varepsilon)$.
\end{itemize}
\end{priklad}

\begin{pozorovani}
	Nechť $f,g$ jsou funkce, $c\in \mathbb R ^*, \delta_0 \in \mathbb R , \delta_0 > 0$ a platí
	$$
	    \forall x\in P(c,\delta_0): f(x)=g(x).
	$$
	Potom $\lim_{x\to c}f(x)=A\in \mathbb R ^*$ právě tehdy, když $\lim_{x\to c}g(x)=A\in \mathbb R ^*$.
\end{pozorovani}

\begin{poznamka}
	Limita funkce v bodě má \uv{lokální charakter}.
\end{poznamka}

\begin{definice}
	Řekneme, že funkce $f$ je \textbf{spojitá v bodě} $c\in \mathbb R $, jestliže platí
	$\lim_{x\to c} f(x)=f(c)$.
\end{definice}

\begin{poznamka}
	Funkce je spojitá v bodě $c\in \mathbb R $ právě tehdy, když platí
	$$
	    \forall \varepsilon \in \mathbb R , \varepsilon > 0 \,\exists \delta \in \mathbb R , \delta > 0 \, \forall x \in B(c,\delta):f(x)\in B(f(c),\varepsilon).
	$$
\end{poznamka}

\begin{definice}
Řekneme, že funkce $f$ je \textbf{spojitá zprava} (resp. \textbf{zleva}) v bodě $c \in \mathbb{R}$,
jestliže platí
\[
\lim_{x \to c^+} f(x) = f(c)
\quad
(\text{resp. } \lim_{x \to c^-} f(x) = f(c)).
\]
\end{definice}

\begin{definice}
Nechť $J \subset \mathbb{R}$ je interval. Funkce $f \colon J \to \mathbb{R}$
je \textbf{spojitá na intervalu} $J$, jestliže platí:
\begin{itemize}
  \item $f$ je spojitá zprava v každém bodě $J$, který není pravým koncovým bodem $J$,
  \item $f$ je spojitá zleva v každém bodě $J$, který není levým koncovým bodem $J$.
\end{itemize}
\end{definice}

\begin{poznamka}
	Nechť $c\in \mathbb R , A\in \mathbb R ^*$. Potom platí, že $\lim_{x\to c}f(x)=A$ právě tehdy,
	když $\lim_{x\to c_-}f(x)=A$ a $\lim_{x\to c_+}f(x)=A$.
	\par Implikace $\implies$ je zřejmá. Implikace $\impliedby$: Zvolme $\varepsilon \in \mathbb R ,\varepsilon > 0$.
	Nalezneme $\delta _1\in \mathbb R , \delta_1 >0$ takové, že
	$$
	    \forall x\in P_1(c,\delta_1) : f(x)\in B(A,\varepsilon).
	$$
	Podobně nalezneme $\delta_2\in \mathbb R , \delta_2 >0$ takové, že
	$$
	    \forall x\in P_2(c,\delta_2):f(x)\in B(A,\varepsilon).
	$$
	Položme $\delta_3=\min \left \{ \delta_1, \delta_2 \right \} $. Potom máme
	$$
	    P(c,\delta_3) \subseteq P_+(c,\delta_1) \cup P_-(c,\delta_2).
	$$
	Pro každé $x \in P(c,\delta_3)$ platí $f(x)\in B(A,\varepsilon)$.
\end{poznamka}

\subsection{Věty o limitách}
\begin{veta}\label{veta3.2}
	Nechť funkce $f$ má v bodě $c\in \mathbb R ^*$ vlastní limitu. Pak existuje
	$\delta\in \mathbb R ,\delta >0$ takové, že $f$ je na $P(c,\delta)$ omezená.
\end{veta}

\begin{proof}
	Označme $A=\lim_{x\to c}f(x)$. Potom platí $A\in \mathbb R $. Položme $\varepsilon =1$. K němu
	nalezneme $\delta\in \mathbb R ,\delta > 0$ takové, že pro každé $x\in P(c,\delta): f(x)\in B(A,\varepsilon)$,
	neboli
	$$
	    \forall x \in P(c,\delta):f(x)\in(A-1, A+1).
	$$
	Tedy $f$ je omezená na $P(c,\delta)$.
\end{proof}

\begin{veta}[Aritmetika limit]\label{veta3.3} % věta 3.3
Nechť $A, B, c \in \mathbb{R}^\star$, $\lim_{x \to c} f(x) = A$ a
$\lim_{x \to c} g(x) = B$. Potom platí:
\begin{itemize}
  \item[(a)]
  $
  \lim_{x \to c} \bigl(f(x) + g(x)\bigr) = A + B,
  $
  pokud je výraz na pravé straně definován,
  \item[(b)]
  $
  \lim_{x \to c} f(x) g(x) = AB,
  $
  pokud je výraz na pravé straně definován,
  \item[(c)]
  $
  \lim_{x \to c} \frac{f(x)}{g(x)} = \frac{A}{B},
  $
  pokud je výraz na pravé straně definován.
\end{itemize}
\end{veta}

\begin{proof}\,
\begin{enumerate}
	\item[(b)] Případ $A,B \in \mathbb R $: Podle věty \ref{veta3.2}
	nalezneme $K\in \mathbb R , K > 0$ takové, že existuje $\delta_0 \in \mathbb R , \delta_0 > 0$ splňující
	$$
	    \forall x \in P(c,\delta_0):|f(x)| < K.
	$$
	Zvolme $\varepsilon \in \mathbb R, \varepsilon > 0$. Nalezneme $\delta_1,\delta_2 \in \mathbb R ,\delta_1, \delta_2 > 0$
	taková, že platí
	\begin{align*}
		&\forall x \in P(c,\delta_1): |f(x)-A|<\varepsilon, \\
		&\forall x \in P(c,\delta_2): |g(x)-B|< \varepsilon.
	\end{align*}
	Položme $\delta = \min \left \{ \delta_0, \delta_1, \delta_2 \right \} $.Pro každé $x\in P(c,\delta)$ platí
	\begin{align*}
	|f(x)g(x)-AB| & = |f(x)g(x)-f(x)B+f(x)B-AB| \\
	& \leq |f(x)(g(x)-B)| + |(f(x)-A)B|++
	& < K\varepsilon + \varepsilon |B| = (K+|B|)\varepsilon.
\end{align*}
Odtud plyne $\lim_{x\to c} f(x)g(x)=AB.$
\item[(c)] Případ $A = \infty, B\in \mathbb R , B> 0$. Potom $\frac{A}{B}\infty$. Zvolme $\varepsilon \in \mathbb R ,\varepsilon > 0$.
K němu nalezneme $\delta_1, \delta_2 \in \mathbb R , \delta_1, \delta_2 > 0$ taková, že platí
\begin{align*}
\forall x \in P(c,\delta_1): f(x)\in B(\infty, \frac{1}{2B}\varepsilon),\\
\forall x \in P(c,\delta_2): g(x) \in B(B,\frac{1}{2}B).
\end{align*}
    Položme $\delta = \min \left \{ \delta_1, \delta_2 \right \} $. Pro každé $x \in P(c,\delta)$ platí
    $$
    \frac{f(x)}{g(x)}> \frac{\frac{2B}{\varepsilon}}{\frac{3}{2}B} = \frac{4}{3\varepsilon} > \frac{1}{\varepsilon},
    $$
    tj. $\frac{f(x)}{g(x)} \in B(\infty, \varepsilon).$ \qedhere
\end{enumerate}
\end{proof}

\begin{poznamka}
	Věta \ref{veta3.3} platí analogicky i pro jednostranné limity.
\end{poznamka}

\begin{dusledek}\label{duslvety3.3}
	Nechť $c\in \mathbb R $ a $f,g$ jsou spojité funkce v bodě $c$. Potom také $f+g$ a $fg$ jsou spojité
	v bodě $c$. Pokud navíc $g(c)\ne 0$, pak i funkce $\frac{f}{g}$ je spojitá v $c$.
\end{dusledek}

\begin{priklad}
	Funkce $x\mapsto c, x \in \mathbb R $ a $x\mapsto x, x \in \mathbb R $ jsou spojité na $\mathbb R $.
	Odtud a z důsledku \ref{duslvety3.3} plyne, že každý polynom je spojitý v každém bodě z $\mathbb R $.
\end{priklad}

\begin{veta}\label{veta3.4} % věta 3.4
Nechť $A \in \mathbb{R}^\star$, $A > 0$, $c \in \mathbb{R}^\star$,
$\lim_{x \to c} g(x) = 0$, $\lim_{x \to c} f(x) = A$ a existuje
$\eta \in \mathbb{R}$, $\eta > 0$, takové, že pro každé
$x \in P(c,\eta)$ platí $g(x) > 0$.
Potom platí
\[
\lim_{x \to c} \frac{f(x)}{g(x)} = \infty.
\]
\end{veta}

\begin{proof}
	Rozlišíme dva případy
	\begin{enumerate}[$i$.]
		\item $A\in \mathbb R $: Zvolme $\varepsilon \in \mathbb R , \varepsilon > 0$. K němu
		nalezneme $\delta_1, \delta_2 \in \mathbb R , \delta_1, \delta_2 > 0$ taková, že
		\begin{align*}
		&\forall x \in P(c, \delta_1): |f(x)-A|<\frac{1}{2}A,\\
		&\forall x \in P(c,\delta_2): |g(x) - 0| = |g(x)| < \frac{1}{2} A\varepsilon.
		\end{align*}
		Položme $\delta = \min \left \{ \delta_1, \delta_2, \eta \right \}. $ Potom pro každé $x \in P(c,\delta)$
		platí
		$$
  \frac{f(x)}{g(x)} > \frac{\frac{1}{2}A}{\frac{1}{2}A\varepsilon}=\frac{1}{\varepsilon},
		$$
		tj. $\frac{f(x)}{g(x)} \in B(\infty, \varepsilon)$. Potom tedy $\lim_{x\to c} \frac{f(x)}{g(x)} = \infty$.
		\item $A = \infty$: Zvolme $\varepsilon \in \mathbb R , \varepsilon > 0$. K němu
		nalezneme $\delta_1, \delta_2\in \mathbb R , \delta_1, \delta_2 >0$ taková, že
		\begin{align*}
			&\forall x \in P(c,\delta_1): f(x) \in B(\infty, \varepsilon), \\
			&\forall x \in P(c,\delta_2): G(x) \in B(0, 1).
        \end{align*}
        Položme $\delta = \min \left \{ \delta_1, \delta_2,\eta \right \} $. Potom pro každé $x \in P(c,\delta)$
        platí
        $$
        \frac{f(x)}{g(x)}> \frac{\frac{1}{\varepsilon}}{1} = \frac{1}{\varepsilon}.
        $$
        Odtud plyne $\lim_{x\to c} \frac{f(x)}{g(x)}=\infty$. \qedhere
	\end{enumerate}
\end{proof}

\begin{veta}[Limita a uspořádání]\label{veta3.5} % 3.5
Nechť $c \in \mathbb{R}^\star$ a funkce $f$ a $g$ mají limity v $c$.
\begin{itemize}
  \item[(a)] Nechť
  \[
  \lim_{x \to c} f(x) > \lim_{x \to c} g(x).
  \]
  Pak existuje $\delta \in \mathbb{R}$, $\delta > 0$, takové, že platí
  \[
  \forall x \in P(c,\delta) : f(x) > g(x).
  \]
  \item[(b)] Nechť existuje $\delta \in \mathbb{R}$, $\delta > 0$, takové, že platí
  \[
  \forall x \in P(c,\delta) : f(x) \le g(x).
  \]
  Potom platí
  \[
  \lim_{x \to c} f(x) \le \lim_{x \to c} g(x).
  \]
\end{itemize}
\end{veta}

\begin{proof}
	Označme $A = \lim_{x\to c} f(x), B = \lim_{x\to c} g(x)$.
	\begin{enumerate}[(a)]
		\item Nalezneme $\varepsilon \in \mathbb R , \varepsilon > 0$
		takové, že $B(A,\varepsilon) \cap B(B,\varepsilon) = \emptyset.$ Potom pro každé $y_1 \in B(A,\varepsilon)$
		a $y_2\in B(B,\varepsilon)$ platí $y_1 > y_2$. Nalezneme $\delta_1, \delta_2 \in \mathbb R , \delta_1, \delta_2 > 0$
		taková, že
		\begin{align*}
			& \forall x \in P(c,\delta_1):f(x)\in B(A,\varepsilon), \\
			& \forall x \in P(c,\delta_2):g(x)\in B(B,\varepsilon).
\end{align*}
    Položme $\delta = \min \left \{ \delta_1, \delta_2 \right \} $. Potom pro každé $x \in P(c,\delta)$
    platí $f(x)\in B(A,\varepsilon),g(x)\in B(B,\varepsilon)$, a tedy $f(x) > g(x)$.
    \item Pro spor předpokládejme, že $A > B$. Podle již dokázaného bodu $(a)$ nalezneme
    $\eta \in \mathbb R , \eta > 0$ takové, že
    $$
        \forall x \in P(c,\eta): f(x) > g(x).
    $$
    Zvolme $z \in P(c, \min \left \{ \delta, \eta \right \} )$. Potom platí
    $$
        f(z) > g(z) \geq f(z),
    $$
    což je spor. (Druhá nerovnost z předpokladu.) \qedhere
\end{enumerate}
\end{proof}

\begin{veta}[O dvou strážnících] % 3.6
Nechť $A \in \mathbb{R}^\star$, $\eta \in \mathbb{R}$, $\eta > 0$, a funkce
$f$, $g$ a $h$ splňují:
\begin{itemize}
  \item $\forall x \in P(c,\eta) : f(x) \le h(x) \le g(x)$,
  \item $\lim_{x \to c} f(x) = \lim_{x \to c} g(x) = A$.
\end{itemize}
Potom existuje $\lim_{x \to c} h(x)$ a rovná se $A$.
\end{veta}

\begin{proof}
	Zvolme $\varepsilon \in \mathbb R , \varepsilon > 0$. K němu nalezneme $\delta_1, \delta_2 \in \mathbb R , \delta_1, \delta_2 > 0$
	taková, že
	\begin{align*}
		& \forall x \in P(c,\delta_1):f(x)\in B(A, \varepsilon), \\
		& \forall x \in P(c,\delta_2):g(x)\in B(A, \varepsilon).
\end{align*}
    Položme $\delta = \min \left \{ \delta_1, \delta_2,\eta \right \} $. Potom pro každé
    $x \in B(c,\delta)$ platí $f(x), g(x)\in B(A, \varepsilon)$ a navíc $f(x)\leq h(x)\leq g(x)$,
    tedy $h(x) \in B(A, \varepsilon)$. Odtud plyne, že $\lim_{x\to c} h(x) = A.$
\end{proof}

\begin{poznamka}
	Pokud $A = \infty$, stačí strážník $f$. Podobně pro $A = \infty$.
\end{poznamka}

\begin{dusledek}
	Nechť $c \in \mathbb R ^*, \lim _{x\to c} f(x) = 0$ a existuje $\eta \in \mathbb R ,\eta > 0$
	takové, že funkce $g$ je na $P(c,\eta)$ omezená. Potom $\lim_{x\to c} f(x)g(x)=0.$
\end{dusledek}

\begin{proof}
	Nalezneme $K\in \mathbb R , K > 0$ takové, že $\forall x \in P(c,\eta): |g(x)| <K$.
	Platí také $\lim_{x\to c} |f(x)|=0.$ Potom $\forall x \in P(c, \eta)$ platí
	$$
	    -K |f(x)| \leq f(x)g(x) \leq K|f(x)|.
	$$
	Odtud podle věty o dvou strážnících máme $\lim_{x\to c} f(x)g(x) = 0.$
\end{proof}

\begin{veta}[Limita složené funkce] \label{veta3.7}
Nechť $x_0,y_0\in\R^*$, $f,g$ jsou funkce, $\lim_{x\to x_0} g(x)=y_0$, $\lim_{y\to y_0} f(y)$ existuje a je splněna alespoň jedna z podmínek:
\begin{enumerate}
	\item[(P)] $\exists\eta \in \mathbb R , \eta >0\ \forall x\in P(x_0,\eta): g(x)\ne y_0$,
	\item[(S)] $f$ je spojitá v bodě $y_0$.
\end{enumerate}
Potom platí $$\lim_{x\to x_0} (f\circ g)(x)=\lim_{y\to y_0} f(y).$$
\end{veta}

\begin{proof}
	Označme $A=\lim_{y\to \_0}f(y)$. Chceme $\lim_{x\to x_0} f(g(x))=A$.
	\begin{enumerate}
		\item[(P)] Zvolme $\varepsilon \in \mathbb R , \varepsilon > 0$. K nim nalezneme
		$\psi \in \mathbb R , \psi > 0$ takové, že
		$$\forall y \in P(y_0, \psi): f(y)\in B(A, \varepsilon).$$
		K $\psi$ nalezneme $\delta^\prime \in \mathbb R , \delta ^\prime > 0$ takové, že platí
		$$
		    \forall x \in P(x_0,\delta^\prime): g(x)\in B(y_0,\psi).
		$$
		Nechť $\eta \in \mathbb R,\eta>0 $ splňuje podmínku (P). Položme
		$\delta = \min \left \{ \delta^\prime, \eta \right \} $.
		Pro každé $x \in P(x_0,\delta)$ platí $g(x\in B(y_0,\psi)$ a $g(x)\ne y_0$, tedy $g(x)\in P(y_0, \psi)$.
		Odtud plyne $f(g(x))\in B(A,\varepsilon).$
		\item[(S)] Zvolme $\varepsilon \in \mathbb R , \varepsilon > 0$. K němu díky (S) nalezneme
		$\psi \in \mathbb R , \psi > 0 $ takové, že
		$$
		    \forall y \in B(y_0,\psi) : f(y) \in B(f(y_0), \varepsilon),
		$$
		neboť $A=f(y_0)$. K $\psi$ nalezneme $\delta \in \mathbb R ,\delta > 0$ takové, že
		$$
		    \forall x \in P(x_0, \delta):g(x) \in B(y_0, \psi).
		$$
		Potom pro každé $x \in P(x_0, \delta)$ platí $f(g(x)) \in B(f(y_0),\varepsilon)=B(A,\varepsilon)$. \qedhere
	\end{enumerate}
\end{proof}

\begin{dusledek}
	Jestliže je funkce $g$ spojitá v bodě $c\in \mathbb R $ a funkce $f$ je spojitá v bodě $g(c)$, potom
	je funkce $f\circ g$ spojitá v bodě $c$.
\end{dusledek}

\begin{poznamka}
	Platí i varianty věty \ref{veta3.7} pro jednostranné limity, např.: Nechť $x_0,y_0 \in \mathbb R \smallsetminus \left \{ -\infty \right \} $,
	$f,g$ jsou funkce, $\lim_{x\to x_0^+} g(x)=y_0,\lim_{y\to y_0^+} f(y)$ existuje a je splněna alespoň jedna z následujících
	podmínke:
	\begin{enumerate}
		\item[(P)] $\exists \eta \in \mathbb R , \eta > 0\, \forall x \in P_+(x_0,\eta):g(x)>y_0$,
		\item[(S)] $f$ je spojitá v bodě $y_0$,
\end{enumerate}
    potom $\lim_{x\to x_0^+} f(g(x))=\lim_{y\to y_0^+}f(y)$.
\end{poznamka}

\begin{veta}[Heine]\label{Heinehoveta}
Nechť $c,A\in \mathbb R ^*$ a $f$ je funkce. Pak jsou následující výroky ekvivalentní:
\begin{enumerate}[(i)]
  \item $\lim_{x\to c} f(x)=A$.
  \item Pro každou posloupnost $\{x_n\}$ splňující
  \begin{itemize}
	\item $\forall n\in \mathbb N :x_n\ne c$ a
	\item $\lim_{n\to\infty} x_n=c$
  \end{itemize}
  platí $\lim_{n\to\infty} f(x_n)=A$.
\end{enumerate}
\end{veta}

\begin{proof}\,
\begin{enumerate}
	\item[$\implies$] Nechť $\left \{ x_n \right \} $ je posloupnost splňující
	\begin{itemize}
		\item $\forall n\in \mathbb N : x_n \ne c,$
		\item $\lim_{n \to \infty} x_n = c$.
	\end{itemize}
	Chceme ukázat, že $\lim_{n \to \infty} f(x_n)=A$. Zvolme $\varepsilon \in \mathbb R , \varepsilon > 0$.
	Díky ($i$) nalezneme k $\varepsilon$ číslo $\delta\in \mathbb R , \delta > 0$ takové, že
	$$
	    \forall x \in P(c,\delta):  f(x) \in B(A,\varepsilon).
	$$
	K $\delta$ nalezneme $n_0 \in \mathbb N $ takové, že
	$$
	    \forall n \in \mathbb N , n \geq n_0: x_n \in B(c,\delta),
	$$
	neboť $\lim_{n \to \infty} x_n = c$. Pak máme
	$$
	    \forall n \in \mathbb N , n \geq n_0: x_n \in P(c,\delta).
	$$
	Potom z předchozího plyne
	$$
	    \forall n \in \mathbb N , n \geq n_0: f(x_n) \in B(A,\varepsilon).
	$$
	Odtud $\lim_{n \to \infty} f(x_n)=A.$
	\item[$\impliedby$] Obměnou. Předpokládejme negaci ($i$):
	$$
	    \exists \varepsilon \in \mathbb R , \varepsilon > 0 \, \forall \delta \in \mathbb R , \delta > 0 \, \exists x \in P(c,\delta): \neg(f(x)\in B(A,\varepsilon)).
	$$
	Nechť $\varepsilon \in \mathbb R , \varepsilon > 0$ a podle $\neg$ ($i$) splňuje
	$$
	    \forall \delta \in \mathbb R , \delta > 0 \, \exists x \in P(c,\delta): \neg (f(x)\in B(A,\varepsilon).
	$$
	Pro každé $n \in \mathbb N $ nalezneme $x_n \in P \left( c, \frac{1}{n} \right ) $ takové, že
	$\neg(f(x_n)\in B(A,\varepsilon)$. Potom pro každé $n \in \mathbb N $ platí $x_n \ne c$.
	Dále máme $\lim_{n \to \infty} x_n = c$. Zvolme $\tau \in \mathbb R , \tau > 0$. K němu nalezneme
	$n_0 \in \mathbb N $ splňující $\frac{1}{n_0} < \tau$. Potom pro každé $n\in \mathbb N , n \geq n_0$
	platí
	$$
	    x_n \in P(c, \frac{1}{n}) \subseteq P(c,\frac{1}{n_0})\subseteq B(c,\tau).
	$$
	Neplatí však $\lim_{n \to \infty} f(x_n)=A$, neboť pro každé $n \in \mathbb N $ máme
	$\neq (f(x_n)\in B(A,\varepsilon))$.\qedhere
\end{enumerate}
\end{proof}

\begin{poznamka}
	Lze formulovat varianty Heineovy věty, např.: Nechť $c\in \mathbb R $ a $f$ je funkce.
	Pak jsou výroky ($i$) a ($ii$) ekvivalentní:
	\begin{enumerate}[($i$)]
		\item $f$ je spojitá zprava v bodě $c$,
		\item pro každou posloupnost $\left \{ x_n \right \} $ splňující
		\begin{itemize}
			\item $\forall n \in \mathbb N : x_n \geq c$,
			\item $\lim_{n \to \infty} x_n = c$,
		\end{itemize}
		platí $\lim_{n \to \infty} f(x_n)=f(c)$.
	\end{enumerate}
\end{poznamka}

\begin{priklad}
	Uvažujme funkci $f(x)=\sin \frac{1}{x}, x \in \mathbb R \smallsetminus \left \{ 0 \right \} $.
	Položme $x_n=\frac{1}{n\pi},n\in \mathbb N $, $y_n = \frac{1}{2n\pi + \frac{\pi}{2}}, n \in \mathbb N $.
	Platí $x_n \ne 0, y_n \ne 0$ pro všechna $n \in \mathbb N $. Navíc je
	\begin{align*}
		\lim_{n \to \infty} f(x_n)=\lim_{n \to \infty} \sin(n\pi)=0, \\
		\lim_{n \to \infty} f(y_n) = \lim_{n \to \infty} \left( 2n\pi + \frac{\pi}{2} \right ) =1,
\end{align*}
tedy $\lim_{x \to 0} f(x)$ neexistuje.
\end{priklad}

\begin{veta}[Limita monotónní funkce]\label{veta3.9}
Nechť $a,b\in \mathbb R ^*$, $a<b$, $f$ je monotónní funkce na intervalu $(a,b)$. Potom existují $\lim_{x\to a^+} f(x)$ a $\lim_{x\to b^-} f(x)$, přičemž platí:
\begin{enumerate}[(a)]
  \item je-li $f$ na $(a,b)$ neklesající, pak
  \begin{align*}
	\lim_{x\to a^+} f(x)=\inf f((a,b)) && \textrm{a} & & \lim_{x\to b^-} f(x)=\sup f((a,b)),
  \end{align*}
  \item je-li $f$ na $(a,b)$ nerostoucí, pak
  \begin{align*}
	\lim_{x\to a^+} f(x)=\sup f((a,b)) && \textrm{a} & & \displaystyle \lim_{x\to b^-} f(x)=\inf f((a,b)).
  \end{align*}
\end{enumerate}
\end{veta}

\begin{proof}\,
\begin{enumerate}[(a)]
	\item Označme $m=\inf f((a,b))$. Předpokládejme, že $m\in \mathbb R $. Zvolme $\varepsilon \in \mathbb R , \varepsilon > 0$.
	Podle definice infima není číslo $m+\varepsilon$ dolní závorou množiny $f((a,b))$. Nalezneme $x^\prime \in (a,b)$
	takové, že $m\leq f(x^\prime)< m + \varepsilon$. Nalezneme $\delta \in \mathbb R , \delta  > 0$ takové, že
	$P_+(a,\delta)\subset (a,x^\prime).$ Potom pro každé $x\in P_+(a,\delta)$ platí, že $a<x<x^\prime$, a tedy
	$$
	    m-\varepsilon < m \leq f(x) \leq f(x^\prime) < m+ \varepsilon.
	$$
	Potom máme $f(x)\in B(m,\varepsilon)$. Podobně postupujeme, pokud $m=-\infty$.
	\item Obdobně.\qedhere
\end{enumerate}
\end{proof}

\subsection{Funkce spojité na intervalu}
\begin{veta}[Bolzano]\label{veta3.10}
Nechť funkce $f$ je spojitá na intervalu $[a,b]$ a předpokládejme, že $f(a)<f(b)$. Potom pro každé $C\in (f(a),f(b))$ existuje $\xi\in(a,b)$ takové, že $f(\xi)=C$.
\end{veta}

\begin{proof}
	Nechť $C\in (f(a),f(b)).$ Položme $M = \left \{ z\in \left [ a,b \right ] , f(z)< C \right \} $.
	Množina $M$ je neprázdná ($a \in M$) a shora omezená, neboť $M \subseteq \left [ a,b \right ] $.
	Označme $\xi=\sup M$ (používáme vlastnost suprema). Přivedeme ke sporu možnost, že $f(\xi) < C$ i
	$f(\xi) < C$.
	\begin{enumerate}[($i$)]
		\item Pro spor předpokládejme, že $f(\xi)<C$. Platí tedy, že $\xi < b$. Funkce $f$ je v
		$\xi$ spojitá zprava. Podle věty \ref{veta3.5} (a) nalezneme $\delta \in \mathbb R ,\delta > 0$
		takové, že
		$$
		    \forall x \in \left [ \xi,\xi<\delta \right ):f(x)< C.
		$$
		Potom máme, že $\left [ \xi, \xi+ \delta \right ) \subset M $, což je spor.
		\item Pro spor předpokládejme, že $f(\xi) < C$. Potom $\xi > a$. Funkce $f$ je zleva spojitá
		v $\xi$. Podle věty \ref{veta3.5} (a) nalezneme $\delta \in \mathbb R , \delta > 0$
		takové, že
		$$\forall x \in (\xi-\delta, \xi]:f(x)>C.$$
		Potom $M\subseteq \left [ a,\xi-\delta \right ] $, což je spor ($\xi-\delta <\xi$  a $ \xi$ je
		horní závora).\qedhere
	\end{enumerate}
\end{proof}

\begin{poznamka}
	Věta \ref{veta3.10} platí i v případě, že $f(a) >f(b)$.
\end{poznamka}

\begin{lemma}\label{lemma3.11} % 3.11
Nechť $M\subset \mathbb R $ a platí
$$\forall x,y\in M\ \forall z\in \mathbb R :\ x<z<y\implies z\in M.$$
Pak $M$ je interval.
\end{lemma}

\begin{proof}
	Pokud $M = \emptyset$, pak $M$ je interval. Předpokládejme tedy, že $M\ne \emptyset$.
	Položme $a=\inf M, b =\sup M$. Ověříme, že interval $(a,b) \subseteq M \subseteq (a,b)\cup \left \{ a,b \right \} $.
	Nejprve ověříme $(a,b) \subseteq M$. Zvolme $z\in (a,b)$. Nalezneme prvky $x,y\in M$ splňující
	$x<z$ a $z<y$. Podle podmínky ze znění máme $z\in M$. Nyní ověříme $M \subseteq (a,b)\cup \left \{ a,b \right \} $.
	Zvolme $z\in M$. Potom $a\leq z\leq b$. Odtud plyne $z\in(a,b) \cup \left \{ a,b \right \} $.
\end{proof}

\begin{veta}\label{veta3.12} % 3.12
Nechť funkce $f\colon J\to\R$ je spojitá na intervalu $J$. Potom je $f(J)$ interval.
\end{veta}

\begin{proof}
	Použijeme lemma \ref{lemma3.11}. Předpokládejme, že $y_1,y_2 \in f(J)$ a $z\in \mathbb R $
	splňují $y_1 < z < y_2$. Nalezneme $x_1,x_2 \in J$ takové, že $f(x_1) = y_1, f(x_2) = y_2$.
	Podle věty \ref{veta3.10} nalezneme $\xi \in (x_1, x_2)$ (resp. $\xi \in (x_2, x_1)$) splňující
	$f(\xi) = z$ ($f$ je spojitá na $J$, a tedy i na $\left [ x_1,x_2 \right ] $, resp. $\left [ x_2,x_1 \right ] $).
	Máme tedy $z\in f(J)$. Podle lemmatu \ref{lemma3.11} je $f(J)$ interval.
\end{proof}

\begin{veta}\label{veta3.13}% 3.13
Nechť $f$ je spojitá funkce na intervalu $[a,b]$. Potom je $f$ na $[a,b]$ omezená.
\end{veta}

\begin{proof}
	Pro spor předpokládejme, že $f$ není omezená shora. Pro každé $n\in \mathbb N $
	nalezneme $x_n \in \left [ a,b \right ] $ takový, že $f(x_n) > n$. Díky Weierstrassově větě
	nalezneme vybranou posloupnost $\left \{ x_{n_k} \right \}_{k=1}^\infty $ z $\left \{ x_n \right \}_{n=1}^\infty $.
	Pak $\lim_{k \to \infty} x_{n_k} = x^* \in \left [ a,b \right ] $. Z varianty Heineovy věty plyne
	$\lim_{k \to \infty} f(x_{n_k})=f(x^*)\in \mathbb R $, $\forall k \in \mathbb N : f(x_{n_k}) > n_k \geq k$.
	odstud plyne $\lim_{k \to \infty} f(x_{n_k}) = \infty$, což je spor.
\end{proof}

\begin{poznamka}
	V důkazu věty \ref{veta3.13} jsme užili následující variantu Heineovy věty:
	Nechť $I\subset \mathbb R $ je interval a $f:I\to \mathbb R $. Pak je ekvivalentní:
	\begin{enumerate}[(i)]
		\item $f$ je spojitá na $I$,
		\item pro každé $c\in I$ a každou posloupnost $\left \{ x_n \right \} $ prvků intervalu $I$
		splňující $\lim_{n \to \infty} x_n=c$ platí $\lim_{n \to \infty} f(x_n)=f(c)$.
\end{enumerate}
\end{poznamka}

\begin{definice}
Nechť $M\subset\R$, $x\in M$ a funkce $f$ je definována alespoň na $M$ (tj. $M\subset\mathcal D(f)$).
Řekneme, že $f$ nabývá v bodě $x$ \textbf{maxima} (resp. \textbf{minima}) na $M$, jestliže platí
\begin{align*}
	\forall y\in M:\ f(y)\le f(x) && \textrm{(resp. }\forall y\in M:\ f(y)\ge f(x)\textrm{)}.
\end{align*}
Bod $x$ nazýváme \textbf{bodem maxima}
(resp. \textbf{minima}) funkce $f$ na množině $M$. Symbol $\max_M f$ (resp. $\min_M f$) označuje největší (resp. nejmenší) hodnotu,
které funkce $f$ na množině $M$ nabývá (pokud taková hodnota existuje).
\end{definice}

\begin{definice}
Nechť $M\subset \mathbb R $, $x\in M$ a funkce $f$ je definována alespoň na $M$ (tj. $M\subset \mathcal D(f)$).
Řekneme, že funkce $f$ má v bodě $x$
\begin{itemize}
	\item \textbf{lokální maximum vzhledem k $M$}, jestliže existuje $\delta>0$ takové,
	že pro každé $y\in P(x,\delta)\cap M$ platí $f(y)\le f(x)$,
	\item \textbf{lokální minimum vzhledem k $M$}, jestliže existuje $\delta>0$ takové,
	že pro každé $y\in P(x,\delta)\cap M$ platí $f(y)\geq f(x)$.
\end{itemize}
\end{definice}

\begin{veta}\label{veta3.14} % 3.14
Nechť $a,b\in \mathbb R $, $a<b$, a $f$ je spojitá funkce na intervalu $[a,b]$. Potom $f$ nabývá na $[a,b]$ svého maxima a svého minima.
\end{veta}

\begin{proof}
	Dokážeme existenci maxima. Podle věty \ref{veta3.13} je $f$ omezená na intervalu $\left [ a,b \right ] $,
	a proto $G=\sup f(\left [ a,b \right ] )$ je prvkem $\mathbb R $. Pro každé $n\in \mathbb N $
	nalezneme $x_n \in \left [ a,b \right ] $ splňující $f(x_n)> G-\frac{1}{n}$. Pomocí Weierstrassovy
	věty nalezneme vybranou posloupnost $\left \{ x_{n_k} \right \}_{k=1}^\infty $ z
	$\left \{ x_n \right \}_{n=1}^\infty $ takovou, že $\lim_{n \to \infty} x_{n_k}\in \mathbb R $.
	Potom platí, že $x^* = \lim x_{n_k} \in \left [ a,b \right ] $. Dále máme, že
	$\lim_{n \to \infty} f(x_n)=G$, neboť
	$$
	    \forall n \in \mathbb N:G-\frac{1}{n} < f(x_n) \leq G,
	$$
	a tedy $\lim_{n \to \infty} f(x_{n_k})=G$. Podle Heineovy věty máme, že $\lim_{n \to \infty} f(x_{n_k})=f(x^*)$.
	Odtud $f(x^*)=G$.
\end{proof}

\begin{veta}[O inverzní funkci]\label{veta3.15} % 3.15
Nechť $f$ je spojitá a rostoucí (resp. klesající) funkce na intervalu $J$. Potom funkce
$f^{-1}$ je spojitá a rostoucí (resp. klesjící) na intervalu $f(J)$.
\end{veta}

\begin{proof}
	Funkce $f$ je rostoucí, a tedy prostá. Funkce $f^{-1}$ je tedy dobře definovaná na $f(J)$.
	Množina $f(J)$ je interval podle věty \ref{veta3.12}. Je-li $f$ rostoucí (resp. klesající),
	pak je snadné ukázat, že $f^{-1}$ je rostoucí (resp. klesající). Předpokládejme, že $y_0\in f(J)$
	není pravým koncovým bodem intervalu $f(J)$. Budeme dokazovat, že $f^{-1}$ je spojitá zprava
	v bodě $y_0$. Označme $x_0= f^{-1}(y_0)$. Zvolme $\varepsilon \in \mathbb R , \varepsilon > 0$.
	Bod $y_0$ není pravý koncový bod $f(J)$, a tedy $x_0$ není pravý koncový bod $J$, neboť $f$
	je rostoucí. Nalezneme $x_1\in (x_0,x_0+\varepsilon) \cap J$. Nalezneme $\delta \in \mathbb R , \delta > 0$
	takové, že $f(x_1)=y_0+ \delta$. Pak pro každé $y\in \left [ y_0,y_0+\delta \right ) $
	platí $x_0=f^{-1}(y_0)\leq f^{-1}(y)\leq f^{-1}(y_0+\delta)=x_1$, takže
	$f^{-1}(y) \in [x_0,x_1) \subseteq B(x_0,\varepsilon)= B(f^{-1}(y_0),\varepsilon)$. Spojitost
	zleva v bodě, který není levým koncovým bodem $f(J)$, by se dokazovala analogicky.
\end{proof}

\section{Elementární funkce}
\begin{veta} % 4.1
	% TODO
\end{veta}

\subsection{Vlastnosti exponenciální funkce}
\begin{enumerate}
	\item[(E3)] $\underline{\exp(0)=1}$
	\par Podle (E1) platí $\exp(0) = \exp(0+0) = \exp(0)\cdot \exp(0)$, neboli $\exp(0)=\exp(0)^2$, tedy
	$\exp(0)$ je buď 0 nebo 1. Kdyby $\exp(0)=0$, potom pro každé $x\in \mathbb R $ platí
	$$
	    \exp(x) = \exp(x+0)=\exp(x)\cdot \exp(0) = 0,
	$$
	což je spor s (E2).
	\item[(E4)] $\underline{\forall x \in \mathbb R : \exp(x)\cdot \exp(-x)=1}$
	\par Pro $x\in \mathbb R $ platí:
	$$
	    1=\exp(0)=\exp(x+(-x))=\exp(x)\cdot \exp(-x).
	$$
	\item[(E5)] $\underline{\forall n \in \mathbb Z \, \forall x \in \mathbb R : \exp(nx)=\exp(x)^n}$
	\begin{itemize}
		\item $n=0$: pro $x \in \mathbb R $ je $\exp(0\cdot x)=\exp(0)=1=\exp(x)^0$
		\item $n \rightsquigarrow n+1$: pro $x \in \mathbb R $ je
		\begin{align*}
			\exp((n+1)x) & =\exp(nx+x) = \exp(nx)\cdot \exp(x) = \exp(x)^n \cdot \exp(x)\\
			&=\exp(x)^{n+1}
		\end{align*}
		\item $n\in \mathbb Z , n < 0$: pro $x \in \mathbb R $ je
		$$
			\exp(nx) = \exp(-(-nx)) = \frac{1}{\exp(-nx)}=\frac{1}{\exp(x)^{-n}} = \exp(x)^n
		$$
	\end{itemize}
	\item[(E6)] $\underline{\forall x \in(0,\infty):\exp(x)>1}$
	\par Z vlastnosti (E2) plyne, že existuje $\delta \in \mathbb R , \delta > 0$ takové, že
	$$
	    \forall z \in (0,\delta]: \exp(z)-1 > 0.
	$$
	Položme $k=\left [ \frac{x}{\delta} \right ] $. Potom můžeme psát
	$$
	    x=k\delta + z, \,\,\, z \in [0,\delta).
	$$
	Potom máme
	$$
	    \exp(x) = \exp(k\delta + z)=\exp(k\delta) \cdot \exp(z) = \underbrace{\exp(\delta)^k}_{>1 \textrm{ pro } k\geq 1}\cdot \exp(z)>1,
	$$
	pokud $k=0$, pak $z>0$.
	\item[(E7)] $\underline{\forall x \in (-\infty, 0): \exp(x)\in(0,1)}$
	\par Plyne z (E4) a (E6).
	\item[(E8)] \underline{$\exp$ je rostoucí na $\mathbb R $}
	\par Zvolme $x,y \in \mathbb R , x<y$. Položme $k=y-x$. Potom $k>0$ a platí
	$$
	    \exp(y) = \exp(x+k) = \underbrace{\exp(x)}_{>0}\cdot \underbrace{\exp(k)}_{>1}>\exp(x).
	$$
	\item[(E9)] $\underline{\lim_{x \to \infty} \exp(x)=\infty}$
	\par Podle věty % TODO 3.9 (A)
	je $\lim_{x \to \infty} \exp(x)$ existuje, neboť $\exp$ je rostoucí. Podle Heineovy věty
	platí, že $\lim_{n \to \infty} \exp(n)=\lim_{x \to \infty} \exp(x)$. Navíc máme
	$$
	    \lim_{n \to \infty} \exp(n)=\lim_{n \to \infty} (\exp(1))^n = \infty, \,\,\,\textrm{(E6)} \implies \exp(1) > 1.
	$$
	\item[(E10)] $\underline{\lim_{x \to -\infty} \exp(x)=0}$
	\par Plyne z (E9) a (E4):
	$$
    \lim_{x \to -\infty} \exp(x)=\lim_{x \to -\infty} \frac{1}{\exp(-x)}=\frac{1}{\infty}=0
	$$
	\begin{align*}
		&g(x) = x & \lim_{x \to -\infty} g(x)&=\infty \\
		&f(y) = \exp(y) & \lim_{y \to \infty} f(y)&=\infty \\
		&\textrm{je splněna podmínka (P)}
	\end{align*}
	\item[(E11)] \underline{$\exp$ je spojitá na $\mathbb R$}
	\par spojitá v 0:
	$$
    \lim_{x \to 0} \exp(x) = \lim_{x\to 0} \left( \frac{\exp(x)-1}{x}\cdot x+1 \right ) = 1\cdot 0 + 1= 1=\exp(0)
	$$
	\par spojitá na $c\in \mathbb R $:
	$$
        \lim_{x \to c} \exp(x) = \lim_{x\to c} \exp(x-c+c) = \lim_{x \to c} \exp(x-c)\cdot \exp(c)=1 \cdot \exp(c) = \exp(c)
	$$
	\item[(E12)] \underline{$\mathcal H(\exp)=(0,\infty)$}
	\par Díky větě \ref{veta3.12} a (E11) máme, že $\mathcal H(\exp)=\exp(\mathbb R )$ je interval.
	Víme, že platí $\mathcal H(\exp)\subseteq (0,\infty)$. Z (E9) a (E10) plyne $\mathcal H(\exp)=(0,\infty)$.
\end{enumerate}

\begin{poznamka}[Označení]
	Eulerovo číslo $e=\exp(1)$.
\end{poznamka}

\begin{definice}
	% TODO
	logaritmus
\end{definice}

\subsection{Vlastnosti logaritmu}
\begin{enumerate}
	\item[(L1)] \underline{$\mathcal D(\log)=(0,\infty)$}
	\par Platí totiž $\mathcal H(\exp) = (0,\infty)$.
	\item[(L2)] \underline{$\mathcal H(\log)=\mathbb R $}
	\par Platí totiž $\mathcal D(\exp)=\mathbb R $.
	\item[(L3)] \underline{$\log(1)=0$}
	\par Je totiž $\exp(0)=1$.
	\item[(L4)] \underline{$\forall x,y \in(0,\infty):\log(xy)=\log x + \log y$}
	\par Pro $x,y\in (0,\infty)$ označme $a=\log x, b = \log y$. Potom
	$$
	    \log(xy)=\log(\exp(a)\exp(b))=\log(\exp(a+b))=a+b=\log x + \log y.
	$$
	\item[(L5)] \underline{$\forall x \in (0,\infty): \log \frac{1}{x} = - \log x$}
	\par Platí
	$$
	    0=\log (1) = log \left( x\cdot \frac{1}{x} \right ) = \log (x)+ \log \left( \frac{1}{x} \right ),
	$$
	odkud máme $\log \frac{1}{x}=-\log x$.
	\item[(L6)] \underline{$\forall n \in \mathbb Z  \, \forall x \in (0,\infty):\log (x^n)=n\log x$}
	\begin{itemize}
		\item $n=0:$ platí
		\item $n \rightsquigarrow n+1:$ ($n\in \mathbb N \cup \left \{ 0 \right \}) , x \in (0,\infty)$
		$$
		    \log(x^{n+1})=\log(x^n\cdot x) = \log(x^n)+ \log (x) = n\log x + \log x = (n+1)\log x.
		$$
		\item $n\in \mathbb Z , n < 0:$ $x \in (0,\infty)$
		$$
		    \log(x^n)=\log \left( \frac{1}{x^{-n}} \right ) = -\log (x^{-n}) = (-1) \cdot (-n)\cdot \log x=n\log x.
		$$
	\end{itemize}
	\item[(L7)] \underline{$\log$ je rostoucí na $(0,\infty)$}
	\par $\exp$ je rostoucí na $\mathbb R $
	\item[(L8)] \underline{$\lim_{x \to \infty} \log x=\infty$}
	\par Z věty o limitě monotónní funkce, faktu že $\log $ je rostoucí a toho, že
	$\mathcal H(\log) = \mathbb R $ plyne, že $\lim_{x \to \infty} \log x = \infty$.
	\item[(L9)] \underline{$\lim_{x \to 0^+} \log x = -\infty$}
	\par Podobně jako v (L8).
	\item[(L10)] \underline{$\log$ je spojitá na $(0,\infty)$}
	\par Plyne ze spojitosti $\exp$, (E5) a věty \ref{veta3.15}.
	\item[(L11)] \underline{$\lim_{x \to 1} \frac{\log x}{x-1}=1$}
	\par Vezměme si
	\begin{align*}
		 f(y)&=\frac{\exp(y)-1}{y} & & \lim_{y \to 0} f(y)=1\\
		 g(x) &=\log x & & \lim_{x \to 1} g(x) = 0
			\textrm{je splněna podmínka (P)},
	\end{align*}
	odkud máme
	$$
        \lim_{x \to 1} f\circ g(x) = 1,
	$$
	tedy že
	$$
	\lim_{x\to 1} \frac{\exp(\log x)-1}{\log x} = \lim_{x\to 1} \frac{x-1}{\log x} = 1.
	$$
\end{enumerate}

\begin{poznamka}\,
\begin{enumerate}[(i)]
	\item Pro $a \in \mathbb R , a > 0, n \in \mathbb Z $ je
	$$
	    a^n= \exp(n\log a)=\exp(\log(a^n))=a^n,
	$$
	tedy definice obecné mocniny je kompatibilní s (L6).
	\item Pro $n\in \mathbb N $ je funkce
	$$
	    x\mapsto \sqrt[n]{x}
	$$
	\begin{itemize}
		\item pro $n$ sudé rovna inverzní funkci k $f(x)=x^n, x \in [0,\infty)$,
		\item pro $n$ liché rovna inverzní funki k $h(x)=x^n, x \in \mathbb R $.
	\end{itemize}
	Platí $x\mapsto \sqrt[n]{x}$ je spojitá na svém definičním oboru (podle věty \ref{veta3.15}).
\end{enumerate}
\end{poznamka}

\begin{veta}[Základní vlastnosti sinu a kosinu]\label{veta3.17} % veta 3.17
	% TODO
\end{veta}

\subsection{Vlastnosti sinu a kosinu}
\begin{enumerate}
	\item[(G6)] $\underline{\cos(0)=1}$
	\par Zvolme $x\in (0,\frac{\pi}{2})$. Podle (G1) platí
	$$\sin x = \sin(x+0)=\sin x \cos 0+\cos x \sin 0 = \sin x \cos 0,$$
	a protože $\sin x \ne 0$ pro $x\in (0,\frac{\pi}{2})$, je
	$$
	    \cos 0 = 1.
	$$
	\item[(G7)] $\underline{\forall x \in \mathbb R : \sin^2 x + \cos ^2 x = 1}$
	\par Platí
	$$
	    1= \cos(0)=\cos(x-x)=\cos x \cos(-x)-\sin x \sin(-x)=\cos^2 x+ \sin^2 x.
	$$
	\item[(G8)] $\underline{\forall x \in \mathbb R : |\sin x|\leq 1 \land |\cos x| \leq 1}$
	\par Plyne z (G7).
	\item[(G9)] $\underline{\cos\left (\frac{\pi}{2}\right)=0}$
	\par Podle (G7) máme
	$$
	    \sin^2 \left( \frac{\pi}{2} \right )  + \cos^2 \left( \frac{\pi}{2} \right ) = 1,
	$$
	ale $\sin^2 \left( \frac{\pi}{2} \right )$ je podle (G4) rovno 1, takže
	$$
	    \cos^2 \left( \frac{\pi}{2} \right ) = 0 \implies \cos \left( \frac{\pi}{2} \right ) =0.
	$$
	\item[(G10)] $\underline{\forall x \in \mathbb R : \sin \left( x+\frac{\pi}{2} \right )  = \cos x}$
	\par Z (G1) a (G9) plyne
	$$
	    \sin \left( x+\frac{\pi}{2} \right ) = \sin x \cos \left( \frac{\pi}{2} \right ) +\cos x \sin \left( \frac{\pi}{2} \right ) =\cos x.
	$$
	\item[(G11)] $\underline{\forall x \in \mathbb R : \cos \left( x+\frac{\pi}{2} \right ) = - \sin x }$
	\par Z (G2) plyne
	$$
	    \cos \left( x+ \frac{\pi}{2} \right ) = \cos x \cos \frac{\pi}{2} - \sin x \sin \frac{\pi}{2} = -\sin x.
	$$
	\item[(G12)] \underline{sinus a kosinus jsou $2\pi$ periodické}
	\par Z (G10) a (G11) plyne
	\begin{align*}
	\sin(x+2\pi) &= \sin \left( x+\frac{3}{2}\pi + \frac{\pi}{2} \right )  = \cos \left( x+ \frac{3}{2}\pi \right )\\
& = - \sin(x+\pi) = -\cos \left( x+\frac{\pi}{2} \right ) = \sin x.
\end{align*}
	\par $\cos(x+2\pi)=\cos x$ analogicky.
	\item[(G13)] \underline{sinus a kosinus jsou spojité na $\mathbb R $}
	\begin{itemize}
		\item sinus je spojitý v nule:
		$$
  \lim_{x \to 0} \sin x = \lim_{x\to 0} \frac{\sin x}{x}\cdot x = 1\cdot 0 = 0 = \sin 0.
		$$
		\item kosinus je spojitý v nule:
		\begin{align*}
		\lim_{x\to 0} \cos x & = \lim_{x\to 0} \cos \left( \frac{x}{2}+\frac{x}{2} \right ) = \lim_{x\to 0} \left( \cos^2 \left( \frac{x}{2} \right ) - \sin^2 \left( \frac{x}{2} \right ) \right) \\
	&= \lim_{x\to 0} \left( 1-2\sin^2 \left( \frac{x}{2} \right ) \right ) = 1-2\cdot 0^2 = 1 = \cos 0.
\end{align*}
		\item spojitost sinu v $x_0\in \mathbb R $:
		\begin{align*}
		\lim_{x\to x_0} \sin x & = \lim_{x\to x_0} \sin \left( \left( x-x_0 \right ) +x_0 \right )\\
	& = \lim_{x\to x_0} \left( \sin(x-x_0)\cos x_0 + \cos (x-x_0)\sin x_0  \right )\\
	& = 0 \cdot \cos x_0 + \cos 0 \cdot \sin x_0 = \sin x_0.
\end{align*}
		\item spojitost kosinu v $x_0 \in \mathbb R $: obdobně s pomocí (G2)
	\end{itemize}
	\item[(G14)] \underline{$\left \{ x\in \mathbb R , \sin x = 0 \right \} = \left \{ k\pi, k \in \mathbb Z  \right \}  $ a $\left \{ x\in \mathbb R , \cos x = 0 \right \} = \left \{ \frac{\pi}{2}+k\pi, k \in \mathbb Z  \right \}  $}
	\par Platí
	\begin{align*}
		\forall x \in \left( 0, \frac{\pi}{2} \right ) : \sin x  > 0 &&  \forall x \in \left( -\frac{\pi}{2}, 0 \right ) :\sin x < 0.
	\end{align*}
	Dále pro $x \in \left( \frac{\pi}{2},\pi \right ) $ je
	$$
	    \sin x = \cos \left( x-\frac{\pi}{2} \right ) = - \sin (x-\pi).
	$$
	Podobně když $x\in \left( \pi, \frac{3\pi}{2} \right ) $, máme $\sin x = - \sin (x-\pi)$.
	Odtud dostáváme, že nulové body jsou právě v $0, \frac{\pi}{2}, \pi$ a $\frac{3\pi}{2}$.
	Pro kosinus dostaneme z (G11) posunutím o $\frac{\pi}{2}$.
\end{enumerate}

\begin{definice} % Tangens a kotangens
	% TODO
\end{definice}

\begin{definice} % cyklometrické funkce
	% TODO
\end{definice}

\begin{veta}\label{veta3.18}
	% TODO
\end{veta}

\begin{proof}
	Spojitosti funkcí $\log, \exp, \sin, \cos$ jsme již dokázali. Funkce $\tg$ a $\cotg$ dostaneme
	pomocí věty o artimetice limit. Zbývá tedy dokázát spojitost funkcí $\arcsin$, $\arccos$, $\arctg$, $\arccotg$.
	Do dostaneme podle věty \ref{veta3.15}.
\end{proof}

\section{Derivace}
\subsection{Definice a základní vztahy}
\begin{definice}
Nechť $a\in\R$ a $f$ je funkce. Existuje-li limita
\[
\lim_{h\to 0} \frac{f(a+h)-f(a)}{h},
\]
pak její hodnotu nazýváme \textbf{derivací} funkce $f$ v bodě $a$. Obdobně definujeme \textbf{derivaci zprava} a \textbf{derivaci zleva} funkce $f$ v bodě $a$ předpisy
\[
\lim_{h\to 0^+} \frac{f(a+h)-f(a)}{h},\qquad \lim_{h\to 0^-} \frac{f(a+h)-f(a)}{h}.
\]
Derivaci zleva v bodě a derivaci zprava v bodě nazýváme \textbf{jednostrannými derivacemi}.
\end{definice}

\begin{oznaceni}
	Pokud derivace existuje, je určena jednoznačně (důsledek věty \ref{3.1}). Používáme značení
	$f^\prime(a)$, $f_+^\prime(a)$, $f_-^\prime(a)$.
\end{oznaceni}

\begin{poznamka}\,
\begin{enumerate}[$i$.]
	\item Derivace
	$$f^\prime(a)
 \begin{cases}
		\textrm{existuje}, \\
		\textrm{neexistuje} \,\,\,\begin{cases}
			\textrm{je vlastní, tj. }\in \mathbb R, \\
			\textrm{je nevlastní} \,\,\, \begin{cases}
				\infty, \\
				-\infty.
\end{cases}
\end{cases}
 \end{cases}
	$$
	\item $f^\prime (a) \implies \exists \delta \in \mathbb R , \delta > 0: f$ je definována na okolí $B(a,\delta)$
	\item $f^\prime(a) = A\in \mathbb R ^* \iff f_+^\prime(a) = f_-^\prime (a) = A\in \mathbb R ^*$
	\item $\lim_{x\to a} \frac{f(x)-f(a)}{x-a} = \lim_{h\to 0} \frac{f(a+h)-f(a)}{h}$
	\item Derivace je lokální pojem.
\end{enumerate}
\end{poznamka}

\begin{poznamka}[Derivace a tečna]
	Nechť $a\in \mathbb R $, $f$ je funkce a $f^\prime(a)\in \mathbb R $. Tečnou ke grafu $f$ v bodě
	$\left [ a,f(a) \right ] $ rozumíme funkci
	$$
	   t: x\mapsto f(a)+ f^\prime(a)(x-a).
	$$
	Platí
	\begin{align*}
	\lim_{x\to a} \frac{f(x)-t(x)}{x-a}&= \lim_{x\to a} \frac{f(x)-f(a)-f^\prime(a)(x-a)}{x-a}\\
	&= \left( \frac{f(x)-f(a)}{x-a}-f(a) \right )= f^\prime(a)-f^\prime(a) = 0.
	\end{align*}
\end{poznamka}

\begin{poznamka}
	Definičním oborem derivace funkce $f$ rozumíme množinu
	$\left \{ x \in \mathbb R , f^\prime (x)\in \mathbb R  \right \} $.
\end{poznamka}

\begin{priklad}\,
\begin{enumerate}[$i.$]
	\item $f(x)=c, c\in B(a, \delta)$, kde $\delta>0, a \in \mathbb R $. Platí
	$$
	    f^\prime(a) = \lim_{x\to a} \frac{f(x)-f(a)}{x-a} = \lim_{x\to a} \frac{c-c}{x-a} = 0.
	$$
	\item $f(x)=x^m$, kde $m\in\mathbb{N}$, $x\in\mathbb{R}$:
	\par  Pro $a\in\mathbb{R}$:
    \[
    f'(a)=\lim_{x\to a}\frac{x^m-a^m}{x-a}=\lim_{x\to a}\frac{(x-a)\left(x^{m-1}+x^{m-2}a+\cdots+a^{m-1}\right)}{x-a}=na^{n-1}.
    \]
    \item $f(x)=\operatorname{sgn}x$, $x\in\mathbb{R}$
    \begin{align*}
    f'_+(0)&=\lim_{x\to0^+}\frac{\operatorname{sgn}x-\operatorname{sgn}0}{x-0}
    =\lim_{x\to0^+}\frac{1}{x}=+\infty\\
    f'_-(0)&=\lim_{x\to0^-}\frac{\operatorname{sgn}x-\operatorname{sgn}0}{x-0}
    =\lim_{x\to0^-}\frac{-1}{x}=+\infty,
    \end{align*}
    odkud máme, že $\operatorname{sgn}^\prime(0)=\infty.$
    \item $f(x)=|x|$
    \[
    f'_+(0)=1, \qquad f'_-(0)=-1,
    \]
    a proto derivace $f'(0)$ neexistuje.
\end{enumerate}
\end{priklad}

\begin{veta}[Derivace a spojitost]\label{veta5.1}
	Nechť funkce $f$ má v bodě $a \in \mathbb R $ vlastní derivaci. Potom je $f$ v bodě $a$ spojitá.
\end{veta}

\begin{proof}
	Počítejme
    $$
	\lim_{x\to a} f(x) = \lim_{x\to a} \left( \frac{f(x)-f(a)}{x-a}(x-a)+f(a) \right ) =\underbrace{f^\prime(a)}_{\in \mathbb R }\cdot 0+f(a)=f(a).
	$$
\end{proof}

\begin{poznamka}\,
\begin{enumerate}[(i)]
	\item Věta \ref{veta5.1} platí i v jednostranné verzi.
	\item Předpoklad, že $f^\prime(a)$ je vlastní, nelze vynechat (vizte funkci signum v nule).
\end{enumerate}
\end{poznamka}

\begin{veta}[Aritmetika derivací]\label{veta5.2}
Nechť $a \in \mathbb{R}$ a $f, g$ jsou funkce, které mají v bodě $a$ derivaci.
\begin{itemize}
\item[(a)] Pak platí $$(f+g)'(a)=f'(a)+g'(a),$$ je-li výraz na pravé straně definován.
\item[(b)] Nechť alespoň jedna z funkcí $f$ a $g$ je spojitá v bodě $a$.
Potom platí $$(fg)'(a)=f'(a)g(a)+f(a)g'(a),$$
je-li výraz na pravé straně definován.
\item[(c)] Nechť funkce $g$ je spojitá v bodě $a$. Potom platí
$$
\left(\frac{f}{g}\right)'(a)=\frac{f'(a)g(a)-f(a)g'(a)}{g^2(a)},$$
je-li výraz na pravé straně definován.
\end{itemize}
\end{veta}

\begin{proof}\,
\begin{enumerate}[(a)]
	\item Platí
	\begin{align*}
	(f+g)'(a) & = \lim_{h\to 0} \frac{(f+g)(a+h)-(f+g)(a)}{h}\\
	&= \lim_{h\to 0} \frac{f(a+h)+g(a+h)-f(a)-g(a)}{h}\\
	&= \lim_{h\to 0} \left( \frac{f(a+h)-f(a)}{h} + \frac{g(a+h)-g(a)}{h} \right)\\
 &= f'(a) + g'(a).
\end{align*}
\item Předpokládejme, že $g$ je spojitá v $a$. Pak
\begin{align*}
(fg)'(a) &= \lim_{h\to 0} \frac{(fg)(a+h)-(fg)(a)}{h}\\
&= \lim_{h\to 0} \frac{f(a+h)g(a+h)-f(a)g(a)}{h}\\
&= \lim_{h\to 0} \frac{(f(a+h)-f(a))g(a+h)+f(a)(g(a+h)-g(a))}{h}\\
&= f'(a)\,g(a) + f(a)\,g'(a).
\end{align*}
\item Platí
\begin{align*}
\left(\frac{f}{g}\right)'(a) &= \lim_{h\to 0} \frac{\left( \frac{f}{g} \right ) (a+h)-\left( \frac{f}{g} \right ) (a)}{h}\\
	&= \lim_{h\to 0} \frac{1}{h}\left( \frac{f(a+h)}{g(a+h)} - \frac{f(a)}{g(a)} \right)\\
	&= \lim_{h\to 0} \frac{1}{h}\,\frac{f(a+h)g(a)-f(a)g(a+h)}{g(a+h)g(a)}\\
	&= \lim_{h\to 0} \frac{1}{g(a+h)g(a)}\left( \frac{f(a+h)-f(a)}{h}g(a) - f(a)\frac{g(a+h)-g(a)}{h} \right)\\
	&= \frac{f'(a)g(a)-f(a)g'(a)}{g(a)^2}.\qedhere
\end{align*}
\end{enumerate}
\end{proof}

\begin{veta}[Derivace složené funkce] \label{veta5.3}
	Nechť funkce $g$ je spojitá v bodě $a \in \mathbb R$, má v tomto bodě derivaci
	a funkce $f$ má derivaci v bodě $g(a)$. Pak platí
	$$
	    (f\circ g)^\prime(a) = f^\prime(g(a))\cdot g^\prime(a),
	$$
	je-li výraz na pravé straně definován.
\end{veta}

\begin{proof}
	Označme $b=g(a)$,
	$$
	\varphi(y)=\frac{f(y)-f(b)}{y-b}, \quad \mathcal D(f)\smallsetminus \left \{ b \right \} .
	$$
	Platí $\lim_{y\to b} \varphi(y)=f^\prime(b)$. Rozlišíme dva případy
	\begin{enumerate}[(i)]
		\item $g^\prime (a) \ne 0$: nalezneme $\eta > 0$ takové, že platí
		$$
		    \forall x in P(a,\eta):\frac{g(x)-g(a)}{x-a} \ne 0.
		$$
		odtud plyne, že
		$$
		    \forall x \in P(a, \eta): g(x)\ne g(a).
		$$
		Počítejme
		\begin{align*}
			(f\circ g)^\prime (a) &= \lim_{x\to a} \frac{f(g(x))-f(g(a))}{x-a}=\lim_{x\to a} \frac{f(g(x))-f(g(a))}{g(x)-g(a)}\cdot \frac{g(x)-g(a)}{x-a}\\
			&=\lim_{x\to a} \varphi(g(x))\cdot \frac{g(x)-g(a)}{x-a}.
		\end{align*}
		\item $g^\prime(a) = 0$: Výraz na pravé straně je definován, a proto
		$f^\prime(b)\in \mathbb R $. Nalezneme $c>0, \delta>0$ taková, že
		$$
		    \forall y \in P(b,\delta): |\varphi(y)|<c \quad (\textrm{věta }\ref{veta3.2}).
		$$
		Dále
		\begin{align*}
			& \forall y \in P(b,\delta): \left | \frac{f(y)-f(b)}{y-b} \right |<c,\\
			& \forall y \in B(b,\delta): |f(y)-f(b)| \leq c\cdot |y-b|.
\end{align*}
    Zvolme $\varepsilon>0$. Nalezneme $\rho_1>0$ takové, že
    $$
        \forall x \in P(a,\rho_1): \left | \frac{g(x)-g(a)}{x-a}-0 \right|<\varepsilon.
    $$
    Dále nalezneme $\rho_2 \in (0,\rho_1)$ takové, že
    $$
        \forall x \in B(a,\rho_2): g(x)\in B(b,\delta).
    $$
    Pro každé $x \in P(a,\rho_2)$ platí
    $$
	\left | \frac{f(g(x))-f(g(a))}{x-a} \right | \leq \frac{c|g(x)-g(a)|}{|x-a|} = c\cdot \left | \frac{g(x)-g(a)}{x-a} \right |<c\varepsilon,
    $$
    odkud plyne $(f\circ g)^\prime(a)=0=f^\prime(b)\cdot 0 = f^\prime(b)\cdot g^\prime(a)$.\qedhere
\end{enumerate}
\end{proof}

\begin{poznamka}
	V čem je problém:
	\begin{align*}
    (f\circ g)^\prime (a) &= \lim_{x\to a} \frac{f(g(x))-f(g(a))}{x-a} = \lim_{x\to a} \frac{f(g(x))-f(g(a))}{g(x)-g(a)}\cdot \frac{g(x)-g(a)}{x-a}\\
    &= f^\prime(b)\cdot g^\prime(a).
    \end{align*}
    Funkce
    \begin{align*}
	&y\mapsto \frac{f(y)-f(b)}{y-b} && \lim_{y\to b} \frac{f(y)-f(b)}{y-b}=f^\prime(b)-f^\prime(g(a))\\
	&g(a)=b && \lim_{x\to a} g(x) = g(a)=b
    \end{align*}
\end{poznamka}

\begin{veta}[Derivace inverzní funkce] \label{veta5.4}
Nechť $a$ je vnitřním bodem intervalu $I$ a $f$ je spojitá a ryze monotónní funkce na $I$. Označme $b=f(a)$.
\begin{enumerate}[(a)]
  \item Nechť $f'(a)$ existuje a je nenulová. Potom platí
  $$(f^{-1})'(b)=\dfrac{1}{f'(a)}.$$
  (Je-li $f'(a)=\pm\infty$, pak $(f^{-1})'(b)=0$.)
  \item Nechť $f'(a)=0$ a $f$ je rostoucí na $I$. Potom $(f^{-1})'(b)=\infty$.
  \item Nechť $f'(a)=0$ a $f$ je klesající na $I$. Potom $(f^{-1})'(b)=-\infty$.
\end{enumerate}
\end{veta}

\begin{proof}
	Funkce $f:I\to \mathbb R $. Označme $J=f(I)$. Podle věty \ref{veta3.12} je $J$ interval.
	Inverzní funkce $f^{-1}$ je tedy $J\to I$. Podle věty \ref{veta3.15} je $f^{-1}$ spojitá
	a ryze monotónní. Označme
	$$
	    \varphi(x) = \frac{f(x)-f(a)}{x-a}, \quad x \in \mathcal D(f)\smallsetminus \left \{ a \right \}.
	$$
	Platí $\lim_{x\to a} \varphi(x)=f^\prime(a)$ v případech (a) až (c). Dále platí
	$\lim_{y\to b} f^{-1}(y)=f^{-1}(b)=a$, neboť $b$ je vnitřním bodem $J$ a $f^{-1}$ je spojitá
	na intervalu $J$.
	\begin{enumerate}[(a)]
		\item $\lim_{y\to b} (\varphi \circ f^{-1})(y)$, a je splněna podmínka (P), protože
		$f^{-1}$ je ryze monotónní. Máme tedy
		$$
  \lim_{y\to b} \frac{f(f^{-1}(y))-f(a)}{f^{-1}(y)-a} = \lim_{y\to b} \frac{y-b}{f^{-1}(y)-f^{-1}(b)},
		$$
		tedy
		$$
		    (f^{-1})^\prime(y) = \lim_{y\to b} \frac{f^{-1}-f^{-1}(b)}{y-b} = \lim_{y\to b} \frac{1}{\frac{y-b}{f^\prime(y)-f^\prime(b)}}=\frac{1}{f^\prime(a)}.
		$$
		\item Podobně jako v (a) máme, že
		$$\lim_{y\to b} \frac{y-b}{f^{-1}(y)-f^{-1}(b)}=f^\prime (a) = 0.$$
		Potom platí podle věty \ref{veta3.4}
		$$
		    (f^{-1})^\prime = \lim_{y\to b} \frac{1}{\frac{y-b}{f^{-1}(y)-f^{-1}(b)}}=\infty.
		$$
		\item Analogicky jako (b).\qedhere
	\end{enumerate}
\end{proof}

\paragraph{Derivace elementárních funkcí}
\begin{itemize}
	\item $f(x)=c, x \in \mathbb R$
	\par $f^\prime(x) = 0, x \in \mathbb R $
	\item $\left( x^n \right )^\prime = nx^{n-1}, x \in \mathbb R , n \in \mathbb N  $
	\item $\left( \exp(x) \right )^\prime  = \exp(x), x \in \mathbb R $
	$$
 \left( \exp(x) \right )^\prime = \lim_{h\to 0} \frac{\exp(x+h)-\exp(x)}{h} = \lim_{h\to 0} \exp (x)\cdot \frac{\exp(h)-1}{h} = \exp(x)
	$$
	\item $(\log x)^\prime = \frac{1}{x}, x \in \left( 0,\infty \right ) $
	$$
	    (\log x)^\prime =\frac{1}{\exp(\log x)} = \frac{1}{x}.
	$$
	\begin{align*}
		f(z) &= \exp(z) \\
		f^{-1}(x) & = \log x
	\end{align*}
	\item $(x^a)^\prime = ax^{a-1}, x \in \left( 0,\infty \right ), a \in \mathbb R  $
	$$
 \left( x^a \right )^\prime = \left( e^{a \log x} \right )^\prime = e^{a\log x}\cdot \frac{a}{x} = ax^{a-1}.
	$$
	\item $\left( a^x \right )^\prime = a^x \log a, x \in \mathbb R , a > 0 $
	$$
 \left( a^x \right )^\prime = \left( e^{x\log a} \right )^\prime = e^{x\log a}\cdot \log a = a^x\log a.
	$$
\end{itemize}

\begin{poznamka}
	Funkce $x^x$ se nehodí derivovat podle žádného z předchozích bodů, ale
$$\left( x^x \right )^\prime = \left( e^{x\log x} \right )^\prime = e^{x\log x}\cdot (x\log x + x\cdot \frac{1}{x}) = x^x (\log x +1), x > 0.  $$
\end{poznamka}

\begin{itemize}
	\item $\left( \sqrt[n]{x} \right )^\prime = \frac{1}{n} \cdot \frac{\sqrt[n]{x}}{x} $ pro $x \in (0,\infty)$ pro $n \in \mathbb N $ sudé, nebo $x \in \mathbb R \smallsetminus \{0\}$ pro $n \in \mathbb N $ liché.
	Dále $(\sqrt[n]{x})^\prime = \infty$ pro $x=0$ a $n \in \mathbb N $ liché,
	$(\sqrt[n]{x})_+^\prime = \infty$ pro $x=0$ a $n\in \mathbb N $ sudé.
	\item $(\sin x )^\prime = \cos x, x \in \mathbb R $
	\begin{align*}
	(\sin x )^\prime &=\lim_{h\to 0} \frac{\sin(x+k)-\sin x}{h} = \lim_{h\to 0} \frac{\sin x \cos h + \cos x \sin h-\sin x}{h} \\
	&= \lim_{h\to 0} \left( \sin x \cdot \frac{\cos h - 1}{h} + \cos x \cdot \frac{\sin h }{h} \right ),
	\end{align*}
	ale
	$$
 \lim_{h\to 0} \frac{\cos h -1}{h} = \lim_{h\to 0} \frac{\cos h-1}{h^2}\cdot h = \left( -\frac{1}{2} \right )\cdot 0,
	$$
	tedy $(\sin x )^\prime = \cos x.$
	\item $(\cos x)^\prime = -\sin x \, x \in \mathbb R $
	\item $(\tg x)^\prime = \frac{1}{\cos^2 x}, x \in \mathcal D(\tg)$
	\item $(\cotg x )^\prime = -\frac{1}{\sin^2 x}, x \in \mathcal D(\cotg)$
	$$
	    (\cotg x) ^\prime = \left( \frac{\cos x}{\sin x} \right )^\prime = \frac{-\sin x \cdot \sin x - \cos x \cos x}{\sin^2 x} = - \frac{1}{\sin^2 x}.
	$$
	\item $(\arcsin x)^\prime = \frac{1}{\sqrt{1-x^2}}, x \in (-1,1)$
	\begin{align*}
		f(z) &= \sin z, z \in \left [ -\frac{\pi}{2}, \frac{\pi}{2} \right ]  \\
		f^{-1}(x) &= \arcsin x, x \in \left [ -1,1 \right ] \\
		(\arcsin x)^\prime &= \frac{1}{\cos(\arcsin x)}, \quad x \in (-1,1)\\
		\cos^2 y &= 1-\sin^2 y, \quad \ in \left( -\frac{\pi}{2}, \frac{\pi}{2} \right ) \implies \cos y = \sqrt{1-\sin ^2y} \\
		\cos(\arcsin x) &= \sqrt{1-\sin^2(\arcsin x)}=\sqrt{1-x^2}.
	\end{align*}
	\item $(\arccos x)^\prime = -\frac{1}{\sqrt{1-x^2}}, x \in (-1,1)$
	\item $(\arctg x)^\prime = \frac{1}{1+x^2},x \in \mathbb R $
	\item $(\arccotg x)^\prime = -\frac{1}{1+x^2}, x \in \mathbb R $
\end{itemize}

\begin{veta}[Nutná podmínka lokálního extrému] \label{veta5.5}
Jestliže $a$ je bodem lokálního extrému funkce $f$, potom buď $f'(a)$ neexistuje, nebo $f'(a)=0$.
\end{veta}

\begin{proof}
	Pro spor předpokládejme, že $f^\prime(a)$ existuje a $f^\prime(a)\ne 0$. Předpokládejme nejprve, že
	$f^\prime(a)>0$. Potom
	$$
	    f^\prime(a) = \lim_{x\to a} \frac{f(x)-f(a)}{x-a} > 0.
	$$
	Nalezneme $\delta > 0$ takové, že
	$$
	    \forall x \in P(a, \delta) : \frac{f(x)-f(a)}{x-a} > 0.
	$$
	Potom
	\begin{align*}
	    \forall x \in (a,a+\delta): f(x)>f(a), && \forall x \in (a-\delta, a): f(x) < f(a).
	\end{align*}
	Proto $a$ není bodem lokálního extrému, což je spor. Případ $f^\prime(a)<0$ lze přivést
	ke sporu analogicky.
\end{proof}

\begin{priklad}[Častá úloha]
	Nechť $f:\left [ a,b \right ] \to \mathbb R $ je spojitá na $\left [ a,b \right ] $ a
	pro každé $x \in (a,b)$ existuje $f^\prime(x)$. Nalezněte extrémy $f$ na $\left [ a,b \right ] $.
\end{priklad}

\begin{reseni}
	Protože je $f$ spojitá na $\left [ a,b \right ] $, má $f$ maximum i minimum.
	Dále ať $$E=\left \{ x\in (a,b), f^\prime(x)=0 \right \} \cup \left \{ a,b \right \}.  $$
	Extrémy mohou být jedině v $E$. Dále vypočteme funkční hodnoty prvků v $E$ a podle toho
	určíme maximum a minimum.
\end{reseni}

\subsection{Věty o střední hodnotě}
\begin{veta}[Rolle] \label{veta5.6}
Nechť $a,b \in \R, a<b$ a $f\colon[a,b]\to\R$ je spojitá funkce. Je-li $f(a)=f(b)$ a $f$ má derivaci v každém bodě intervalu $(a,b)$, pak existuje $c\in(a,b)$ takové, že $f'(c)=0$.
\end{veta}

\begin{proof}
	Podle věty \ref{veta3.14} $f$ nabývá v intervalu $\left [ a,b \right ] $ svého maxima i minima.
	Označme $m=\min f (\left [ a,b \right ])  $ a $M = \max f (\left [ a,b \right ])$.  Potom platí
	$$
	    m \leq f(a) = f(b) \leq M.
	$$
	\par Pokud $m=M$, pak je $f$ konstantní na $\left [ a,b \right ] $ a stačí volit $c \in (a,b)$ libovolně.
	\par Pokud $m < M$, pak platí $m< f(a)$ nebo $f(a)< M$. Pokud $f(a)<M$, pak existuje $c\in \left [ a,b \right ] $
	takové, že $f(c)=M.$ Podle předpokladů věty a podle věty \ref{veta5.5} máme $f^\prime(c)=0$.
	Pokud $m<f(a)$, postupujeme obdobně.
\end{proof}

\begin{veta}[Lagrange] \label{veta5.7}
Nechť $a,b \in \R, a<b$ a $f\colon[a,b]\to\R$ je spojitá funkce, která má v každém bodě intervalu $(a,b)$ derivaci. Pak existuje $c\in(a,b)$ takové, že
\[
 f'(c)=\frac{f(b)-f(a)}{b-a}.
\]
\end{veta}

\begin{proof}
Definujme pomocnou funkci
$$
    g(x)=f(x) - \frac{f(b)-f(a)}{b-a}(c-a), \quad x \in \left [ a,b \right ] .
$$
Platí, že
\begin{itemize}
	\item $g$ je spojitá na $\left [ a,b \right ] $,
	\item $g(a)= f(a)$, $g(b)=f(b)-\frac{f(b)-f(a)}{b-a}(b-a)=f(a)$,
\item $g^\prime(x)=\underbrace{f^\prime(x)}_{\in \mathbb R ^*}-\underbrace{\frac{f(b)-f(a)}{b-a}}_{\in \mathbb R }$, $x \in (a,b)$.
\end{itemize}
Podle Rolleovy věty \ref{veta5.6} nalezneme $c\in(a,b)$ takové, že $g^\prime(c)=0$. Potom platí
$$
    \underbrace{f^\prime(c)}_{\in \mathbb R }-\frac{f(b)-f(a)}{b-a}=0.
$$
Odsud dostáváme
$$
    f^\prime(c)=\frac{f(b)-f(a)}{b-a},
$$
což jsme chtěli dokázat.
\end{proof}

\begin{veta}[Cauchy] \label{veta5.8}
Nechť $a,b \in \R, a<b$, $f$ a $g$ jsou funkce spojité na $[a,b]$, $f$ má v každém bodě $x\in(a,b)$ derivaci a $g$ má v každém bodě $x\in(a,b)$ vlastní nenulovou derivaci. Potom $g(a)\ne g(b)$ a existuje $c\in(a,b)$ takové, že
\[
 \frac{f'(c)}{g'(c)}=\frac{f(b)-f(a)}{g(b)-g(a)}.
\]
\end{veta}

\begin{proof}
	Podle Lagrangeovy věty \ref{veta5.7} nalezneme $d \in (a,b)$ takový, že
	$$
	    g^\prime(d)=\frac{g(b)-g(a)}{b-a}.
	$$
	Platí $g^\prime(d)\ne 0$, a tedy $g(a)\ne g(b)$. Položme
	$$
	\varphi(x) = (f(x)-f(a))(g(b)-g(a))-(f(b)-f(a))(g(x)-g(a)), \quad x \in \left [ a,b \right ].
	$$
	Platí, že
	\begin{itemize}
		\item $\varphi$ je spojitá na $\left [ a,b \right ] $,
		\item $\varphi(a)=0$, $\varphi(b)=0$,
		\item $\varphi^\prime(x)=\underbrace{f^\prime(x)}_{\in \mathbb R ^*}\underbrace{(g(b)-g(a))}_{\ne 0}-(f(b)-f(a))\underbrace{g^\prime(x)}_{\in \mathbb R}$, $x \in (a,b)$.
	\end{itemize}
	Podle Rolleovy věty \ref{veta5.6} nalezneme $c \in (a,b)$ takové, že $\varphi^\prime(c)=0$,
	neboli
	$$
	    \underbrace{f^\prime(c)}_{\in \mathbb R }\underbrace{(g(b)-g(a))}_{\ne 0}-(f(b)-f(a))\underbrace{g^\prime(c)}_{\ne 0}=0.
	$$
	Odtud plyne
	$$
 \frac{f^\prime(c)}{g^\prime(c)}=\frac{f(b)-f(a)}{g(b)-g(a)},
	$$
	což jsme chtěli dokázat.
\end{proof}

\begin{veta}[l'H\^ospitalovo pravidlo] \label{veta5.9}
Nechť $a\in\R\cup\{-\infty\}$, $f,g$ jsou funkce, existuje
$$\displaystyle \lim_{x\to a^+}\frac{f'(x)}{g'(x)}$$ a platí buď
\begin{enumerate}[(a)]
  \item $\lim_{x\to a^+} f(x)=\lim_{x\to a^+} g(x)=0$, nebo
  \item $\lim_{x\to a^+} \abs{g(x)}=\infty$.
\end{enumerate}
Potom platí $$ \lim_{x\to a^+}\frac{f(x)}{g(x)}=\lim_{x\to a^+}\frac{f'(x)}{g'(x)}.$$
\end{veta}

\begin{proof}
	Provedeme důkaz (a), důkaz (b) umět nemusíme.
	\par Označme $L=\lim_{x\to a_+} \frac{f^\prime(x)}{g^\prime(x)}.$
	\paragraph{Pomocné tvrzení.}
	\begin{enumerate}
		\item[(i)] Pro každé $\alpha \in \mathbb R ,\alpha > L$ existuje $\delta > 0$ takové, že
		$$
		    \forall x \in P_+(a,\delta):\frac{f(x)}{g(x)}<\alpha.
		$$
		\item[(ii)] Pro každé $\beta \in \mathbb R , \beta < L$ existuje $\delta^\prime>0$ takové,
		že
		$$
		    \forall x \in P_+(a,\delta^\prime):\frac{f(x)}{g(x)}>\beta.
		$$
	\end{enumerate}
		\paragraph{Důkaz pomocného tvrzení.}
		\begin{enumerate}
			\item[(i)] Nalezneme $r\in (L,\alpha)$. Nalezneme $\delta > 0$ takové, že platí
			\begin{itemize}
				\item $f$, $f^\prime$, $g$ a $g^\prime$ jsou definovány na $P_+(a,\delta)$,
				\item $f^\prime$, $g^\prime$ jsou vlastní na $P_+(a,\delta)$,
				\item $g^\prime$ je nenulové na $P_+(a,\delta)$,
				\item $\forall x \in P_+(a,\delta):\frac{f^\prime(x)}{g^\prime(x)}<r$,
				\item $g$ je nenulová na $P_+(a,\delta)$; pokud by existoval nulový bod $f$ v
				$P_+(a,\delta)$, pak je nejvýše jeden a $\delta$ vhodně zmenšíme.
			\end{itemize}
		Zvolme $x\in P_+(a,\delta)$. Předpokládejme, že $y \in P_+(a,\delta)$ a $y<x$. Podle Cauchyovy
		věty \ref{veta5.8} nalezneme $c \in (y,x)$ takové, že
		$$
  \frac{f^\prime(c)}{g^\prime(c)}=\frac{f(x)-f(y)}{g(x)-g(y)}.
		$$
		Potom $\frac{f^\prime(c)}{g^\prime(c)}<r$, a tedy
		$$
  \frac{f(x)-f(y)}{g(x)-g(y)}<r.
		$$
		Potom máme, že
		$$
  \lim_{y\to a^+} \frac{f(x)-f(y)}{\underbrace{g(x)-g(y)}_{<r}}=\underbrace{\frac{f(x)}{\underbrace{g(x)}_{\ne 0}}}_{\leq r}<\alpha.
		$$
		\item[(ii)] Dokážeme analogicky.
	\end{enumerate}
	\paragraph{Vlastní důkaz.} Pokud $L= \infty$, resp. $L=-\infty$, pak věta plyne z (ii), resp. z (i).
	Pokud $L\in \mathbb R $ a $\varepsilon > 0$, položíme $\alpha = L+\varepsilon$, $\beta = L-\varepsilon$
	a nalezneme příslušná $\delta$ a $\delta^\prime$. Položme $\delta_1=\min \left \{ \delta, \delta^\prime \right \} $.
	Pak pro každé $x \in P_+(a,\delta_1)$ platí
	$$
	    L-\varepsilon =\beta < \frac{f(x)}{g(x)}<\alpha = L+\varepsilon, \quad \textrm{tj. } \frac{f(x)}{g(x)}\in B(L,\varepsilon),
	$$
	což jsme chtěli dokázat.
\end{proof}

\begin{poznamka}\,
\begin{enumerate}[$i.$]
	\item $\lim_{x\to 0} \frac{2x+1}{3x+1} \ne \lim_{x\to 0} \frac{2}{3}$
	\item $\lim_{x\to \infty} \frac{x}{x+\sin x}=1$, ale l'H\^ospitalovým pravidlem výsledek nedostaneme
	\item $\lim_{x\to 0} \frac{e^{-\frac{1}{x^2}}}{x}= \lim_{x\to 0} \frac{e^{-\frac{1}{x^2}}\cdot \frac{2}{x^3}}{1}=\lim_{x\to 0} \frac{2e^{-\frac{1}{x^2}}}{x^3}$
\end{enumerate}
\end{poznamka}

\begin{priklad}
	Spočtěte $\lim_{x\to 0} \frac{\sin x -x}{x^3}$.
\end{priklad}

\begin{reseni}
	L'H\^ospitalovým pravidlem dostáváme
	$$
	    \lim_{x\to 0} \frac{\sin x -x}{x^3} = \lim_{x\to 0} \frac{\cos x -1}{3x^2} = \frac{1}{3}\cdot \frac{-1}{2} = -\frac{1}{6}.
	$$
\end{reseni}

\begin{veta}[Vztah derivace a monotonie]\label{veta5.10}
Nechť $f$ je spojitá na intervalu $J$ a v každém vnitřním bodě $J$ (množinu všech vnitřních bodů intervalu $J$ označme jako $\operatorname{int} J$) existuje derivace $f$.
\begin{enumerate}[(i)]
  \item Je-li $f'(x)>0$ pro všechna $x\in\operatorname{int} J$, pak $f$ je rostoucí na $J$.
  \item Je-li $f'(x)<0$ pro všechna $x\in\operatorname{int} J$, pak $f$ je klesající na $J$.
  \item Je-li $f'(x)\ge 0$ pro všechna $x\in\operatorname{int} J$, pak $f$ je neklesající na $J$.
  \item Je-li $f'(x)\le 0$ pro všechna $x\in\operatorname{int} J$, pak $f$ je nerostoucí na $J$.
\end{enumerate}
\end{veta}

\begin{proof}
Dokážeme (i), ostatní případy se dokážou analogicky.
	Předpokládejme, že $x,y\in J$, $x<y$. Podle Lagrangeovy věty \ref{veta5.7} nalezneme
	$c \in (x,y)$ takové, že
	$$
	    f^\prime(c) = \frac{f(y)-f(x)}{y-x}>0.
	$$
	Odtud plyne $f(x)<f(y)$.
\end{proof}

\begin{poznamka}[Varování!]
	Věta \ref{veta5.10} je formulovaná pro \emph{interval}.
\end{poznamka}

\begin{veta}\label{veta5.11}
Nechť $f$ je spojitá zprava v bodě $a\in\R$ a existuje $\displaystyle \lim_{x\to a^+} f'(x)$. Potom existuje $f'_+(a)$ a platí $f'_+(a)=\displaystyle \lim_{x\to a^+} f'(x)$.
\end{veta}

\begin{proof}
    Je vlastně
	$$
	f_+^\prime(a)=\lim_{x\to a_+} \frac{f(x)-f(a)}{x-a}.
	$$
	Platí, že
	$$
 \lim_{x\to a_+} (f(x)-f(a))=0, \qquad \lim_{x\to a_+} (x-a)=0.
	$$
	Odtud z l'H\^ospitalova pravidla dostáváme
	$$
	    f_+^\prime(a) = \lim_{x\to a_+} \frac{f^\prime(x)}{1}=\lim_{x\to a_+} f^\prime(x).
	$$
\end{proof}

\begin{definice}
Nechť $n\in\N$, $a\in\R$ a $f$ má vlastní $n$-tou derivaci na okolí bodu $a$. Potom \textbf{$(n+1)$-ní derivací funkce} $f$ v bodě $a$ rozumíme
\[
 f^{(n+1)}(a)=\lim_{h\to 0} \frac{f^{(n)}(a+h)-f^{(n)}(a)}{h}.
\]
\end{definice}

\begin{poznamka}
	Derivace značíme
	$$
	    f^\prime \quad f^{\prime\prime}\quad f^{\prime\prime\prime}\quad f^{(4)}\quad f^{(5)}\quad\dots\quad f^{(n)}.
	$$
\end{poznamka}

\section{Taylorův polynom}
\begin{motivace}
	Máme-li funkci $f$, která má v bodě $a$ vlastní derivaci, pak je tečna
	$$
	    t(x) = f(a) + f^\prime(a)(x-a).
	$$
	Podíváme-li se na limitu
	$$
 \lim_{x \to a} \frac{f(x)-t(x)}{x-a} = 0 = \lim_{x \to a} \frac{f(x)-f(a)-f^\prime(a)(x-a)}{x-a},
	$$
	zjistíme, že aproximuje chování funkce.
\end{motivace}

\subsection{Základní vlastnosti}
\begin{definice}
Nechť $f$ je funkce, $a \in \mathbb{R}$ a $f^{(n)}(a) \in \mathbb{R}$. Pak polynom
\[
T^{f,a}_n(x) = f(a) + f'(a)(x - a) + \dots + \frac{1}{n!} f^{(n)}(a)(x - a)^n
\]
nazýváme \textbf{Taylorovým polynomem} řádu $n$ funkce $f$ v bodě $a$.
\end{definice}

\begin{poznamka}\,
\begin{enumerate}[(1)]
	\item $\operatorname{st} \left( T_n^{f,a} \right ) \leq n $
	\item $T:=T_n^{f,a}$, pak $T(a) = f(a)$,
	$$T^\prime(a) = f^\prime(a)+ f^{\prime\prime}(a)(x-a)+\dots+ \frac{1}{(n-1)!}f^{(n)}(a)(x-a)^{n-1}=T_{n-1}^{f^\prime,a}(a)=f^\prime(a)$$
	$$
	    T^{\prime\prime}(a) = f^{\prime\prime}(a)\quad \dots\quad T^{(n)}(a)=f^{(n)}(a)
	$$
\end{enumerate}
\end{poznamka}

\section{Číselné řady}
\subsection{Základní pojmy}
\begin{definice}
Nechť $\{a_n\}_{n=1}^\infty$ je posloupnost.
\begin{enumerate}[(a)]
\item Pro $m\in\mathbb{N}$ položme $s_m=\sum_{n=1}^m a_n$. \textbf{Součet řady} $\sum_{n=1}^\infty a_n$ je limita posloupnosti $\{s_m\}$, pokud tato limita existuje.
\item Řekneme, že řada $\sum_{n=1}^\infty a_n$ \textbf{konverguje} (\textbf{je konvergentní}), je-li $\lim_{m\to \infty} s_m$ vlastní.
Řekneme, že řada $\sum_{n=1}^\infty a_n$ \textbf{diverguje} (\textbf{je divergentní}), jestliže $\lim_{m\to \infty} s_m$ neexistuje nebo je nevlastní.
Pro jemnější rozlišení budeme někdy říkat, že řada $\sum_{n=1}^\infty a_n$ \textbf{diverguje k $\infty$}, respektive
\textbf{diverguje k $-\infty$}, jestliže $\lim s_m = \infty$, resp. $\lim s_m = -\infty$.
\item Číslo $a_n, n \in \mathbb N $ je \textbf{$n$-tým členem} řady $\sum_{n=1}a_n$
a číslo $s_m, m \in \mathbb N $ je jejím \textbf{$m$-tým částečným součtem}.
\end{enumerate}
\end{definice}

\begin{poznamka}
	Nekonečná řada
	$$
	\sum_{n=1}^\infty a_n \begin{cases}
	\textrm{má součet } \begin{cases}
		\textrm{konverguje}\\
		\textrm{diverguje } \begin{cases}
			\textrm{k } \infty,\\
			\textrm{k } -\infty
\end{cases}
\end{cases}	\\
\textrm{nemá součet}.
\end{cases}
	$$
\end{poznamka}

\begin{priklad}\,
\begin{enumerate}[(1)]
	\item $\sum_{n=1}^\infty (-1)^n$ nemá součet (diverguje)
	\item $\sum_{n=1}^\infty q^{n-1}$ konverguje $\iff q \in (-1,1)$, je totiž
	$$
	    s_n = 1+q+q^2+\dots + q^{n-1}= \begin{cases}
					\frac{q^n-1}{q-1}, & q \in \mathbb R  \smallsetminus \left \{ 1 \right \}\\
				 n, & q=1.
					\end{cases}
	$$
	\item $\sum_{n=1}^\infty \frac{1}{n(n+1)}=1$, neboť
	\begin{align*}
	\frac{1}{n(n+1)} &= \frac{1}{n}-\frac{1}{n+1},\\
	s_m &= \left (\frac{1}{1} - \frac{1}{2}\right) + \left( \frac{1}{2} - \frac{1}{3} \right ) + \dots + \left( \frac{1}{m} - \frac{1}{m+1}\right ) = 1 - \frac{1}{m+1} .
\end{align*}
\end{enumerate}
\end{priklad}

\begin{veta}[Nutná podmínka konvergence řady] \label{veta7.1}
Nechť řada $\sum_{n=1}^\infty a_n$ konverguje. Pak platí $\lim a_n=0$.
\end{veta}

\begin{proof}
	Označme $x_n = \sum_{j=1}^n a_j$. Počítáme
	$$
 \lim a_n = \lim (s_n - s_{n-1}).
	$$
	Protože $\sum_{n=1}^\infty a_n$ konverguje, pak $\lim_{n \to \infty} s_n = s \in \mathbb R $.
	Odtud máme
	$$
	\lim a_n = \lim (s_n - s_{n-1}) = s - s = 0.
	$$
\end{proof}

\begin{veta}[Bolzanova--Cauchyova podmínka konvergence řady] \label{7.2}
Nechť $\sum_{n=1}^\infty a_n$ je řada. Pak jsou následující tvrzení ekvivalentní:
\begin{enumerate}[(i)]
	\item $\sum_{n=1}^\infty a_n$ je konvergentní,
	\item platí výrok
	\[
\forall \varepsilon \in \mathbb R , \varepsilon>0 \, \exists n_0\in \mathbb N  \, \forall m, n\in \mathbb N , m\ge n\ge n_0:\left| \sum_{k=n}^m a_k \right| < \varepsilon.
\]
\end{enumerate}
\end{veta}

\begin{proof}
	Ať je $s_n = \sum_{j = 1}^n a_j$.
	\begin{itemize}
		\item[$\impliedby$] Platí
		\begin{equation}\label{bcpodminka}
		\forall \varepsilon > 0 \exists n_0 \in \mathbb N  \forall m, n \in \mathbb N , m \geq n \geq n_0: s_m - s_{n-1}| < \varepsilon.
\end{equation}
		Ověříme BC-podmínku pro posloupnost $\left \{ s_m \right \} $. Zvolme $\varepsilon > 0$. Podle
		(\ref{bcpodminka}) nalezneme k $\varepsilon$ příslušné $n_0\in \mathbb N $. Zvolme
		libovolně $m_1, n_1 \in \mathbb N , m_1 \geq n_0, n_1 \geq n_0$. Rozlišíme následující možnosti:
		\begin{enumerate}[(i)]
			\item $m_1 > n_1$: $|s_{m_1} - s_{n_1}|< \varepsilon $, neboť $m_1 \geq n_1 + 1 \geq n_0$,
			\item $m_1 = n_1$: $|s_{m_1} - s_{n_1}| = 0 < \varepsilon$,
			\item $m_1 < n_1$: $|s_{m_1} - s_{n_1}| = < |s_{n_1} - s_{m_1}|\varepsilon $, neboť $n_1 \geq m_1 + 1 \geq n_0$.
		\end{enumerate}
		Máme tedy že $\left \{ a_n \right \} $ splňuje BC-podmínku, tedy $\left \{ s_n \right \} $ je konvergentní,
		a proto $\sum_{n=1}\infty a_n$ konverguje.
		\item[$\implies$] Protože $\sum a_n$ konverguje, splňuje $\left \{ s_n \right \} $ BC-podmínku, odkud plyne (ii).\qedhere
	\end{itemize}
\end{proof}

\begin{priklad}
    Řada $\sum_{n=1}^\infty \frac{1}{n}$ (harmonická řada) je divergentní.
    \par Vezměme
    $$
    s_{2n}-s_n=\underbrace{\frac{1}{n+1} + \frac{1}{n+2}+ \dots + \frac{1}{n+n}}_{n \textrm{ členů}} \geq n \cdot \frac{1}{2n} = \frac{1}{2}.
    $$
    Protože $\sum_{n=1}^\infty \frac{1}{n}$ nesplňuje BC-podmínku, a tedy diverguje (k nekonečnu).
    Zároveň je to protipříklad na opačnou implikaci ve větě \ref{veta7.1}.
\end{priklad}

\begin{poznamka}
	Platí $\lim_{n \to \infty} \left( \sum_{j=1}^\infty \frac{1}{j} - \log n \right ) = \gamma \in \mathbb R  $.
	Konstantě $\gamma$ říkáme Eulerova-Mascheroniho konstanta.
\end{poznamka}

\begin{veta}\label{veta7.3}
Nechť řady $\sum_{n=1}^\infty a_n$ a $\sum_{n=1}^\infty b_n$ mají součet.
\begin{enumerate}[(a)]
	\item Nechť $\alpha \in \mathbb R $ a výraz $\alpha\sum_{n=1}^\infty a_n$ je definován.
	Potom má řada $\sum_{n=1}^\infty \alpha a_n$ součet a platí
	$$
	    \sum_{n=1}^\infty \alpha a_n = \alpha\sum_{n=1}^\infty a_n.
	$$
	\item Nechť je výraz $\sum_{n=1}^\infty a_n+\sum_{n=1}^\infty b_n$ definován.
	Potom má řada $\sum_{n=1}^\infty (a_n+b_n)$ součet a platí
	$$
	    \sum_{n=1}^\infty (a_n+b_n) = \sum_{n=1}^\infty a_n +\sum_{n=1}^\infty b_n.
	$$
\end{enumerate}
\end{veta}

\begin{proof}
	\,
	\begin{enumerate}[(a)]
		\item Označme $s_n = \sum_{j=1}^n a_j, \sigma_n= \sum_{j=1}^n \alpha a_j$. Platí $\sigma_n = \alpha s_n$.
		Platí $\lim s_n = s \in \mathbb R ^*$. Podle předpokladu je výraz definován, takže
		$\lim \sigma _n = \lim \alpha s_n$, tedy $\sum_{j=1}^\infty \alpha a_j = \alpha \sum_{j=1}^\infty  a_j$.
		\item Označme $s_n = \sum_{j=1}^n a_j, \sigma_n= \sum_{j=1}^n b_j$. Pak je
		$t_n = \sum_{j=1}^n \left( s_n + \sigma s_n \right ) $. Pak $\lim t_n = \lim \left( s_n + \sigma_n \right )  = s + \sigma$,
		neboť tento výraz je definován.\qedhere
	\end{enumerate}
\end{proof}

\subsection{Řady s nezápornými členy}
\begin{veta}\label{veta7.4}
Nechť je $\sum_{n=1}^\infty a_n$ řda s nezápornými členy. Pak má $\sum_{n=1}^\infty a_n$ součet.
\end{veta}

\begin{proof}
	Poněvadž pro každé $n \in \mathbb N $ platí $a_n \geq 0$, je posloupnost $\left \{ s_j \right \}_{j=1}^\infty $
	neklesající. Odtud máme podle věty \ref{v213}, že $\lim_{j\to \infty} s_j \in \mathbb R ^*$.
\end{proof}

\begin{veta}[Srovnávací kritérium] \label{veta7.5}
Nechť $\sum_{n=1}^{\infty} a_n$ je řada s nezápornými členy, $\sum_{n=1}^{\infty} b_n$ je řada s kladnými členy a
existuje $n_0 \in \mathbb{N}$ takové, že pro každé $n \in \mathbb{N}$, $n \ge n_0$, platí
$a_n \le b_n$.
\begin{enumerate}[(a)]
\item Je-li $\sum_{n=1}^{\infty} b_n$ konvergentní, pak je i
$\sum_{n=1}^{\infty} a_n$ konvergentní.
\item Je-li $\sum_{n=1}^{\infty} a_n$ divergentní, pak je i
$\sum_{n=1}^{\infty} b_n$ divergentní.
\end{enumerate}
\end{veta}

\begin{proof}\,
\begin{enumerate}[(a)]
	\item Máme
	\begin{align*}
	s_n &= a_1 + a_2 + \dots + a_{n_0} + a_{n_0+1} + \dots + a_n, \quad n > n_0\\
	t_n &= b_1 + b_2 + \dots + b_{n_0} + b_{n_0+1} + \dots + b_n.
\end{align*}
    Platí
    $$
        s_n \leq s_{n_0-1} + t_n - t_{n_0-1} = t_n + s_{n_0-1} - t_{n_0-1}.
    $$
    Dále platí $\lim s_n = s \in \mathbb R ^*$ a $\lim t_n = t \in \mathbb R ^*$ podle věty \ref{veta7.4}.
    Potom $s \leq t + s_{n_0-1} - t_{n_0-1}$. Podle předpokladu je $t \in \mathbb R $, takže
    celá pravá strana je prvekm $\mathbb R $, a tedy $s \in \mathbb R $, neboť všechny členy jsou nezáporné.
    \item Plyne z (a) obměnou, což jsme již dokázali. \qedhere
\end{enumerate}
\end{proof}

\begin{veta}[Limitní srovnávací kritérium] \label{veta7.6}
Nechť $\sum_{n=1}^{\infty} a_n$ je řada s nezápornými členy,
$\sum_{n=1}^{\infty} b_n$ je řada s kladnými členy a
$\lim \frac{a_n}{b_n}$ existuje. Označme
$
A = \lim \frac{a_n}{b_n}.
$
\begin{enumerate}[(a)]
	\item Nechť $A \in (0,\infty)$. Pak $\sum_{n=1}^{\infty} a_n$ konverguje právě tehdy,
když konverguje $\sum_{n=1}^{\infty} b_n$.
\item Nechť $A = 0$. Pokud $\sum_{n=1}^{\infty} b_n$ konverguje,
pak konverguje i $\sum_{n=1}^{\infty} a_n$.
\item Nechť $A = \infty$. Pokud $\sum_{n=1}^{\infty} a_n$ konverguje,
pak konverguje i $\sum_{n=1}^{\infty} b_n$.
\end{enumerate}
\end{veta}

\begin{proof}
	\,
	\begin{enumerate}[(a)]
		\item
		\begin{itemize}
			\item[$\implies$]
			Nalezneme $n_0 \in \mathbb N $ takové, že
			$$\forall n \in \mathbb N , n \geq n_0 : \frac{a_n}{b_n} > \frac{1}{2}A,$$
			tedy že
			$$
			    \forall n \in \mathbb N , n \geq n_0 : a_n > \frac{1}{2}A b_n.
			$$
			Věta \ref{veta7.5} dává, že řada
			$$
			    \sum_{n=1}^\infty \left( \frac{1}{2}Ab_n \right )
			$$
			konverguje. Dále podle věty \ref{veta7.3}(a) máme, že
			$$
			    \sum_{n=1}^\infty b_n
			$$
			konverguje.
			\item[$\impliedby$] Nalezneme $n_0 \in \mathbb N $ takové, že
			$$
			    \forall n  \geq n_0: \frac{a_n}{b_n} < 2A.
			$$
			Potom máme
			$$
			    \forall n \geq n_0: a_n < 2Ab_n.
			$$
			Protože $\sum b_n$ konverguje, konverguje i $\sum 2Ab_n$
			podle věty \ref{veta7.3}. Odtud podle věty \ref{veta7.5} dostáváme,
			že konverguje i $\sum a_n$.
		\end{itemize}
		\item Nalezneme $n_0 \in \mathbb N $ takové, že
		$$
		    \forall n \geq n_0: \frac{a_n}{b_n} < 1.
		$$
		Odtud máme
		$$
		    \forall n \geq n_0: a_n < b_n.
		$$
		Protože $\sum b_n$ konverguje, konverguje podle věty \ref{veta7.5} i $\sum a_n$.
		\item Obdobně. \qedhere
	\end{enumerate}
\end{proof}

\begin{priklad}\,
\begin{enumerate}[(1)]
	\item $\sum a_n=\sum_{n=1}^\infty \frac{n+ \sqrt{n}}{n^2+3n}$
	\par Zvolme $b_n = \frac{1}{n}$. Z limitního srovnávacího kritéria dostaneme
	$$
 \lim_{n \to \infty} \frac{a_n}{b_n} = \lim_{n \to \infty} \frac{n^2+n \sqrt{n}}{n^2+3n} = 1 \in (0,\infty),
	$$
	a tedy $\sum a_n$ diverguje, neboť $\sum \frac{1}{n}$ diverguje.
	\item $\sum a_n=\sum_{n=1}^\infty \frac{1}{n^2}$
	\par Zvolme $b_n = \frac{1}{n(n+1)}$, přičemž $\sum b_n = 1$. Pak máme
	$$
 \lim_{n \to \infty} \frac{n(n+1)}{n^2} = 1 \in (0,\infty),
	$$
	takže $\sum a_n$ konverguje, neboť $\sum b_n$ konverguje.
\end{enumerate}
\end{priklad}

\begin{veta}[Cauchyovo odmocninové kritérium] \label{veta7.7}
Nechť $\sum_{n=1}^{\infty} a_n$ je řada s nezápornými členy.
\begin{enumerate}[(a)]
	\item Jestliže platí
\[
\exists q \in (0,1)\ \exists n_0 \in \mathbb{N}\ \forall n \in \mathbb{N}, n \ge n_0 :
\sqrt[n]{a_n} \le q,
\]
pak je $\sum_{n=1}^{\infty} a_n$ konvergentní.

\item  Jestliže $\limsup \sqrt[n]{a_n} < 1$, pak je
$\sum_{n=1}^{\infty} a_n$ konvergentní.

\item  Jestliže $\lim \sqrt[n]{a_n} < 1$, pak je
$\sum_{n=1}^{\infty} a_n$ konvergentní.

\item  Jestliže $\limsup \sqrt[n]{a_n} > 1$, pak není pravda, že $\lim a_n = 0$,
a tedy $\sum_{n=1}^{\infty} a_n$ je divergentní.

\item Jestliže $\lim \sqrt[n]{a_n} > 1$, pak není pravda, že $\lim a_n = 0$,
a tedy $\sum_{n=1}^{\infty} a_n$ je divergentní.
\end{enumerate}
\end{veta}

\begin{proof}\,
\begin{enumerate}[(a)]
	\item Platí
	$$
	    \forall n \geq n_0: 0 \leq a_n \leq \underbrace{q^n}_{b_n}.
	$$
	Protože je $\sum b_n$ konvergentní, podle věty \ref{veta7.5} dostáváme, že
	$\sum a_n$ je konvergentní.
	\item Nalezneme $q\in (\limsup \sqrt[n]{a_n}, 1)$. Potom nalezneme
	$n_0\in \mathbb N $ takové, že
	$$
	    \sup \left \{ \sqrt[n]{a_n}, n \geq n_0 \right \} < q.
	$$
	Odsud plyne
	$$
	    \forall n \in \mathbb N , n \geq n_0: \sqrt[n]{a_n} < q.
	$$
	Z (a) plyne, že $\sum a_n$ konverguje.
	\item $\limsup a_n = \lim a_n$, a tedy můžeme použít (b).
	\item Nalezneme rostoucí posloupnost přirozených čísel $\left \{ n_k \right \}_{k=1}^\infty $ takovou, že
	$$
 \lim_{k \to \infty} \sqrt[n]{a_{n_k}} = \limsup_{n\to \infty} \sqrt[n]{a_n} > 1.
	$$
	Nalezneme $k_0 \in \mathbb N $ takové, že
	$$
	    \forall k \geq k_0: \sqrt[n]{a_{n_k}} > 1.
	$$
	Zároveň tedy platí
	$$
	    \forall k \geq k_0: a_{n_k} > 1.
	$$
	Neplatí tedy, že $\lim a_{n_k} =0$, a tedy neplatí, že $\lim a_n = 0$. Proto
	podle věty \ref{veta7.1} $\sum a_n$ diverguje.
	\item Je $\limsup \sqrt[n]{a_n} = \lim \sqrt[n]{a_n} > 1$. Dále použitím (d).\qedhere
\end{enumerate}
\end{proof}

\begin{poznamka}
	Pokud $\limsup \sqrt[n]{a_n} = 1$, pak nelze z této informace odvodit ani konvergenci, ani divergenci.
	Například řada $\sum \frac{1}{n}$ diverguje, přičemž $\lim \sqrt[n]{\frac{1}{n}} = 1$. Naopak
	řada $\sum \frac{1}{n^2}$ konverguje, přičemž $\lim \sqrt[n]{\frac{1}{n^2}}= 1$.
\end{poznamka}

\begin{veta}[D'Alembertovo podílové kritérium]\label{veta7.8}
Nechť $\sum_{n=1}^{\infty} a_n$ je řada s kladnými členy.
\begin{enumerate}[(a)]
	\item Jestliže platí
\[
\exists q \in (0,1)\ \exists n_0 \in \mathbb{N}\ \forall n \in \mathbb{N}, n \ge n_0 :
\frac{a_{n+1}}{a_n} \le q,
\]
pak je řada $\sum_{n=1}^{\infty} a_n$ konvergentní.

\item  Jestliže $\limsup \frac{a_{n+1}}{a_n} < 1$, pak je řada
$\sum_{n=1}^{\infty} a_n$ konvergentní.
\item  Jestliže $\lim \frac{a_{n+1}}{a_n} < 1$, pak je řada
$\sum_{n=1}^{\infty} a_n$ konvergentní.

\item  Jestliže $\lim \frac{a_{n+1}}{a_n} > 1$, pak neplatí $\lim a_n = 0$,
a tedy je řada $\sum_{n=1}^{\infty} a_n$ divergentní.
\end{enumerate}
\end{veta}

\begin{proof}
\begin{enumerate}[(a)]
	\item Pro každé $n \in \mathbb N , n \geq n_0$ platí
	$$
	    a_n \leq a_{n_0}q^{n-n_0}.
	$$
	Použijeme matematickou indukci.
	\begin{itemize}
		\item $n = n_0$: Platí.
		\item $n \rightsquigarrow n+ 1$: Je
		$$a_{n+1} \leq qa_n \leq qa_{n_0}q^{n-n_0}=a_{n_0}q^{(n+1)-n_0}.$$
		Protože $\sum a_{n_0}a^{n-n_0}=\sum_{n=1}^\infty a_{n_0}q^{-n_0+1}q^{n-1}$ konverguje,
		konverguje i $\sum a_n$ podle věty \ref{veta7.5}.
	\end{itemize}
	\item Nalezneme $q \in \left( \limsup \frac{a_{n+1}}{a_n},1 \right ) $ a $n_0 \in \mathbb N $
	taková, že
	$$
	    \sup \left \{ \frac{a_{n+1}}{a_n}, n \geq n_0 \right \} < q \in (0,1).
	$$
	Potom
	$$
	    \forall n \geq n_0: \frac{a_{n+1}}{a_n} < q.
	$$
	Z již dokázané části (a) pak vyplývá, že $\sum a_n$ je konvergentní.
	\item Je $\limsup \frac{a_{n+1}}{a_n} = \lim \frac{a_{n+1}}{a_n} < 1$, dále podle (b).
	\item Nalezneme $n_0 \in \mathbb N $ takové, že
	$$
	    \forall n \geq n_0 : \frac{a_{n+1}}{a_n} > 1,
	$$
	neboli
	$$
	    \forall n \geq n_0: a_{n+1} > a_n.
	$$
	Odsud plyne, že
	$$
	    \forall n \geq n_0: a_n \geq a_{n_0} > 0.
	$$
	Tedy neplatí, že $\lim a_n = 0$, a proto $\sum a_n$ není konvergentní podle věty \ref{veta7.1}.\qedhere
\end{enumerate}
\end{proof}

\begin{veta}[Kondenzační kritérium] \label{veta7.9}
Nechť $\{a_n\}$ je nerostoucí posloupnost s nezápornými členy.
Pak řada
$
\sum_{n=1}^{\infty} a_n
$
konverguje právě tehdy, když konverguje řada
$
\sum_{n=0}^{\infty} 2^n a_{2^n}.
$
\end{veta}

\begin{proof}
Označme
\begin{align*}
s_n & = a_1+ a_2+ \dots a_n,\\
t_k &= a_1 + 2a_2 + \dots + 2^ka_{2^k}, \quad k \in \mathbb N _0.
\end{align*}
	Pro $n \leq 2^k$ platí
	\begin{align}\label{prvni}
	\begin{split}
	s_n &\leq a_1 + (a_2 + a_3) + (a_4 + a_5 + a_6 + a_7) + \dots + (a_{2^k}+ \dots + a_{2^{k+1}-1})\\
	& \leq a_1 + (a_2 + a_2) + (a_4 + a_4 + a_4 + a_4) + \dots + (a_{2^k}+ \dots + a_{2^k})\\
	&= a_1 + 2a_2 + \dots + 2^ka_{2^k} = t_k.
\end{split}
	\end{align}
	Pokud $n \geq 2^k$, platí
	\begin{align}\label{druha}
	\begin{split}
	s_n& \geq a_1 + a_2 + (a_3 + a_4) + (a_5+ a_6+a_7+a_8) + \dots + (a_{2^{k-1}+1}+a_{2^k})\\
	& \leq \frac{1}{2} a_1+ a_2 + 2 a_4+ 4a_8 + \dots + 2^{k-1} a_{2^k} \\
	&= \frac{1}{2} \cdot \left( a_1 + 2a_2 + 4a_4+8a_8+ \dots + 2^ka_{2^k} \right ) = \frac{1}{2} \cdot t_k.
\end{split}
\end{align}

\begin{enumerate}
	\item[$\implies$] Máme, že $\lim s_n = s\in \mathbb R $, neboť $\sum a_n$ konverguje. Pro každé
	$n, k$ splňující $n\geq 2^k$ platí podle (\ref{druha}) $t_k \leq 2s_n \leq 2s$. Odsud máme, že
	$\lim t_k \in \mathbb R$, neboť $\sum 2^k a_{2^k}$ je konvergentní
	\item[$\impliedby$] Podle předpokladu $\lim_{k\to \infty}t_k = t\in \mathbb R $.
	Zvolme $n\in \mathbb N $. Potom pro každé $k\in \mathbb N, n \leq 2^k $ platí (\ref{druha}),
	tedy $s_n \leq t_k \leq t$. Odtud plyne $\lim s_n \leq t$. Proto je $\sum a_n$ konvergentní. \qedhere
\end{enumerate}
\end{proof}

\begin{veta}\label{veta7.10}

\end{veta}

\begin{proof}
	Rozlišíme dva případy.
	\begin{enumerate}[(i)]
		\item $\alpha > 1$: Použijeme větu \ref{veta7.9}. Místo řady $\sum_{n=1}^\infty \frac{1}{n^\alpha}$
		vyšetřujeme řadu
		$$\sum_{k=0}^\infty 2^k \frac{1}{{2^k}^\alpha} = \sum_{k=0}^\infty \left( 2^k \right )^{1-\alpha} = \sum_{k=0}^\infty \left( 2^{1-\alpha} \right )^k , $$
		která konverguje, neboť $2^{1-\alpha < 0} \in (0,1)$; proto konverguje i původní řada.
		\item $\alpha \leq 1$: Platí
		$$
		    \forall n \in \mathbb N : \frac{1}{n^\alpha} \geq \frac{1}{n} \geq 0.
		$$
		Řada $\sum \frac{1}{n}$ diverguje, a tedy máme podle srovnávacího kritéria \ref{veta7.5},
		že i řada $\sum \frac{1}{n^\alpha}$ diverguje.\qedhere
\end{enumerate}
\end{proof}

\subsection{Řady s obecnými členy}
\begin{veta}[Leibniz] \label{veta7.11}

\end{veta}

\begin{proof}
	BÚNO $\left \{ a_n \right \} $ je nerostoucí. Potom
	$$
	    \lim a_n = \inf \left \{ a_n, n \in \mathbb N  \right \}.
	$$
	Z předpokladu víme, že $\lim a_n = 0$, a tedy $\forall n \in \mathbb N : a_n \geq 0$.
	Pak je
	$$
	    s_n = -a_1 + a_2 - a_3 +a_4 + \dots + (-1)^na_n.
	$$
	Pro každé $n\in \mathbb N $ platí
	\begin{align*}
	s_{2n} &= \underbrace{(a_1+a_2)}_{\leq 0}+\underbrace{(-a_3+a_4)}_{\leq 0}+\dots +\underbrace{(-a_{2n-1}+a_{2n})}_{\leq 0},\\
    s_{2n+2} &= s_{2n} + \underbrace{(-a_{2n+1} + a_{2n+2})}_{\leq 0} \leq s_{2n}.
    \end{align*}
	Dále platí
	$$
	    s_{2n} = -a_1 + \underbrace{(a_2 - a_3)}_{\leq 0} + \underbrace{(a_4  -a_5)}_{\leq 0}+ \dots + \underbrace{(a_{2n-2}+a_{2n-1})}_{\leq 0} + a_{2n}\geq -a_1,
	$$
	odkud máme, že $\lim s_{2n}=s \in \mathbb R $.
	\par Pro každé $n \in \mathbb N $ platí
	$$
	    s_{2n+1} = s_{2n}-a_{2n+1}.
	$$
	Potom
	$$
	    \lim s_{2n+1} = \lim(s_{2n}-a_{2n+1}) = s - 0 = s.
	$$
	Chceme dokázat, že $\lim s_n = s$. Zvolme $\varepsilon > 0$. K němu nalezneme
	$k_1, k_2 \in \mathbb N $ taková, že
	$$
	    \forall k \in \mathbb N , k \geq k_1: |s_{2k} - s| < \varepsilon.
	$$
	Podobně
	$$
	    \forall k \in \mathbb N , k \geq k_2: |s_{2k+1}-s| < \varepsilon.
	$$
	Položme $n_0 = \max \left \{ 2k_1, 2k_2 +1 \right \} $. Vezměme $n\geq n_0$. Pokud je $n$ sudé,
	je $n = 2k$ a $k\geq k_1$. Pak platí
	$$
	    |s_n - s| = |s_{2k}-s| < \varepsilon.
	$$
	Pokud je $n$ liché, je $n = 2k+1$, $k\geq k_2$. Pak platí
	$$
	    |s_n-s| = |s_{2k+1}-s| < \varepsilon.
	$$
\end{proof}

\begin{priklad}
Řada $\sum_{n=1}^\infty (-1)^n \frac{1}{n}$ konverguje.
\end{priklad}

\subsection{Absolutní konvergence řad}
\begin{definice}

\end{definice}

\begin{veta}\label{veta7.12}

\end{veta}

\begin{proof}
	Použijeme větu \ref{7.2}, tj. ověříme BC-podmínku pro řadu $\sum_{n=1}^\infty a_n$.
	\par Zvolme $\varepsilon > 0$. Protože $\sum |a_n|$ konverguje, platí
	$$
	    \exists n_0 \, \forall m \geq n \geq n_0: \left | \sum_{j=n}^m |a_j| \right | < \varepsilon.
	$$
	Pro každé $m,n \in \mathbb N , m \geq n \geq n_0$ platí
	$$
	\left | \sum_{j=n}^m a_j  \right | \leq \sum _{j=n}^m |a_j| < \varepsilon
	$$
	díky trojúhelníkové nerovnosti. Odtud máme podle věty \ref{7.2}, že $\sum a_n$ konverguje.
	\par Označme
	\begin{align*}
	s_n &= a_1 + a_2 + \dots + a_n,\\
		\sigma_n &= |a_1| +|a_2| + \dots + |a_n|.
\end{align*}
	Potom $\lim s_n = s \in \mathbb R $ a $\lim \sigma_n = \sigma \in \mathbb R $ a platí
	$$
	\forall n \in \mathbb N : s_n \leq \sigma _n,
	$$
	odkud máme, že $s \leq \sigma$.
\end{proof}

\begin{poznamka}
Někdy je možné ukázat konvergenci řady $\sum_{n=1}^\infty a_n$ tako, že ukážeme konvergenci
řady $\sum_{n=1}^\infty |a_n|$, což je řada s nezápornými členy.
\end{poznamka}

\subsection{Přerovnání řad}
\begin{definice}

\end{definice}

\begin{veta}\label{veta7.13}

\end{veta}

\begin{proof}\,
\begin{enumerate}[(1)]
	\item Předpokládejme nejprve, že $\sum a_n$ má nezáporné členy.
	Pro každé $n$ platí
	$$
	    a_{\pi(1)}+a_{\pi(2)} + \dots + a_{\pi(n)} \leq a_1 + a_2 + +\dots a_m \leq \sum_{j=1}^\infty a_j,
	$$
	pokud $m \geq \max \left \{ \pi(1), \pi(2), \dots, \pi(n) \right \} $. Odtud plyne
	$$
	    \sum_{k=1}^\infty a_{\pi(k)} \leq \sum_{j=1}^\infty a_j.
	$$
	Protože je $\sum a_j$ konvergentní, je taky $\sum a_{\pi(k)}$ konvergentní. Poněvadž
	$\sum a_j$ je přerovnáním řady $\sum a_{\pi(k)}$, dostáváme podle předchozího
	$$
	\sum_{k=1}^\infty a_{\pi(k)} \geq \sum_{j=1}^\infty a_j,
	$$
	a tedy
	$$
	    \sum_{k=1}^\infty a_{\pi(k)} = \sum_{j=1}^\infty a_j
	$$
	\item Obecný případ. Pro $a \in \mathbb R $ budeme značit $a^+ = \max \left \{ a,0 \right \} $.
	Obdobně $a^- = \max \left \{ -a, 0 \right \} $. Platí $a=a^+-a^-$, $0\leq a^+\leq |a|$ a konečně
	$0\leq a^-\leq |a|$.
	\par Máme-li $\sum a_n$ absolutně konvergentní řadu, je i $\sum |a_n|$ konvergentní.
	Dostáváme, že $\sum a_n^+$ a $\sum a_n^-$ jsou konvergentní řady s nezápornými členy.
	Podle bodu (1) máme
	\begin{align*}
	\sum a_{\pi(n)}^+ &= \sum a_n^+,\\
	\sum a_{\pi(n)}^- &= \sum a_n^-.
\end{align*}
Dále je
$$
    \sum a_{\pi(n)} = \sum  (a_{\pi(n)}^+ - a_{\pi(n)}^-) = \sum a_{\pi(n)}^+ - \sum a_{\pi(n)}^-
$$
podle věty \ref{veta7.3}. Dále je
$$
    = \sum a_n^+ - \sum a_n^- = \sum (a_n^+ - a_n^-) = \sum a_n,
$$
čímž je důkaz u konce. \qedhere
\end{enumerate}
\end{proof}

\begin{veta}[Riemann] \label{veta7.14}

\end{veta}

\begin{proof}
	Hlavní myšlenka důkazu. Označme
	\begin{align*}
	P = \left \{ n \in \mathbb N , a_n \geq 0 \right \},  & Q = \left \{ n \in \mathbb N , a_n < 0 \right \}.
\end{align*}
Pak je $\mathbb N = P \cup Q$ a taky $P\cap Q = \emptyset$. Množiny $P$ a $Q$ jsou nekonečné (kdyby nějaká z nich
byla konečná, řada by byla absolutně konvergentní). Nalezneme posloupnosti přirozených čísel
$\left \{ p_n \right \}_{n=1}^\infty $ a $\left \{ q_n \right \}_{n=1}^\infty $ takové, že jsou prosté
a $P = \left \{ p_n, n \in \mathbb N  \right \} $, $Q = \left \{ q_n, n \in \mathbb N  \right \} $.
Dále je $\sum_{n=1}^\infty a_{p_n} = \infty$ a $\sum_{n=1}^\infty a_{q_n} = -\infty$, jinak by řada nebyla
neabsolutně konvergentní. Bereme členy z $P$, dokud částečný součet není větší než $s$. Pak beru z $Q$,
dokud částečný součet není menší než $s$, a tak podobně. Platí $\lim a_n = 0$, neboť $\sum a_n$ je konvergentní
řada. Proto částečné součty konvergují k $s$. Pro $s=\infty$ začneme s hranicí 1, pak zařadíme záporný člen,
pak máme hranici 2, a tak podobně.
\end{proof}

\subsection{Součin řad}
\begin{definice}

\end{definice}

\begin{veta}[Mertens] \label{veta7.15}

\end{veta}

\begin{proof}
	Označme
	\begin{align*}
		A_n &= \sum_{i=1}^na_i, & A &= \sum_{n=1}^\infty, & \alpha = \sum_{n=1}^\infty|a_n|,\\
	 B_m &= \sum_{i=1}^m b_i & B &= \sum_{n=1}^\infty b_n, & \beta _n = B_n-B,\\
	 C_n &= \sum_{i=1}^n, \quad c_i = \sum_{j=1}^i a_{i+1-j}b_j.
\end{align*}
Je
\begin{align*}
C_n &= \underbrace{(a_1b_1)}_{c_1} + \underbrace{(a_1b_2+a_2b_1)}_{c_2} + \dots + \underbrace{(a_1b_n + a_2 + b_{n-1}+ \dots + a_nb_1)}_{c_n}\\
&= a_1(b_1+b_2+\dots +b_n) + a_2(b_1+b_1+\dots+b_{n-1}) +\dots + a_nb_1\\
&= a_1 (B+ \beta_n) + a_2 (B+ \beta_{n-1}) +\dots +a_n(B+\beta_1)\\
&= (a_1+a_2+\dots+a_n)B + \undebrace{(a_1\beta_n + a_2\beta_{n-1} +\dots + a_n\beta_1)}_{\gamma_n}\\
&= A_nB+\gamma_n.
\end{align*}
Chceme, že $\lim C_n = A\cdot B$. Stačí ukázat, že $\lim \gamma_n = 0$, neboť
$
    \lim A_n = A
$ a
$$
    \lim C_n = \lim (A_nB+\gamma_n) = A\cdot B + 0 = A\cdot B.
$$
Budeme dokazovat $\lim_{n\to \infty} \gamma_n = 0$. Víme, že
\begin{equation}\label{limitarovnic}
    \lim_{n \to \infty} \beta _n = 0.
\end{equation}
Zvolme $\varepsilon > 0$. Díky (\ref{limitarovnic}) nalezneme $N \in \mathbb N $ takové, že
$$
    \forall n \geq N: |\beta_n|< \varepsilon.
$$
Potom pro každé $n\geq N$ platí
$$
    \gamma_n = (a_1\beta_n+a_2\beta_{n-1}+\dots+a_{n+1-N}\beta_N)+ (a_{n+2-N}\beta_{N-1} +\dots + a_n\beta_1).
$$
Odhadneme $|\gamma_n|$ jako
\begin{align*}
|\gamma_n| &\leq (|a_1|\cdot \underbrace{|\beta_n|}_{<\varepsilon} + \dots + |a_{n+1-N}|\cdot \underbrace{|\beta_N|}_{<\varepsilon})+ |a_{n+2-N}\beta_{N-1} +\dots + a_n\beta_1|\\
& \leq \alpha \varepsilon + |a_{n+2-N}\beta_{N-1} +\dots + a_n\beta_1|.
\end{align*}
Poněvadž $\lim_{j\to \infty} a_j = 0$, dostáváme, že
$$
0\leq \limsup |\gamma_n| \leq \alpha \cdot \varepsilon,
$$
a proto $\limsup |\gamma_n| = 0$. Dále platí
$$
    0 \leq \liminf |\gamma_n| \leq \limsup |\gamma_n|  = 0,
$$
a tedy $\liminf |\gamma _n| = 0$ a taky $\lim \gamma _n = 0$.
\end{proof}

\begin{priklad}
	Řada
	$$c_k\sum_{n=1}\infty \underbrace{\frac{(-1)^n}{\sqrt{n}}_{a_n = b_n}}$$ konverguje podle Leibnitzova kritéria, neboť $\frac{1}{\sqrt{n}}$
	konverguje. Řada
	$$
	c_k\sum _{i=1}^k a_{k+1-i}b_i = \sum_{i=1}^k \frac{(-1)^{k+i-1}}{\sqrt{k+1-i}}\cdot \frac{(-1)^i}{\sqrt{i}}= (-1)^{k+1} \sum_{i=1}^k \frac{1}{\sqrt{(k+1-i)i}},
	$$
	díky AG nerovnosti máme
	$$
 \frac{1}{\sqrt{(k+1-i)i}} \geq \frac{1}{\frac{(k+1-i)+i}{2}}=\frac{2}{k+1},
	$$
	a tedy i
	$$
	c_k \geq \frac{2}{k+1},
	$$
	a proto $\sum c_k$ diverguje, neboť není splněna nutná podmínka konvergence.
\end{priklad}

\begin{veta}[Abel] \label{veta7.16}

\end{veta}

\subsection{Taylorovy řady elementárních funkcí}
\begin{definice}

\end{definice}

\paragraph{Exponenciální funkce.} Použijeme Lagrangeův tvar zbytku (věta \ref{veta6.4}). Máme funkce
$f(x)=\exp (x)$ v bodě $a=0$. Uvažujme $x\ne a$. Pro každé $n \in \mathbb N $ nalezneme podle
věty \ref{veta6.4} bod $\xi _n$ mezi 0 a $x$ splňující
$$
    f(x)-T_n^{f,a}=\frac{1}{(n+1)!}f^{(n+1)}(\xi_n)x^{n+1}.
$$
Potom
$$
    |f(x)-T_n^{f,0}(x)| \leq  \frac{1}{(n+1)!}|x|^{n+1}e^{|x|}
$$
jde pro $n \to \infty$ k nule. Tedy
$$
    f(x) = \lim_{n \to \infty} T_n^{f,0}(x) = \lim_{n \to \infty} \sum_{j=0}^n \frac{1}{j!}x^j = \boxed{\sum_{j=0}^\infty \frac{1}{j!}x^j}.
$$
Odtud plyne pro $x=1$:
$$
    e = \sum_{j=0}^\infty \frac{1}{j!}.
$$
\paragraph{Logaritmická funkce.} Opět použijeme Lagrangeovu větu. Máme $f(x)=\log(1+x)$, $x \in \left ( 0,1 \right ] $.
Pro každé $n \in \mathbb N $ nalezneme podle věty \ref{veta6.4} $\xi_n \in (0,x)$ splňující
$$
    |f(x) - T_n^{f,0}(x)| = \left | \frac{1}{(n+1)!}f^{(n+1)}(\xi_n) x^{n+1} \right|.
$$
Je
\begin{align*}
	f^\prime(x) &= \frac{1}{1+x} \\
	f^{\prime\prime}(x) &= -\frac{1}{(1+x)^2} \\
	f^{\prime\prime\prime}(x) &= (-1)(-2)\frac{1}{(1+x)^3} \\
	& \dots \\
	f^{(n)}(x) &= (-1)(-2) \dots (-(n-1)) \frac{1}{(1+x)^n} = (-1)^{n-1}(n-1)! \frac{1}{(1+x)^n}.
\end{align*}
Přepíšeme tedy původní výraz
$$
    = \left |\frac{1}{(n+1)!} (-1)^n n! \frac{1}{(1+\xi_n)^{n+1}}x^{n+1} \right |\leq \frac{1}{n+1},
$$
a tedy
$$
\lim_{n \to \infty} T_n^{f,0}(x)=f(x).
$$

\par Dále ať je $x \in (-1,0)$. Podle věty \ref{veta6.5} nalezneme $\xi_n \in (x,0)$ splňující
\begin{align*}
|f(x)-T_n^{f,0}(x)| &= \left |\frac{1}{n!}f^{(n+1)}(\xi_n)(x-\xi_n)^n x\right |\\
	&= \left | \frac{1}{n!}(-1)^n \frac{1}{(1+\xi_n)^{n+1}}(x-\xi_n)^n x\right| = \left | \left( \frac{x-\xi_n}{1+\xi_n} \right )^n \frac{x}{1+\xi_n}  \right|,
\end{align*}
ale
\begin{align*}
\frac{\xi_n - x}{1+\xi_n} &\leq -x, \qquad \textrm{neboť} \\
\xi_n -x&\leq -x-x\xi_n \\
\xi_n(1+x)=\xi_n+x\xi_n &\leq 0.
\end{align*}
Původní výraz je tedy roven
$$
    = \left( \frac{\xi_n - x}{1+\xi_n} \right )^n \frac{|x|}{1-|x|} \leq \frac{|x|^{n+1}}{1-|x|},
$$
což pro $n\to \infty$ konverguje k nule.

Odtud plyne pro $x=1$:
$$
    \log 2 = \sum_{n=1}^\infty \frac{(-1)^{n-1}}{n}.
$$

\paragraph{Funkce sinus.} Máme $f(x)=\sin x$. Vezměme $x \in \mathbb R $, $x\ne 0$. Pro každé
$n \in \mathbb N $ nalezneme $\xi_n$ mezi 0 a $x$ splňující
$$
    f(x)-T_n^{f,0}(x) = \frac{1}{(n+1)!}f^{(n+1)}(\xi_n)x^{n+1}.
$$
Pak máme
$$
    |f(x) - T_n^{f,0}(x)| \leq \frac{1}{(n+1)!}|x|^{n+1},
$$
což pro $n \to \infty$ jde k nule.

\subsection{Posloupnosti a řady s komplexními členy}
\begin{definice}

\end{definice}

\begin{poznamka}
	\,
	\begin{enumerate}[(1)]
		\item $\left \{ a_n \right \} $ konverguje k $z$ právě tehdy, když $\lim_{n\to \infty} |a_n - z|=0$.
		\item Limita posloupnosti $\left \{ a_n \right \} $ je určena jednoznačně. Kdyby
		$\left \{ a_n \right \} $ konvergovala k $z_1$ a $z_2$, je
		$$
		    |z_1 -z_2| \leq \underbrace{|z_1-a_n|}_{\to 0} + \underbrace{|a_n-z_2|}_{\to 0},
		$$
		odkud máme $z_1=z_2$.
		\item Platí věta o aritmetice limit.
		\item $\lim a_n = z \iff (\lim \operatorname{Re} a_n = \operatorname{Re} z \land \lim \operatorname{Im}  a_n = \operatorname{Im} z)$.
	\end{enumerate}
\end{poznamka}

\begin{definice}

\end{definice}

\begin{definice}

\end{definice}

\begin{veta}\label{veta7.17}

\end{veta}

\begin{definice}

\end{definice}

\section{Primitivní funkce}
\subsection{Základní vlastnosti}
\begin{definice}

\end{definice}

\section{Metrické prostory I}
\begin{definice}

\end{definice}

\begin{priklad}
\begin{enumerate}[(1)]\,
	\item $(\mathbb R ^n, \rho_e)$, kde
	$$\rho_e(x,y) = \rho_e((x_1,\dots,x_n,(y_1,\dots,y_n)) = \sqrt{\sum_{i=1}^n (x_i-y_i)^2}$$
	je euklidovská metrika (někdy taky $l_2$-metrika).
	\item $(\mathbb R ^n, \rho_1)$, kde
	$$
	    \rho_1(x,y)=\sum_{i=1}^{n} |x_i-y_i|
	$$
	je $l_1$-metrika (newyorská metrika).
	\item $(\mathbb R ^n,\rho_\infty)$, kde
	$$
	    \rho_\infty(x,y) = \max \left \{ |x_i-y_i|, i \in \left \{ 1,\dots, n \right \}  \right \}
	$$
	je maximová metrika ($l_\infty$-metrika).
	\item $(\mathscr C([a,b]), \rho_{\sup{}})$ je prostor spojitých funkcí na $[a,b]$, kde
	$$
	    \rho_{\sup} (f,g) = \sup_{[a,b]}|f-g| = \sup \left \{ |f(x)-g(x)|, x \in [a,b] \right \}
	$$
	je supremová metrika.
	\item $(\mathscr C([a,b]),\rho_i$, kde
	$$
	    \rho_i(f,g) = \int_a^b |f-g|
	$$
	je integrální metrika.
	\item $(P,\rho_d)$, kde
	$$
	    \rho_d(x,y) = \begin{cases}
						0 & x=y,\\
						1 & x=y.
					\end{cases}
	$$
	je diskrétní metrika.
\end{enumerate}
\end{priklad}

\begin{definice}

\end{definice}

\begin{definice}

\end{definice}

\begin{veta}[Vlastnosti konvergence]\label{veta9.1}

\end{veta}

\begin{proof}\,
\begin{enumerate}[(a)]
	\item Předpokládejme, že $x,y \in P$ jsou limity $\left \{ x_n \right \}_{n=1}^\infty $.
	Potom platí
	$$
	    \rho(x,y) \leq \underbrace{\rho(x,x_n)}_{\to \, 0} + \underbrace{\rho(x_n, y)}_{\to \, 0},
	$$
	a proto $\rho(x,y)\leq 0$, a tedy $\rho(x,y) = 0$, a tedy $x=y$, což jsme chtěli dokázat.
	\item Pro každé $n \geq n_0$ platí $\rho(x_n,x)=\rho(x,x)=0$. Odtud máme
	$\lim_{n \to \infty} \rho(x_n,x)=0$, a tedy $x$ je limitou $\left \{ x_n \right \} $.
	\item Posloupnost $\left \{ \rho(x_{n_k},x) \right \} $ je vybraná z $\left \{ \rho(x_n,x) \right \} $.
	Protože je $\lim \rho(x_n,x)=0$, pak i $\lim \rho(x_{n_k},x)=0$, tedy $x$ je limita $\left \{ x_{n_k} \right \} $.\qedhere
\end{enumerate}
\end{proof}

\begin{oznaceni}

\end{oznaceni}

\begin{definice}

\end{definice}

\begin{veta}[Vlastnosti uzavřených množiny] \label{veta9.2}

\end{veta}

\begin{proof}\,
\begin{enumerate}[(a)]
	\item Zřejmé.
	\item Předpokládejme, že $\left \{ x_n \right \} $ je posloupnost prvků $\bigcap \mathcal F$
	a platí $\lim x_n= x \in P$. Zvolme $F \in \mathcal F$. Potom $\left \{ x_n \right \} $ je
	posloupnost prvků $F$. $F$ je podle předpokladu uzavřená, takže $x \in F$. Proto i $x \in \bigcap \mathcal F$.
	\item Předpokládejme, že $\left \{ x_n \right \} $ je posloupnost prvků $\bigcup _{i=1}^m F_i$
	a platí $\lim x_n=x$. Položme
	$$A_i = \left \{ n \in \mathbb N , x_n \in F_i \right \}, \qquad i=1,\dots,m. $$
	Platí $\mathbb N =\bigcup_{i=1}^m A_i$. Nalezneme $i_0\in \left \{ 1,\dots,m \right \} $
	takové, že $A_{i_0}$ je nekonečná. Existuje rostoucí posloupnost $\left \{ n_k \right \}_{k=1}^\infty $
	taková, že $A_{i_0}=\left \{ n_k, k \in \mathbb N  \right \} $. Potom $\left \{ x_{n_k} \right \}_{k=1}^\infty $
	je vybraná z $\left \{ x_n \right \} $ a platí $\lim x_{n_k} = x$ (podle věty \ref{veta9.2}) a
	pro každé $k \in \mathbb N $ platí $x_{n_k} \in F_{i_0}$. $F_{i_0}$ je uzavřená, a tedy
	$x \in F_{i_0} \subseteq \bigcup_{i=1}^m F_i$.\qedhere
\end{enumerate}
\end{proof}

\begin{priklad}
	Pro $(\mathbb R ^2, \rho_\infty)$ se jedná o \uv{otevřený čtverec}. Pro $(\mathbb R ^2, \rho_i)$
	jedná o otočený otevřený čtverec. Dále př. $(\mathscr C([a,b]), \rho_{\sup})$.
\end{priklad}

\begin{pozorovani}
	Nechť $(P,\rho)$ je metrický prostor, $x \in P, r >0$. Pro každé $y \in B(x,r)$ existuje $s>0$
	takové, že $B(y,s)\subseteq B(x,r)$. Stačí volit $0 < s < r-\rho(x,y)$.
\end{pozorovani}

\begin{definice}

\end{definice}

\section{Funkce více proměnných I}
\begin{oznaceni}
    \,
    \begin{itemize}
	\item $\mathbb R ^n = \underbrace{\mathbb R \times \mathbb R  \times \dots \times \mathbb R}_{k\ \text{krát}} $
	\item $x \in \mathbb R ^n$ znamená, že $x=[x_1,x_2,\dots,x_n]$
	\item $e^i = [0,0,\dots,0, \underset{i-\text{tý}}{1},0,\dots,0]$
	\item $o=[0,\dots,0]$
	\item $\nm{x}= \sqrt{\sum_{i=1}^\infty x_i^2}$
	\item někdy píšeme $\rho_e(x,y)=\nm{x-y}$
    \end{itemize}
\end{oznaceni}

\begin{definice}
Nechť $f$ je reálná funkce $n$ proměnných,
$a\in\mathbb{R}^n$
a $1\le i\le n$.
Pak \textbf{parciální derivaci funkce $f$ v bodě $a$
podle $i$-té proměnné} definujeme jako limitu

\[
\frac{\partial f}{\partial x_i}(a)
=
\lim_{t\to0}
\frac{
f(a+t e^i)-f(a)
}{t}.
\]
Symbolem
$
\frac{\partial f}{\partial x_i}
$
označujeme \textbf{parciální derivaci funkce $f$
podle $i$-té proměnné}, tj. funkci definovanou
předpisem
\[
\frac{\partial f}{\partial x_i}
:
x\mapsto
\frac{\partial f}{\partial x_i}(x).
\]
\end{definice}

\begin{poznamka}\,
\begin{enumerate}[(1)]
	\item Značení: místo symbolu $\frac{\partial f }{\partial x_i}(a)$ někdy používáme symbol $\partial_i f(a)$.
	\item Limita v definici
	$$
 \begin{cases}
		\text{neexistuje},\\
		\text{existuje} \ \begin{cases}
			\in \mathbb R \ \text{vlastní},\\
			\text{nevlastní} \ \begin{cases}
				\infty,\\
				-\infty.
\end{cases}
\end{cases}
 \end{cases}
	$$
	\item Pokud $\frac{\partial f }{\partial x_i}(a)$ existuje, pak
	$$
	    \exists \delta>0: \left \{ a+te^i, |t|<\delta \right \}\subset \mathcal D(f).
	$$
	\item Početní technika. Nechť $f$ je funkce $n$ proměnných, $a \in \mathbb R ^n,$ $1\leq i \leq n$.
	Parciální funkce je definována jako
	$$
	g(x)=f(a_1,\dots,a_{i-1},x,a_{i+1},\dots,a_n).
	$$
	Fukce $g$ je z $\mathbb R $ do $\mathbb R $.
	Platí
    $$
    g^\prime(a)=\frac{\partial f }{\partial x_i}(a),
    $$
    pokud má jedna strana smysl.
\end{enumerate}
\end{poznamka}

\begin{priklad}
	Vezměme funkci $f(x_1,x_2)=e^{x_1x_2^2}$. Spočteme parciální derivace
	\begin{align*}
		\frac{\partial f }{\partial x_1} (x_1,x_2) &= e^{x_1x_2^2}x_2^2,\\
		\frac{\partial f }{\partial x_2} (x_1,x_2) &= e^{x_1x_2^2}2x_1x_2.
\end{align*}
\end{priklad}

\begin{definice}
Nechť $f$ je reálná funkce $n$ proměnných,
$a\in\mathbb{R}^n$
a $L:\mathbb{R}^n\to\mathbb{R}$ je lineární zobrazení.
Řekneme, že $L$ je \textbf{totální diferenciál funkce
$f$ v bodě $a$}, jestliže platí
\[
\lim_{h\to0}
\frac{
f(a+h)-f(a)-L(h)
}{
\nm{h}
}
=
0.
\]
\end{definice}

\begin{poznamka}
	Pro $n=1$ máme
	$$
 \frac{f(a+h)-f(a)-Ah}{|h|} \to 0.
	$$
	Tento výraz můžeme ekvivalentně přepsat jako
	$$
 \frac{|f(a+h)-f(a)-Ah|}{|h|} \to 0.
	$$
	Jinými slovy,
	$$
	    \left |\frac{f(a+h)-f(a)-Ah}{h}\right | \to 0.
	$$
	Konečně
	$$
	    \left |\frac{f(a+h)-f(a)}{h} - A\right | \to 0.
	$$
	Proto $A=f^\prime(a)$.
	\par Pro $n=2$ máme
	$$
	    (x,y)\mapsto f(x,y).
	$$
	Lineární zobrazení bude
	$$
	    (x,y)\mapsto f(a_1,a_2) + L(x-a_1,y-a_2).
	$$
	\par Pro $n \in \mathbb N $ máme zobrazení $x \mapsto f(x)$ a lineární zobrazení je
	$$
	    x \mapsto f(a)+L(x-a).
	$$
	Pak
	$$
 \frac{f(x)-f(a)-L(x-a)}{\nm{x-a}}
	$$
	nám po jednoduché substituci dává definiční výraz.
\end{poznamka}

\begin{poznamka}
	Podmínka z definice ohledně aproximace je ekvivalentní podmínce
	$$
 \lim_{x\to a } \frac{f(x)-f(a)-L(x-a)}{\nm{x-a}}=0.
	$$
\end{poznamka}

\begin{veta}[Vztah totálního diferenciálu a parciální derivace]\label{veta10.1}
Nechť $L$ je totální diferenciál funkce
$f$ v bodě $a\in\mathbb{R}^n$.
Potom existují parciální derivace
\[
\frac{\partial f}{\partial x_1}(a),
\dots,
\frac{\partial f}{\partial x_n}(a)
\]
a pro každé $h\in\mathbb{R}^n$ platí
\[
L(h_1,\dots,h_n)
=
\frac{\partial f}{\partial x_1}(a)h_1
+\cdots+
\frac{\partial f}{\partial x_n}(a)h_n.
\]
\end{veta}

\begin{proof}
	Zobrazení $L$ pišme ve tvaru
	$$
	    L(h_1,h_2,\dots,h_n)=A_1h_1+ A_2h_2 + \dots + A_nh_n,
	$$
	kde $A_1,\dots,A_n \in \mathbb R $. Zvolme $i \in \left \{ 1,\dots,n \right \} $.
	Zobrazení $\varphi(t)=te^i$ je spojité. Pak je
	$$
	    \nm{\varphi(t)-\varphi(0)}=\nm{te^i-0e^i} = \nm{ [ 0,\dots, 0, \underset{i-\text{tý}}{t},0\dots,0 ] }=|t|.
	$$
	Potom podle věty \ref{veta9.10} platí
	$$
 \lim_{t\to 0} \frac{f \left( a+\varphi(t) - f(a) -L(\varphi(t)) \right ) }{\nm{\varphi(t)}} = 0,
	$$
	podmínka (P) je splněna. Upravíme
	$$
 \lim_{t\to 0} \frac{f(a+te^i)-f(a)-A_it}{|t|}=0.
	$$
	Odsud dostáváme, že
	$$
 \lim_{t \to 0} \left | \frac{f(a+te^i)-f(a)}{t}-A_i \right | = 0,
	$$
	a tedy
	$$
	\lim_{t \to 0} \frac{f(a+te^i) - f(a)}{t} = A_i.
	$$
	Proto $\frac{\partial f }{\partial x_i}(a)=A_i$.
\end{proof}

\begin{poznamka}
	Z věty \ref{veta10.1} plyne, že totální diferenciál je určen jednoznačně.
\end{poznamka}

\begin{oznaceni}
	Totální diferenciál značíme $f^\prime(a)$. Formálně je $f^\prime(a)\in \left( \mathbb R ^n \right )^* $.
	\par Pozor na kolizi značení! Pro $n=1$ je
	\begin{align*}
		f^\prime(a)\in \mathbb R && f^\prime(a) \in (\mathbb R)^* : \mathbb R \to \mathbb R
	\end{align*}
	třeba rozlišovat mezi derivací a totálním diferenciálem.
\end{oznaceni}

\begin{veta}\label{veta10.2}
Má-li funkce $f$ v bodě
$a\in\mathbb{R}^n$
totální diferenciál, je
$f$ v bodě $a$ spojitá.
\end{veta}

\begin{proof}
    Stačí ukázat, že následující limita je rovna $f(a)$:
    $$
    \lim_{x \to a } f(x) = \lim_{x \to a }  \left (\frac{f(x)-f(a)-f^\prime(a)(x-a) }{\nm{x-a}} \cdot \nm{x-a} + f(a) + f^\prime(a)(x-a)\right ),
    $$
    kde $f^\prime(a)$ je zobrazení z $\mathbb R ^n$ do $\mathbb R $, $h \mapsto f^\prime(a)(h)$.
    Protože je $x \mapsto \nm{x-a}$ vlastně
    $$x \mapsto \sqrt{\sum_{i=1}^n (x_i-a_i)^2},$$ je spojité.
    Konečně třetí člen je vlastně $$x \mapsto x-a = [x_1-a_1, \dots, x_n-a_n],$$ a tedy je též spojité.
    Po substituci je $$y \to f^\prime(a)(y) = \sum_{i=1}^n \frac{\partial f }{\partial x_i}(a)\cdot y_i.$$
    Máme tedy
    $$
    \lim_{x \to a } f(x) = 0\cdot 0 + f(a) + f^\prime(a) (\vec o) = f(a),
    $$
    což jsme chtěli dokázat.
\end{proof}

\begin{poznamka}
	\,
	\begin{enumerate}[(1)]
		\item Existuje-li $f^\prime(a)$, již jsme dokázali, že $\frac{\partial f }{\partial x_i}(a)\in \mathbb R $
		a také že $f$ je v $a$ spojité. Opačná implikace neplatí.
		\item Uvažujme
		$$
		    f(x,y) = \begin{cases}
							0 & x=0 \lor y=0, \\
							1 & \text{jinak}.
						\end{cases}
		$$
		Potom
		\begin{align*}
		\frac{\partial f }{\partial x }(0,0) = 0, & & \frac{\partial f }{\partial y}(0,0) = 0,
\end{align*}
ale $f$ není spojitá v $(0,0)$.
	\end{enumerate}
\end{poznamka}

\begin{lemma}\label{lemma10.3}
Nechť $f$ je reálná funkce $n$ proměnných,
$
I=
(\alpha_1,\beta_1)
\times\cdots\times
(\alpha_n,\beta_n)
\subset\mathbb{R}^n,
$
$a,b\in I$.
Nechť v každém bodě $I$ existují vlastní
parciální derivace $f$
podle všech proměnných.
Potom existují body
$
\xi_1,\dots,\xi_n\in I
$
takové, že
\[
f(b)-f(a)
=
\sum_{i=1}^{n}
\frac{\partial f}{\partial x_i}
(\xi_i)
(b_i-a_i).
\]
\end{lemma}

\begin{proof}
	Označme
	\begin{align*}
	p^0 &= (a_1,a_2,\dots,a_n) = a,\\
	p^1 &= (b_1,a_2,\dots, a_n),\\
	p^2 &= (b_1,b_2, \dots, a_n),\\
	&\dots \\
	p^n &= (b_1,b_2,\dots,b_n) = b.
\end{align*}
Potom platí
$$
    f(b) - f(a) = \sum_{k=1}^n \left( f(p^k)-f(p^{k-1}) \right ).
$$
Položme pro $k\in \left \{ 1,\dots,n \right \} $
$$
    g^k(x)=f(b_1, b_2,\dots, b_{k-1}, x, a_{k+1}, \dots, a_n).
$$
Potom má $g^k$ v každém bodě $(\alpha_k, \beta_k)$ vlastní derivaci, takže existuje $z_k \in [a_k, b_k]$
splňující (Lagrangeova věta \ref{veta5.7})
$$
    \underbrace{g^k(b_k)-g^k(a_k)}_{f(p^k)-f(p^{k-1})}=\underbrace{(g^k)^\prime(z_k)(b_k-a_k)}_{\frac{\partial f }{\partial x_k}(\xi^k)\cdot (b_k-a_k)}.
$$
Položíme $\xi^k = [b_1,b_2,\dots,b_{k-1},z_k, a_{k+1}, \dots,a_n]$.
\end{proof}

\begin{veta}\label{veta10.4}
Nechť $f$ je reálná funkce
$n$ proměnných,
$a\in\mathbb{R}^n$
a
\[
\frac{\partial f}{\partial x_1},
\dots,
\frac{\partial f}{\partial x_n}
\]
jsou spojité funkce v bodě $a$.
Potom má $f$ v bodě $a$
totální diferenciál.
\end{veta}

\begin{proof}
	Máme funkci $f:\mathbb R ^n \to \mathbb R $ a spolu s ní $n$-tici funkcí
	\begin{align*}
	\frac{\partial f }{\partial x_1} &: B(\alpha_1,\beta_1) \to \mathbb R , \\
	\frac{\partial f }{\partial x_2} &: B(\alpha_2,\beta_2) \to \mathbb R , \\
	& \dots\\
\frac{\partial f }{\partial x_n} &: B(\alpha_n,\beta_n) \to \mathbb R .
	\end{align*}
	Dále máme $L: \mathbb R ^n \to \mathbb R $ takové, že
	$$
	    L: k \mapsto \sum_{i=1}^n \frac{\partial f }{\partial x_i}(a)k_i.
	$$
	Ukážeme, že $L = f^\prime(a)$, tj. pro $L$ platí
	$$
 \lim_{x \to a} \frac{f(x)-f(a)-L(x-a)}{\nm{x-a}} = 0.
	$$
	Zvolme $\varepsilon>0$. K němu nalezneme $\tilde \delta >0$ takové, že
	$$
	    \forall i \in \left \{ 1, \dots, n \right \} \, \forall x \in B(a,\tilde \delta): \left | \frac{\partial f }{\partial x_i }(x) - \frac{\partial f }{\partial x_i} (a) \right |  <\varepsilon.
	$$
	Položme $\delta =\frac{\tilde \delta }{\sqrt n}$. Pak
	$$
	    I = (a_1-\delta, a_1+\delta)\times \dots \times (a_n-\delta,a_n+ \delta) \subset (a,\tilde,\delta).
	$$
	Zvolme $x \in B(a,\delta) \subseteq I \subseteq B(a,\tilde \delta)$. Podle lemmatu
	\ref{lemma10.3} nalezneme body $\xi^i (x)\in I$, $i =1,\dots,n $ takové, že
	$$
	    f(x)-f(a) = \sum_{i=1}^n \frac{\partial f }{\partial x_i} (\xi^i (x))(x_i-a_i).
	$$
	Pak máme
	\begin{align*}
	|f(x)-f(a)-L(x-a)|&=\left | \sum_{i=1}^n \frac{\partial f }{\partial x_i} (\xi^i(x))(x_i-a_i) - \sum_{i=1}^n \frac{\partial f }{\partial x_i}(a)(x_i-a_i) \right | \\
	&= \left | \sum_{i=1}^n \left( \frac{\partial f }{\partial x_i}(\xi^i(x)) - \frac{\partial f }{\partial x_i}(a) \right ) \cdot (x_i-a_i) \right |\\
	& \leq \sum_{i=1}^n \underbrace{\left | \frac{\partial f }{\partial x_i} (\xi^i(x))-\frac{\partial f }{\partial x_i}(a) \right |}_{< \, \varepsilon}  \cdot \underbrace{|x_i-a_i|}_{\leq\,\nm{x-a}}\\
	& \leq n\cdot  \varepsilon\cdot  \nm{x-a}.
\end{align*}
	Odsud plyne, že
	$$
 \lim_{x \to a } \frac{f(x)-f(a)-L(x-a)}{\nm{x-a}} = 0.
	$$
\end{proof}

\begin{veta}[Řetízkové pravidlo]\label{veta10.5}
Nechť funkce
$
f_1,\dots,f_k
$
z $\mathbb{R}^n$ do $\mathbb{R}$
mají v bodě
$a\in\mathbb{R}^n$
totální diferenciál
a funkce $g$ z $\mathbb R ^k$ do $\mathbb R $
má v bodě
$
b=
(f_1(a),\dots,f_k(a))
$
totální diferenciál.
Definujme funkci
$h:\mathbb{R}^n\to\mathbb{R}$ předpisem
\[
h(x)
=
g(
f_1(x),
\dots,
f_k(x)
).
\]
Potom má $h$ v bodě $a$
totální diferenciál a pro
$
i\in\{1,\dots,n\}
$
platí
\[
\frac{\partial h}{\partial x_i}(a)
=
\sum_{j=1}^{k}
\frac{\partial g}{\partial y_j}(b)
\frac{\partial f_j}{\partial x_i}(a).
\]
\end{veta}

\section{Obyčejné diferenciální rovnice}
\subsection{Základní pojmy}

\begin{definice}
Diferenciální rovnicí rozumíme rovnici tvaru
\begin{equation}\label{odrdef}
F(x,y,y',y'',\ldots,y^{(n-1)},y^{(n)})=0.
\end{equation}
kde $F$ je reálná funkce $n+2$ proměnných.

\textbf{Řešením diferenciální rovnice} (\ref{odrdef}) rozumíme funkci $y$ definovanou na nějakém neprázdném
otevřeném intervalu $I$, která má v každém bodě intervalu $I$ vlastní $n$-tou derivaci a jejíž hodnoty
spolu s hodnotami derivací splňují rovnici (\ref{odrdef}) v každém bodě intervalu $I$, tj. pro každé $x\in I$
platí
\[
F\bigl(x,y(x),y'(x),y''(x),\ldots,y^{(n-1)}(x),y^{(n)}(x)\bigr)=0.
\]
\end{definice}

\begin{priklad}
	Rovnice $p^\prime = ap - bp^2$ nebo $y^{\prime\prime} +y^\prime +\ = \sin x$. Tu lze převést na
$$y^{\prime\prime} +y^\prime +\ -\sin x= 0.$$
Tady je
	$$F(z_1,z_2,z_3,z_4)=--\sin z_1 + z_2 + z_3 + z_4.$$
\end{priklad}

\begin{definice}
Je-li funkce $y$ řešením rovnice (\ref{odrdef}) na intervalu $I$ a funkce $\widetilde y$ řešením rovnice (\ref{odrdef}) na
intervalu $\widetilde I$, kde $I\subset\widetilde I$, $I\ne\widetilde I$, a $y(x)=\widetilde y(x)$ pro všechna $x\in I$,
pak říkáme, že řešení $\widetilde y$ je \textbf{prodloužením řešení $y$} na interval $\widetilde I$.
Řešení rovnice (\ref{odrdef}), které nemá prodloužení, nazýváme \textbf{maximálním řešením} rovnice (\ref{odrdef}).
\end{definice}

\subsection{Rovnice se separovanými proměnnými}
\begin{definice}
Rovnice se separovanými proměnnými je rovnice tvaru
\begin{equation}\label{separ}
y'=g(y)h(x).
\end{equation}
kde $g,h$ jsou dané reálné funkce.
\end{definice}

\begin{pozorovani}
Označme $\mathcal D(h)=I$.
	Předpokládejme, že $g(c)=0$. Položme $y(x)=c$, $x \in I$. To je maximální řešení, neboť
	po dosazení je na levé straně $y^\prime(x)=0$ a na pravé straně
	$$
	    g(y(x))\cdot h(x) = g(c)\cdot h(x) = 0\cdot h(x)=0.
	$$
\end{pozorovani}

\begin{pozorovani}
Předpokládejme, že $J$ je otevřený interval, kde $g$ nemění znaménko a $y: L \to J$ řeší na otevřeném
	intervalu $L\subseteq I$ rovnici (\ref{separ}). Potom
	\begin{align*}
	\forall x \in L: \qquad y^\prime (x) &= g(y(x))\cdot h(x),\\
	\frac{y^\prime(x)}{g(y(x))}&= h(x).
	\end{align*}
	Nechť $H$ je primitivní k funkci $h$, tj. $H^\prime(x) = h(x)$, $x \in I$. Nechť $G$ je primitivní
	k funkci $\frac{1}{g}$, tj. $G^\prime(z)=\frac{1}{g(z)}$, $z\in J$.
	Protože je
	\begin{align*}
	(G(y(x)))^\prime &= \frac{1}{g(y(x))}\cdot y^\prime(x), \qquad x \in L \\
	H^\prime(x)&=h(x), \qquad x \in L.
	\end{align*}
	Existuje tedy $k \in \mathbb R $ takové, že
	$$
	    \forall x \in L: \qquad G(y(x)) = H(x) + k.
	$$
	Pak je
	$$
	    y(x) = G^{-1} (H(x)+k).
	$$
\end{pozorovani}

\paragraph{Metoda řešení}
Předpokládejme, že $\mathcal D(g)$ je interval, $\mathcal D(h)$ je otevřený interval, $g$ je spojitá na
$\mathcal D(g)$, a $h$ je spojitá
na $\mathcal D(h)$.

\paragraph{1. krok (kořeny a znaménko $g$).}
Označme $I=\mathcal D(h)$. Nalezneme všechny nulové body funkce $g$.
Je-li totiž $g(c)=0$, pak funkce $y(x)=c$, $x\in I$, je řešením rovnice (\ref{separ}).
Těmto řešením se říká \textbf{singulární nebo také stacionární}. Určíme maximální otevřené intervaly,
na kterých je $g$ nenulová.

\paragraph{2. krok (výpočet primitivních funkcí).}
Nechť $J$ je interval z 1. kroku, kde je $g$ nenulová.
Budeme hledat řešení, která jsou definovaná na nějakém otevřeném intervalu, který je obsažen
v $I$, a mají hodnoty v intervalu $J$. Nalezneme primitivní funkci $H$ k funkci $h$ na $I$ a primitivní
funkci $G$ k funkci $\frac{1}{g}$ na $J$. Existuje inverzní funkce
$G^{-1}\colon G(J)\to J$.

\paragraph{3. krok (určení definičních oborů).}
Zvolme $c\in\mathbb R$. Nalezneme maximální neprázdné otevřené intervaly obsažené v množině
\[
\{x\in I;\ H(x)+c\in G(J)\}.
\]

Nechť $L$ je takový interval. Na $L$ definujeme funkci $y$ předpisem
\[
y(x)=G^{-1}\bigl(H(x)+c\bigr),\qquad x\in L.
\]
Funkce $y$ je dobře definována a je řešením rovnice (\ref{separ}).

Ověření, že funkce $y(x)$ je řešením rovnice (\ref{separ}).

Na levé straně je
\begin{align*}
	y^\prime(x) &= \left( G^{-1}(H(x)+c) \right ) ^\prime = \frac{1}{G^\prime \left( G^{-1}(H(x)+c) \right ) }\cdot (H(x)+c)\\
	&=g ( \underbrace{G^{-1}(H(x)+c)}_{y(x)} )\cdot h(x) = g(y(x))\cdot h(x).
\end{align*}

\paragraph{4. krok (lepení).}
Z řešení nalezených ve 3. kroku a singulárních řešení z 1. kroku slepíme všechna maximální řešení.
Nechť $y_1$ a $y_2$ jsou řešení rovnice (\ref{separ}), první na intervalu $(a,b)$ a druhé na intervalu $(b,c)$,
přičemž $b\in I$. Předpokládejme, že
\[
\lim_{x\to b-}y_1(x)=\lim_{x\to b+}y_2(x)=\alpha\in \mathcal D(g).
\]
Pak funkce
\[
y(x)=
\begin{cases}
y_1(x), & x\in(a,b),\\
\alpha, & x=b,\\
y_2(x), & x\in(b,c),
\end{cases}
\]
je řešením rovnice (\ref{separ}) na intervalu $(a,c)$.


Kontrola, že po čtvrtém kroku je v bodě $b$ nová funkce opravdu řešení. Zřejmě je (\ref{separ}) splněna pro každé
$x \in (a,b) \cup (b,c)$. Zbývá ověřit (\ref{separ}) pro $x=b$. Funkce $y$ je spojitá zleva v bodě $b$. Je
\begin{align*}
	y_-^\prime(b) &= \lim_{x \to b-} y^\prime(x) = \lim_{x \to b-} y_1^\prime(x) = \lim_{x \to b-} g(y_1(x))\cdot h(x)\\
	&= g(\alpha)\cdot h(b)=g(y(b))\cdot h(b).
\end{align*}
Z pravé strany provedeme obdobně.

% \listoffigures

\backmatter

\end{document}
